% !TeX program = xelatex
\documentclass[11pt,a4paper,oneside]{article}

% --- Idioma, fuentes y formato ---
\usepackage[spanish,es-nodecimaldot]{babel}
\usepackage{fontspec}
\setmainfont{Arimo}
\setsansfont{Arimo}
\setmonofont{DejaVu Sans Mono}
\usepackage{microtype}
\usepackage[a4paper,top=2.35cm,bottom=2.25cm,left=2.3cm,right=2.3cm,headheight=16pt,headsep=0.55cm,footskip=0.9cm]{geometry}
\usepackage{graphicx}
\usepackage[table]{xcolor}
\definecolor{azulPrincipal}{HTML}{253E67}
\definecolor{azulTabla}{HTML}{4675C1}
\definecolor{negroTabla}{HTML}{080A0D}
\definecolor{grisClaro}{HTML}{EDF1F6}
\definecolor{grisTexto}{HTML}{606A76}
\definecolor{rangoNulo}{HTML}{D4DEEF}
\definecolor{rangoBajo}{HTML}{AFC3E3}
\definecolor{rangoModerado}{HTML}{8AA8D5}
\definecolor{rangoFuerte}{HTML}{315B9C}
\definecolor{rangoMuyFuerte}{HTML}{203C68}
\usepackage{array,tabularx,multirow,makecell,booktabs,longtable}
\usepackage{ragged2e}
\usepackage{enumitem}
\usepackage{float}
\usepackage{caption}
\usepackage{tikz}
\usepackage{titlesec}
\usepackage{fancyhdr}
\usepackage{setspace}
\usepackage{hyperref}
\usepackage{etoolbox}

% Denominaciones exigidas para los indices preliminares.
\addto\captionsspanish{%
  \renewcommand{\contentsname}{Tabla de contenido}%
  \renewcommand{\listfigurename}{Listado de Figuras}%
  \renewcommand{\listtablename}{Listado de Tablas}%
  \renewcommand{\figurename}{Figura}%
  \renewcommand{\tablename}{Tabla}%
}

\hypersetup{
    pdftitle={Acciones UD - Documento de Diseño y Arquitectura de la Solución},
    pdfauthor={CapitalFlow - Equipo de desarrollo},
    pdfsubject={Borrador estructurado del SAD},
    colorlinks=true, linkcolor=azulPrincipal, urlcolor=azulPrincipal
}
\setlength{\parindent}{0pt}
\setlength{\parskip}{6pt}
\emergencystretch=2em
\setstretch{1.10}
\setlist[itemize]{leftmargin=1.5em,topsep=2pt,itemsep=3pt}
\setlist[enumerate]{leftmargin=1.6em,topsep=2pt,itemsep=3pt}
\captionsetup{font=small,labelfont={bf,color=azulPrincipal},textfont=normal,skip=7pt}
\captionsetup[table]{position=top}
\titleformat{\section}{\Large\bfseries\color{azulPrincipal}}{\thesection.}{0.65em}{}[\vspace{-5pt}\titlerule]
\titleformat{\subsection}{\large\bfseries\color{azulPrincipal}}{\thesubsection}{0.65em}{}
\titleformat{\subsubsection}{\normalsize\bfseries\color{azulPrincipal}}{\thesubsubsection}{0.65em}{}
\titlespacing*{\section}{0pt}{19pt}{12pt}
\titlespacing*{\subsection}{0pt}{14pt}{5pt}
\titlespacing*{\subsubsection}{0pt}{12pt}{4pt}
\pagestyle{fancy}
\fancyhf{}
\fancyhead[L]{\footnotesize\color{grisTexto}\textbf{ACCIONES UD} \quad / \quad Documento SAD}
\fancyhead[R]{\footnotesize\color{grisTexto}Borrador de trabajo}
\fancyfoot[L]{\footnotesize\color{grisTexto}Diseño y arquitectura de software}
\fancyfoot[R]{\footnotesize\color{grisTexto}\thepage}
\renewcommand{\headrulewidth}{0.35pt}
\renewcommand{\footrulewidth}{0pt}

% --- Configuracion editable del proyecto ---
\newcommand{\nombreSolucion}{Acciones UD}
\newcommand{\nombreEmpresaDesarrollo}{CapitalFlow}
\newcommand{\integrantesEquipo}{%
  Juan José León Gómez\par
  Daniel Sanabria Rivera\par
  Davidson Esfleider Sánchez Gordillo\par
  Jan Henrik Sánchez Jerez\par
  María Fernanda Vélez Benítez%
}
\newcommand{\fechaCorte}{[DD/MM/AAAA]}
\newcommand{\versionSAD}{0.1 -- Borrador}
\newcommand{\pendiente}[1]{\textcolor{grisTexto}{\emph{[#1]}}}

% --- Auxiliares para formularios ---
\newcommand{\filaVacia}[1]{\rule{0pt}{#1}}
\newcommand{\rotuloAzul}[2]{\cellcolor{azulTabla}\parbox[c][#1][c]{\linewidth}{\RaggedRight\hyphenpenalty=10000\exhyphenpenalty=10000\footnotesize #2}}
\newcommand{\rotuloNegro}[2]{\cellcolor{negroTabla}\parbox[c][#1][c]{\linewidth}{\RaggedRight\hyphenpenalty=10000\exhyphenpenalty=10000\bfseries\footnotesize\color{white}#2}}
\newcommand{\cabeceraNegra}[1]{\cellcolor{negroTabla}\bfseries\color{white}#1}
\newlength{\anchoTres}
\newlength{\anchoCuatro}
\newlength{\anchoCinco}
\setlength{\anchoTres}{\dimexpr\textwidth-6\tabcolsep-4\arrayrulewidth\relax}
\setlength{\anchoCuatro}{\dimexpr\textwidth-8\tabcolsep-5\arrayrulewidth\relax}
\setlength{\anchoCinco}{\dimexpr\textwidth-10\tabcolsep-6\arrayrulewidth\relax}
\newcommand{\separadorTabla}{\vspace{0.35cm}}

% Los diagramas no se inventan. Esta macro deja un espacio claro y cuenta como figura.
\newcommand{\figuraPendiente}[3]{%
\begin{figure}[H]
\centering
\begin{tikzpicture}
  \draw[rounded corners=2pt,dashed,draw=grisTexto,line width=0.7pt] (0,0) rectangle (\dimexpr\textwidth-3mm\relax,#2);
  \node[align=center,text=grisTexto,text width=.80\textwidth] at ({(\textwidth-3mm)/2},{#2/2}) {\large\bfseries Diagrama pendiente\\[3pt]\normalsize Insertar el archivo o diagrama correspondiente\\[2pt]\small #1};
\end{tikzpicture}
\caption{#1 (pendiente de elaboración y validación).}\label{#3}
\end{figure}
}

% Plantilla de requisitos funcionales
\newcommand{\tablaRequisito}[6]{%
\begin{table}[H]
\centering
\begin{tabularx}{\linewidth}{|p{4cm}|X|}
\hline
\textbf{Identificador:} & #1 \\ \hline
\textbf{Nombre del requisito:} & #2 \\ \hline
\textbf{Actor / Fuente:} & #3 \\ \hline
\textbf{Prioridad:} & #4 \\ \hline
\textbf{Descripción del requisito:} &
#5 \\ \hline
\textbf{Método de verificación:} & #6 \\ \hline
\end{tabularx}
\caption{Requisito funcional #1}
\label{tab:#1}
\end{table}
}

\begin{document}

% ============================================================================
% PORTADA
% ============================================================================
\begin{titlepage}
\thispagestyle{empty}
\begin{tikzpicture}[remember picture,overlay]
  \fill[azulPrincipal] (current page.north west) rectangle ([yshift=-1.0cm]current page.north east);
  \fill[azulTabla] ([xshift=2.3cm,yshift=-6.55cm]current page.north west) rectangle ([xshift=4.0cm,yshift=-6.40cm]current page.north west);
\end{tikzpicture}
\vspace*{1.2cm}
{\footnotesize\bfseries\color{azulPrincipal} UNIVERSIDAD DISTRITAL FRANCISCO JOSÉ DE CALDAS}\\[5pt]
{\footnotesize INGENIERÍA DE SISTEMAS \quad / \quad DISEÑO ARQUITECTURAL DE SOFTWARE Y PATRONES}
\vspace*{1.2cm}

% Logo de la empresa de desarrollo (CapitalFlow)
\includegraphics[height=2.2cm,keepaspectratio]{images/capitalflow.png}
\vspace*{0.8cm}

{\Huge\bfseries\color{azulPrincipal}\nombreSolucion}\\[0.35cm]
{\LARGE\bfseries Documento de Diseño y\\Arquitectura de la Solución}\\[0.26cm]
{\Large (SAD)}
\vspace{1.0cm}

\rule{0.62\textwidth}{0.4pt}\\[0.45cm]
\begin{tabular}{@{}>{\bfseries}p{4.45cm}p{9.4cm}@{}}
Empresa de desarrollo: & \nombreEmpresaDesarrollo\\[0.35cm]
Nombre de la solución: & \nombreSolucion\\[0.35cm]
Versión del documento: & \versionSAD\\[0.35cm]
Fecha de corte: & \fechaCorte\\[0.35cm]
\end{tabular}
\vfill
{\bfseries Integrantes (ingenieros)}\\[0.25cm]
\integrantesEquipo
\vspace*{1.0cm}

{\small Bogotá, Colombia \hfill Proyecto final académico -- 2026}
\end{titlepage}

% ============================================================================
% ELEMENTOS PRELIMINARES
% ============================================================================
\pagenumbering{roman}
\tableofcontents
\clearpage
\listoffigures
\clearpage
\listoftables
\clearpage
\pagenumbering{arabic}

% ============================================================================
\section{Descripción del Documento}
\subsection{Propósito y Audiencia}
Este documento tiene como propósito reunir de manera verificable la definición inicial del problema, los requerimientos, los motivadores arquitecturales, los escenarios y las decisiones de diseño de \nombreSolucion. Servirá como referencia común para acordar el alcance, evaluar alternativas técnicas, construir el sistema y verificar que la arquitectura propuesta responde a las necesidades de negocio y a los atributos de calidad.

Está dirigido a los representantes del cliente y fundadores de Acciones UD, el equipo de ingeniería, los responsables de calidad y pruebas, los administradores de la solución, los comisionistas y los responsables de cumplimiento normativo que deban revisar aspectos de su funcionamiento. Su uso académico también permite documentar la trazabilidad entre el enunciado del proyecto y los entregables del curso.

\subsection{Organización del Documento}
El documento inicia con el contexto y los objetivos del proyecto. Después presenta los requerimientos funcionales y los motivadores arquitecturales, incluidos los requerimientos no funcionales y los escenarios de calidad. Continúa con el contexto operacional y los casos de uso, la propuesta arquitectural, los patrones de diseño y los puntos de vista y modelos solicitados. Finalmente, expone la relación entre dichos modelos, reserva el espacio de aceptación y reúne en los anexos los seis formatos base del enunciado.

\subsection{Convenciones}
\begin{itemize}
    \item Los identificadores \textbf{RF-XX} se reservan para requerimientos funcionales; \textbf{RNF-XX} para requerimientos no funcionales; \textbf{CU-XX} para casos de uso; y \textbf{EC-XX} para escenarios de calidad.
    \item Los apartados con la marca \pendiente{Por definir} requieren una decisión del equipo o validación del cliente. No se consideran requisitos aprobados.
    \item Las operaciones de compra y venta se describen desde la perspectiva de Acciones UD; la ejecución bursátil efectiva depende del proveedor y del mercado integrados.
    \item Los tiempos objetivos expresados en el enunciado deben interpretarse con su condición de medición especificada antes de incorporarse como criterios formales de aceptación.
    \item Los diagramas y tablas de anexos constituyen espacios editables que se actualizarán a medida que se tomen decisiones.
\end{itemize}

\subsection{Terminología y Definiciones}
\begin{description}[leftmargin=3.5cm,labelwidth=3.2cm,font=\normalfont\bfseries,style=nextline]
    \item[Acción] Instrumento financiero que representa una participación en la propiedad de una empresa.
    \item[Ticker] Símbolo utilizado para identificar un instrumento cotizado en el contexto de un mercado o proveedor de información.
    \item[Day trading] Actividad de negociación que busca abrir y cerrar posiciones dentro de una misma jornada de mercado.
    \item[Orden] Instrucción para comprar o vender un activo bajo determinadas condiciones de ejecución.
    \item[Market / Limit] Órdenes de mercado y de precio límite, respectivamente.
    \item[Stop-loss / Take-profit] Condiciones u órdenes destinadas a gestionar una salida por pérdida o por objetivo de ganancia, sujetas a la capacidad real de la integración.
    \item[Portafolio] Conjunto de posiciones o activos mantenidos por un inversionista.
    \item[Cotización] Información disponible sobre precios del instrumento, que puede comprender último precio negociado, precios bid/ask y otros datos, según la fuente.
    \item[Comisionista] Especialista de Acciones UD que asesora a inversionistas asociados y puede gestionar órdenes según las autorizaciones aplicables.
    \item[API] Interfaz de programación utilizada para integrar sistemas externos.
    \item[MFA] Autenticación multifactor: verificación de identidad mediante dos o más factores.
    \item[Trazabilidad] Capacidad de reconstruir el origen, los cambios y el resultado de una operación o evento.
    \item[SAD] Documento de Diseño y Arquitectura de Software (\emph{Software Architecture Document}).
\end{description}

% ============================================================================

\section{Generalidades del Proyecto}

\subsection{Problema por Resolver}

El acceso a los mercados bursátiles requiere información oportuna y confiable, herramientas de consulta, mecanismos para la compra y venta de acciones, administración de fondos y seguimiento de las operaciones. Para los inversionistas principiantes, la complejidad de los instrumentos financieros y los riesgos asociados a la toma de decisiones generan la necesidad de orientación y acompañamiento; mientras que los inversionistas experimentados requieren herramientas que faciliten el análisis de información y la gestión eficiente de sus inversiones.

AccionesUD busca atender estas necesidades mediante una plataforma web que centralice las principales actividades del inversionista, integrando servicios de consulta de información bursátil, negociación de activos, seguimiento de portafolios y asesoría financiera. La solución también contempla aspectos relacionados con la protección de datos, la seguridad de las operaciones, la transparencia de las comisiones y el control de acceso según los roles de inversionistas, comisionistas y administradores.

Sin embargo, la disponibilidad de información y la ejecución de operaciones bursátiles dependen de las capacidades técnicas y comerciales de los proveedores externos. Por esta razón, el proyecto debe delimitar su cobertura a los mercados e instrumentos efectivamente soportados por la integración seleccionada, garantizando que las funcionalidades ofrecidas correspondan con las posibilidades reales del sistema.

\subsection{Descripción General del Sistema a Desarrollar}

\nombreSolucion{} será una plataforma web responsiva orientada a inversionistas principiantes y experimentados, que permitirá consultar información bursátil, administrar perfiles de usuario, gestionar portafolios de inversión y realizar operaciones de compra y venta de acciones, de acuerdo con las funcionalidades habilitadas por el proveedor de servicios financieros integrado. Asimismo, contemplará la participación de comisionistas asociados a determinados inversionistas, junto con funcionalidades administrativas, de auditoría y generación de informes.

Para la consulta de cotizaciones, el acceso a datos de mercado y la gestión de operaciones bursátiles, se utilizará exclusivamente la API de \textbf{Alpaca}. Esta decisión responde a las restricciones de tiempo del cronograma académico y a los costos operativos y de licenciamiento asociados con la integración de múltiples proveedores, como Bloomberg o servicios avanzados de Alpha Vantage.

En consecuencia, el alcance de la plataforma estará condicionado por la cobertura de mercados, los instrumentos financieros disponibles y las capacidades de consulta y negociación ofrecidas por Alpaca. Aunque el proyecto se enmarca en el contexto de los mercados bursátiles internacionales, no se garantizará el acceso a todos ellos, sino únicamente a aquellos que puedan ser atendidos mediante la integración implementada. La disponibilidad de datos y operaciones dependerá, además, de las condiciones de acceso, los permisos y las limitaciones establecidas por el proveedor.

Entre las principales funcionalidades del sistema se contemplan la consulta de cotizaciones e información histórica, visualización de datos bursátiles, listas de observación, alertas configurables, gestión de órdenes de compra y venta, actualización de portafolios, administración de fondos, notificaciones, generación de reportes, seguimiento de transacciones y mecanismos de autenticación y autorización. También se incluirán controles relacionados con los horarios de negociación, días no hábiles y restricciones aplicables a los instrumentos y mercados soportados.

Adicionalmente, la plataforma incorporará servicios de asesoría mediante comisionistas y funcionalidades administrativas que permitan gestionar usuarios, supervisar operaciones y mantener la trazabilidad de las actividades realizadas dentro del sistema.


\subsection{Objetivos de la Solución}
\subsubsection{Objetivo General}
Diseñar e implementar una plataforma web de negociación y seguimiento de acciones que permita a inversionistas y comisionistas interactuar con mercados bursátiles compatibles a través de integraciones seguras, ofreciendo información clara, gestión controlada de operaciones y trazabilidad de las transacciones.

\subsubsection{Objetivos Específicos}
\begin{enumerate}
    \item Permitir el registro, la autenticación multifactor y la administración del perfil de los usuarios, con permisos acordes a su rol.
    \item Ofrecer consulta de instrumentos, cotizaciones, información histórica y visualizaciones del portafolio, en función de los datos que permitan la API integrada (Alpaca).
    \item Gestionar el ciclo de vida de las órdenes de compra y venta, las validaciones previas, las confirmaciones y el seguimiento de estados.
    \item Facilitar los servicios de asesoría de comisionistas, preservando la separación de información entre los portafolios de sus clientes asociados.
    \item Presentar al usuario el detalle de valores, cargos y comisiones, con reportes y notificaciones comprensibles.
    \item Proteger el acceso al sistema y conservar evidencia auditable de eventos y transacciones.
    \item Habilitar la configuración administrativa de mercados y ofrecer indicadores para el seguimiento operativo y del negocio.
    \item Diseñar la solución para los objetivos de concurrencia y respuesta establecidos en el enunciado, previa definición de las condiciones de prueba.
\end{enumerate}

\subsection{Stakeholders}

A continuación se identifican los principales grupos de interés
(\textit{stakeholders}) de Acciones UD, considerando tanto aquellos
que aparecen explícitamente en el enunciado como los que se derivan
directamente de la operación, administración y mantenimiento de
la plataforma.

\begingroup
\small
\setlength{\tabcolsep}{7pt}
\renewcommand{\arraystretch}{1.20}

\begin{table}[H]
\caption{Identificación y descripción de stakeholders.}
\begin{tabularx}{\textwidth}{|
>{\RaggedRight\arraybackslash}X|
>{\RaggedRight\arraybackslash}X|}\hline

\rowcolor{negroTabla}
\centering\color{white}\bfseries Stakeholder &
\centering\arraybackslash\color{white}\bfseries Descripción
\\\hline

Inversionistas principiantes y experimentados &
Usuarios finales que consultan información de los mercados
financieros, administran sus fondos y portafolios, configuran
preferencias y realizan solicitudes de compra y venta de acciones.
\\\hline

Comisionistas &
Especialistas en mercados financieros que asesoran a inversionistas
vinculados y, según las autorizaciones establecidas, revisan,
aprueban o gestionan órdenes en representación de sus clientes.
\\\hline

Administradores de la plataforma &
Personal responsable de gestionar la configuración operativa
del sistema, los mercados disponibles, los horarios de operación,
las restricciones y los parámetros de funcionamiento.
\\\hline

Área legal y de cumplimiento &
Responsables de supervisar el cumplimiento de las disposiciones
aplicables, revisar comportamientos anómalos y verificar la
trazabilidad y auditoría de las operaciones realizadas.
\\\hline

Fundadores, directivos y responsables del negocio &
Interesados en la sostenibilidad y crecimiento de Acciones UD,
así como en el seguimiento de indicadores de uso, volumen de
operaciones, ingresos y comisiones.
\\\hline

Equipo de desarrollo, calidad y operación &
Profesionales encargados del diseño, implementación, pruebas,
despliegue, mantenimiento y supervisión técnica de la plataforma,
garantizando su correcto funcionamiento.
\\\hline

Proveedores y aliados externos &
Entidades que suministran servicios de información bursátil,
ejecución de operaciones, gestión de fondos, procesamiento
de pagos y notificaciones mediante integraciones con
la plataforma.
\\\hline

\end{tabularx}
\end{table}

\separadorTabla

\begin{table}[H]
\caption{Expectativas de los stakeholders.}
\begin{tabularx}{\textwidth}{|
>{\RaggedRight\arraybackslash}X|
>{\RaggedRight\arraybackslash}X|}\hline

\rowcolor{negroTabla}
\centering\color{white}\bfseries Stakeholder &
\centering\arraybackslash\color{white}\bfseries Expectativas
\\\hline

Inversionistas principiantes y experimentados &
Disponer de una plataforma intuitiva, segura y disponible,
con información bursátil actualizada, ejecución confiable
de órdenes, protección de sus fondos y datos personales,
y acceso al historial y estado de sus inversiones.
\\\hline

Comisionistas &
Consultar exclusivamente la información de sus clientes
asociados, gestionar las órdenes autorizadas, verificar
su estado y disponer de mecanismos seguros de aprobación,
autorización y seguimiento.
\\\hline

Administradores de la plataforma &
Contar con herramientas para configurar y controlar
la operación del sistema, administrar restricciones,
supervisar su funcionamiento y detectar o atender
incidencias operativas.
\\\hline

Área legal y de cumplimiento &
Garantizar la trazabilidad de las transacciones,
la conservación de registros de auditoría, el control
de accesos y la disponibilidad de información para
investigar operaciones sospechosas y verificar el
cumplimiento normativo.
\\\hline

Fundadores, directivos y responsables del negocio &
Obtener información confiable sobre el crecimiento
de la plataforma, la actividad de los usuarios,
los volúmenes negociados y los ingresos por comisiones,
facilitando la toma de decisiones estratégicas.
\\\hline

Equipo de desarrollo, calidad y operación &
Disponer de una arquitectura mantenible, escalable,
segura y observable, con mecanismos de monitoreo,
diagnóstico, pruebas y recuperación ante fallos que
faciliten la evolución y continuidad del servicio.
\\\hline

Proveedores y aliados externos &
Contar con interfaces de integración claras y seguras,
intercambio confiable de información, mecanismos de
autenticación y manejo de errores, y condiciones de
interoperabilidad acordes con los servicios contratados.
\\\hline

\end{tabularx}
\end{table}

\endgroup

% ============================================================================
\section{Especificación y Recabación de Requerimientos Funcionales}
La siguiente es una \textbf{identificación inicial} de capacidades derivadas del enunciado académico. No reemplaza el levantamiento de requerimientos con el cliente ni determina prioridades definitivas. Cada requisito deberá evolucionar hacia una especificación verificable con actores, reglas, entradas, salidas, excepciones y criterios de aceptación.

\subsection{Gestión de Usuarios, Seguridad de Acceso y Perfil de Inversión}

\tablaRequisito
{RF-01}
{Registro de inversionistas}
{Inversionista}
{Alta}
{El sistema deberá permitir el registro de inversionistas solicitando nombres, apellidos, correo electrónico, teléfono de contacto, país de residencia y credenciales de acceso. Durante el registro se podrá configurar la recepción del reporte mensual de inversiones; las demás preferencias de notificación, portafolio y operaciones se gestionarán en la configuración del perfil. La necesidad de información adicional para verificar la identidad queda sujeta a recabación.}
{Pruebas funcionales de registro, validación de datos personales, consulta y actualización del perfil y elección inicial del reporte mensual.}


\tablaRequisito
{RF-02}
{Autenticación segura de usuarios}
{Inversionista / Comisionista / Administrador de la plataforma / Analista de Seguridad y Auditoría / Directivo del negocio}
{Alta}
{El sistema deberá autenticar a los usuarios registrados de los roles Inversionista, Comisionista, Administrador de la plataforma, Analista de Seguridad y Auditoría y Directivo del negocio. La autenticación deberá incorporar mecanismos de autenticación multifactor (MFA) antes de permitir el acceso a las funcionalidades protegidas.}
{Pruebas funcionales de autenticación multifactor y pruebas de seguridad de acceso para cada tipo de usuario.}


\tablaRequisito
{RF-03}
{Gestión de roles y permisos}
{Administrador de la plataforma}
{Alta}
{El sistema deberá permitir al Administrador de la plataforma gestionar los roles y permisos de Inversionista, Comisionista, Administrador de la plataforma, Analista de Seguridad y Auditoría y Directivo del negocio, restringiendo el acceso a funcionalidades y datos según las autorizaciones establecidas.}
{Pruebas funcionales de control de acceso y autorización por rol.}


\tablaRequisito
{RF-04}
{Construcción y actualización del perfil de actividad inversora}
{Inversionista}
{Alta}
{El sistema deberá construir y actualizar el perfil de actividad inversora a partir de las transacciones realizadas por el inversionista. El análisis deberá identificar su comportamiento de negociación, como frecuencia, volumen y distribución de operaciones por mercado e instrumento. Si la cuenta todavía no tiene operaciones suficientes, se indicará que el perfil se encuentra sin historial suficiente. Esta clasificación basada en transacciones no se considerará, por sí sola, una medición de la tolerancia al riesgo.}
{Pruebas funcionales de análisis transaccional, actualización del perfil y manejo de cuentas sin historial suficiente.}


\tablaRequisito
{RF-05}
{Consulta de estadísticas descriptivas de inversión}
{Inversionista}
{Media}
{El sistema deberá presentar al inversionista estadísticas descriptivas calculadas a partir de sus transacciones, como número de operaciones, volumen negociado en USD y distribución porcentual de la actividad por mercado, instrumento financiero y, cuando la información esté disponible, sector económico. Los porcentajes deberán corresponder al período consultado y mostrar claramente la base de cálculo utilizada. No se generarán recomendaciones personalizadas a partir de estas estadísticas.}
{Pruebas funcionales de estadísticas por mercado, instrumento y sector, y verificación de cantidades y porcentajes con datos conocidos.}


\subsection{Gestión de Cuentas y Suscripciones}

\tablaRequisito
{RF-06}
{Consulta de información de la cuenta}
{Inversionista}
{Alta}
{El sistema deberá permitir al inversionista consultar la información general de su cuenta, incluidos sus datos de identificación, estado de la cuenta y vinculaciones con los servicios habilitados de la plataforma.}
{Pruebas funcionales de consulta de cuenta y validación de acceso a datos propios.}


\tablaRequisito
{RF-07}
{Configuración de preferencias de notificaciones}
{Inversionista}
{Alta}
{El sistema deberá permitir al inversionista configurar desde su perfil sus preferencias para recibir avisos de apertura y cierre de mercados, resúmenes diarios de operaciones y confirmaciones de órdenes ejecutadas, mediante los canales habilitados. La configuración de alertas por umbrales de precio se realizará en el módulo de listas de observación y requerirá una suscripción premium.}
{Pruebas funcionales de actualización y persistencia de preferencias y de generación de avisos según la configuración.}


\tablaRequisito
{RF-08}
{Configuración de la visualización del portafolio}
{Inversionista}
{Media}
{El sistema deberá permitir al inversionista seleccionar desde su perfil la visualización predeterminada del portafolio (gráfico de barras, lista o vista detallada), ordenar los activos por valor total o porcentaje de ganancia o pérdida y mostrar u ocultar los activos de menor peso en el portafolio. La interfaz deberá aplicar las preferencias guardadas.}
{Pruebas funcionales de selección de vista, ordenamiento, ocultación de activos y persistencia de preferencias.}


\tablaRequisito
{RF-09}
{Configuración de preferencias y controles de operaciones}
{Inversionista}
{Alta}
{El sistema deberá permitir al inversionista definir desde su perfil el tipo de orden predeterminado entre las modalidades habilitadas (\textit{market}, \textit{limit}, \textit{stop-loss} y \textit{trailing stop}), establecer un número máximo de operaciones diarias y activar una autorización adicional para órdenes que superen el umbral definido como operación de valor elevado. Estas preferencias deberán aplicarse durante la creación y validación de órdenes.}
{Pruebas funcionales de configuración de tipo de orden, límites diarios, segunda autorización y aplicación en las operaciones.}


\tablaRequisito
{RF-10}
{Contratación y pago de suscripciones premium}
{Inversionista / Stripe}
{Media}
{El sistema deberá permitir contratar una suscripción premium durante el registro inicial o posteriormente desde la sección de Gestión de Perfil. Deberá ofrecer las modalidades de USD \$12 mensuales y USD \$120 anuales, procesando los pagos mediante Stripe. Una vez confirmado el estado de la suscripción, se habilitarán las listas de observación y las alertas de precios por umbrales, identificadas como funcionalidades por suscripción.}
{Pruebas funcionales de contratación en ambos puntos de acceso, verificación de precios y pruebas de integración con Stripe.}


\tablaRequisito
{RF-11}
{Cancelación de suscripciones premium}
{Inversionista / Stripe}
{Media}
{El sistema deberá permitir al inversionista solicitar la cancelación de su suscripción premium y consultar el resultado de la solicitud. La fecha efectiva de cancelación y la conservación temporal de beneficios dependerán de las condiciones del plan contratado.}
{Pruebas funcionales de cancelación y verificación de los estados informados por Stripe.}


\tablaRequisito
{RF-12}
{Consulta del estado de suscripciones}
{Inversionista / Stripe}
{Media}
{El sistema deberá permitir consultar el plan de suscripción, su estado, período de vigencia y la información disponible sobre pagos asociados, de acuerdo con los datos suministrados por Stripe.}
{Pruebas funcionales de consulta del estado y verificación de datos con la integración de pagos.}


\subsection{Gestión de Comisionistas}

\tablaRequisito
{RF-13}
{Selección y vinculación con comisionistas}
{Inversionista / Comisionista}
{Alta}
{Durante el registro, el sistema deberá sugerir al inversionista un comisionista habilitado y permitirle seleccionar uno desde un listado de comisionistas disponibles. El listado mostrará, cuando se disponga de ellos, nombre, especialidad, experiencia y calificaciones. El inversionista podrá acceder posteriormente a esta sección para seleccionar otro comisionista mediante el procedimiento de desvinculación y nueva vinculación. La validación de certificaciones, la vigencia y las reglas de aceptación quedan pendientes de definir.}
{Pruebas funcionales de sugerencia, listado, selección, vinculación y verificación de permisos.}


\tablaRequisito
{RF-14}
{Desvinculación de comisionistas}
{Inversionista / Comisionista}
{Media}
{El sistema deberá permitir al inversionista solicitar la desvinculación de su comisionista asociado, aplicando las autorizaciones y restricciones operativas que se definan. Una vez efectiva la desvinculación, el comisionista dejará de acceder a la información y operaciones del inversionista como cliente activo, sin afectar los registros históricos sujetos a conservación.}
{Pruebas funcionales de desvinculación, actualización de permisos y preservación de registros históricos.}


\tablaRequisito
{RF-15}
{Consulta del historial de comisionistas}
{Inversionista}
{Media}
{El sistema deberá permitir consultar el historial de asociaciones entre inversionistas y comisionistas, incluidos los comisionistas anteriores, períodos de vigencia y estado de las vinculaciones, respetando los permisos del usuario.}
{Pruebas funcionales de consulta del historial y validación de acceso a asociaciones registradas.}


\tablaRequisito
{RF-16}
{Creación y envío de propuestas de inversión}
{Comisionista / Inversionista}
{Alta}
{El sistema deberá permitir que un comisionista genere y remita propuestas de inversión a sus clientes autorizados. Una propuesta podrá incluir una o varias operaciones, con empresa o instrumento, símbolo, tipo de operación, cantidad, precio indicativo de referencia, costos estimados y justificación de la estrategia, según los campos que se acuerden. Al presentarla, se notificará al inversionista; el precio mostrado en la propuesta tendrá una referencia temporal identificable.}
{Pruebas funcionales de creación, asociación a clientes autorizados, envío y notificación de propuestas.}


\tablaRequisito
{RF-17}
{Modificación de propuestas de inversión}
{Comisionista}
{Alta}
{El sistema deberá permitir al comisionista modificar el contenido de propuestas propias que aún no hayan sido aceptadas por el inversionista. Cuando se modifique una propuesta ya notificada, deberá conservarse la trazabilidad de los cambios y notificarse nuevamente su contenido al inversionista.}
{Pruebas funcionales de edición, restricciones por estado, registro de cambios y notificación.}


\tablaRequisito
{RF-18}
{Eliminación o retiro de propuestas no aprobadas}
{Comisionista}
{Media}
{El sistema deberá permitir al comisionista retirar propuestas propias que aún no hayan sido aceptadas, evitando que puedan ser aprobadas posteriormente. Para conservar la trazabilidad, el retiro no deberá eliminar los registros requeridos para auditoría.}
{Pruebas funcionales de retiro, bloqueo de aceptación posterior y conservación del registro.}


\tablaRequisito
{RF-19}
{Aceptación o rechazo de propuestas de inversión}
{Inversionista}
{Alta}
{El sistema deberá permitir al inversionista revisar las propuestas recibidas y aceptarlas o rechazarlas expresamente. La aceptación deberá registrarse y comunicarse al comisionista correspondiente. Esta aceptación no implicará por sí sola el envío ni la ejecución de las órdenes de compra o venta.}
{Pruebas funcionales de consulta, aceptación, rechazo, registro de decisiones y notificación al comisionista.}


\tablaRequisito
{RF-20}
{Autorización del comisionista y tramitación de propuestas aceptadas}
{Comisionista / Proveedor de ejecución bursátil}
{Alta}
{El sistema deberá permitir que el comisionista autorizado confirme o firme lógicamente su conformidad con una propuesta previamente aceptada por el inversionista. Se registrarán la identidad de quien confirma y el momento de la acción. Cumplidas las validaciones de fondos, precio y horario, las órdenes asociadas se tramitarán mediante el mecanismo de ejecución integrado. El significado técnico y jurídico de la firma queda pendiente de recabación.}
{Pruebas funcionales de doble autorización, registro de evidencias y bloqueo de envío sin aprobaciones requeridas.}


\tablaRequisito
{RF-21}
{Consulta del historial de propuestas (propuesto)}
{Inversionista}
{Media}
{El sistema deberá permitir al inversionista consultar todas las propuestas registradas para su cuenta por los comisionistas con los que haya estado vinculado, incluyendo propuestas pendientes, aceptadas, rechazadas o retiradas, sus fechas y estados, sin permitir el acceso a propuestas de otros inversionistas.}
{Pruebas funcionales de consulta histórica, filtros por estado y controles de acceso a propuestas propias.}


\subsection{Calificaciones de Comisionistas}

\tablaRequisito
{RF-22}
{Registro de calificaciones y comentarios de comisionistas}
{Inversionista}
{Media}
{El sistema deberá permitir que los inversionistas que mantengan o hayan mantenido una vinculación con un comisionista registren una calificación y un comentario sobre el servicio recibido. Los criterios de puntuación, la frecuencia permitida y las reglas de publicación deberán definirse con el cliente.}
{Pruebas funcionales de registro de calificación y comentario y validación de vinculación previa.}


\tablaRequisito
{RF-23}
{Consulta de calificaciones y comentarios de comisionistas}
{Inversionista / Comisionista}
{Media}
{El sistema deberá permitir consultar las calificaciones y comentarios publicados sobre cada comisionista, incluyendo una calificación agregada cuando existan suficientes datos, como apoyo para la selección del profesional.}
{Pruebas funcionales de consulta de reseñas y verificación del cálculo de la calificación agregada.}


\subsection{Mercados e Información Bursátil}

\tablaRequisito
{RF-24}
{Consulta de mercados, empresas e instrumentos financieros}
{Inversionista / Comisionista / Proveedor de datos bursátiles}
{Alta}
{El sistema deberá permitir buscar empresas e instrumentos financieros y consultar los mercados habilitados en los que se encuentren disponibles. Presentará información como símbolo, nombre de la empresa, sector económico, precio y volumen negociado, según los datos proporcionados por el proveedor bursátil.}
{Pruebas funcionales de búsqueda de empresas, consulta de instrumentos y pruebas de integración bursátil.}


\tablaRequisito
{RF-25}
{Visualización de cotizaciones e información histórica}
{Inversionista / Comisionista / Proveedor de datos bursátiles}
{Alta}
{El sistema deberá presentar tableros para visualizar cotizaciones e históricos de instrumentos financieros. Mostrará, cuando el proveedor los suministre, el mejor precio de compra (\textit{bid}), el mejor precio de venta (\textit{ask}) y su diferencial (\textit{spread}), calculado como el precio \textit{ask} menos el precio \textit{bid}. Los precios y diferenciales se presentarán en USD; cuando la fuente utilice otra moneda, se aplicará la conversión disponible. Se indicará la fecha y hora de la última cotización recibida. La actualización estará sujeta a los permisos del usuario y a las características de la fuente de datos.}
{Pruebas funcionales de visualización, verificación de diferencial y pruebas de integración con datos bursátiles.}


\tablaRequisito
{RF-26}
{Gestión de listas de observación y alertas de precios}
{Inversionista}
{Media}
{El sistema deberá permitir que los inversionistas con suscripción premium agreguen y quiten acciones o instrumentos de sus listas de observación. También deberá permitir configurar alertas de precio por encima o por debajo de umbrales definidos para los activos observados. Al cumplirse una condición, se generará la notificación correspondiente conforme a las preferencias y canales habilitados del inversionista.}
{Pruebas funcionales de incorporación y eliminación de activos, umbrales superiores e inferiores, notificaciones y bloqueo sin suscripción premium.}


\tablaRequisito
{RF-27}
{Consulta de horarios y condiciones de mercado}
{Inversionista / Comisionista / Administrador de la plataforma}
{Alta}
{El sistema deberá permitir consultar horarios de negociación, feriados, cierres y restricciones operativas de los mercados habilitados, utilizando la configuración vigente y la información de los proveedores integrados. También deberá generar avisos de apertura y cierre de mercado para los inversionistas que hayan activado esas notificaciones en su perfil.}
{Pruebas funcionales de consulta de calendarios y horarios, y avisos de apertura y cierre según preferencias de notificación.}


\tablaRequisito
{RF-28}
{Comparación de mercados y selección de vía de ejecución}
{Inversionista / Comisionista / Proveedor de ejecución bursátil}
{Alta}
{Cuando una empresa o instrumento pueda negociarse en más de un mercado habilitado, el sistema deberá comparar la información disponible de las alternativas y determinar la vía de tramitación aplicable a la operación. La posibilidad de seleccionar el mercado de ejecución y el criterio de comparación dependerán de la API y las reglas de negociación que se acuerden.}
{Pruebas funcionales de comparación y pruebas de integración del mecanismo de enrutamiento disponible.}


\subsection{Gestión de Órdenes Bursátiles}

\tablaRequisito
{RF-29}
{Creación de órdenes de compra y venta}
{Inversionista / Comisionista}
{Alta}
{El sistema deberá permitir crear órdenes de compra y venta de acciones, registrando el instrumento, sentido, cantidad y parámetros de ejecución correspondientes. Se contemplarán inicialmente órdenes \textit{market} y \textit{limit}, así como mecanismos \textit{stop-loss}, \textit{take-profit} y \textit{trailing stop}, según las modalidades habilitadas por el proveedor. Cuando exista una preferencia de tipo de orden predeterminado, se aplicará al iniciar el formulario, sin impedir su modificación dentro de las opciones permitidas. Los tipos de vigencia y las modalidades adicionales de ejecución deberán confirmarse con el cliente. Crear una orden no implicará su aceptación ni su ejecución en el mercado.}
{Pruebas funcionales de creación y validación de órdenes, aplicación de preferencias y comprobación de modalidades admitidas.}


\tablaRequisito
{RF-30}
{Validación de fondos y tenencia de activos}
{Inversionista / Comisionista / Proveedor financiero}
{Alta}
{Antes de cursar una compra, el sistema deberá verificar la disponibilidad de fondos suficientes para cubrir el importe y las comisiones aplicables. Antes de cursar una venta, deberá validar la tenencia suficiente del activo. Además, deberá aplicar el límite de operaciones diarias configurado por el inversionista y exigir la autorización adicional cuando una orden cumpla las condiciones definidas para operaciones de valor elevado. Las verificaciones utilizarán la información financiera disponible y las restricciones operativas correspondientes.}
{Pruebas funcionales con fondos y posiciones suficientes e insuficientes, límites diarios y autorización adicional para órdenes grandes.}


\tablaRequisito
{RF-31}
{Verificación de variaciones de precio y nueva autorización}
{Inversionista / Comisionista / Proveedor de datos bursátiles}
{Alta}
{Antes de enviar una orden, el sistema deberá contrastar el precio de referencia mostrado durante la autorización con la cotización disponible para la tramitación. Si se detecta una variación que, según la regla acordada, exija una nueva aceptación, se deberá detener el envío y solicitar confirmación actualizada del inversionista o descartar la solicitud según corresponda. El criterio de variación permitido queda pendiente de definición.}
{Pruebas funcionales con precios sin variación, cambios que exijan ratificación y operaciones no autorizadas.}


\tablaRequisito
{RF-32}
{Envío y seguimiento de órdenes bursátiles}
{Inversionista / Comisionista / Proveedor de ejecución bursátil}
{Alta}
{El sistema deberá transmitir órdenes validadas y autorizadas al proveedor de ejecución bursátil, registrar sus respuestas y permitir consultar su evolución. Deberá conservar estados, confirmaciones, rechazos y resultados de ejecución, y generar las confirmaciones de órdenes ejecutadas según las preferencias y los canales habilitados del inversionista. Los demás cambios relevantes de estado deberán poder consultarse en la plataforma.}
{Pruebas de integración de envío, seguimiento de estados, registro de respuestas y confirmaciones de ejecución según preferencias.}


\tablaRequisito
{RF-33}
{Control del plazo de respuesta del proveedor de ejecución}
{Proveedor de ejecución bursátil}
{Alta}
{El sistema deberá controlar un plazo de cinco segundos para obtener la respuesta del proveedor que se defina como confirmación de la operación. Si se supera dicho plazo, deberá impedir nuevos envíos automáticos de la misma solicitud, consultar el estado de la orden y solicitar la cancelación cuando todavía sea posible. La plataforma no deberá declarar una orden cancelada ni no ejecutada hasta recibir o verificar el estado efectivo del proveedor. El punto de inicio del plazo y la respuesta que lo satisface deben precisarse con el cliente.}
{Pruebas de integración con demoras simuladas, estados desconocidos, cancelaciones pendientes y confirmaciones tardías.}


\tablaRequisito
{RF-34}
{Cancelación de órdenes bursátiles}
{Inversionista / Comisionista / Proveedor de ejecución bursátil}
{Alta}
{El sistema deberá permitir solicitar la cancelación de órdenes que se encuentren en un estado cancelable, conforme a las reglas del proveedor. Deberá registrar la solicitud y reflejar el resultado real comunicado por la API, sin considerar la cancelación efectiva hasta recibir la confirmación correspondiente.}
{Pruebas funcionales y de integración con órdenes cancelables, ejecutadas y pendientes de cancelación.}


\tablaRequisito
{RF-35}
{Programación de órdenes fuera del horario de mercado}
{Inversionista / Comisionista / Proveedor de ejecución bursátil}
{Media}
{El sistema deberá permitir que, previa confirmación del interesado y cuando la integración lo admita, las órdenes iniciadas fuera del horario operativo permanezcan pendientes para su tramitación el siguiente día hábil aplicable. Antes de liberarlas se comprobarán las condiciones de mercado, fondos y variaciones de precio que puedan requerir una nueva autorización.}
{Pruebas funcionales de cola fuera de horario, cambio de sesión, días hábiles y revalidación antes del envío.}


\subsection{Gestión de Fondos}

\tablaRequisito
{RF-36}
{Vinculación con proveedor financiero y consulta de saldo}
{Inversionista / Proveedor financiero}
{Alta}
{El sistema deberá permitir vincular la cuenta del inversionista con el proveedor financiero acordado y consultar los saldos habilitados por dicha integración. Los importes monetarios se presentarán en dólares estadounidenses (USD), de acuerdo con las conversiones y datos disponibles.}
{Pruebas funcionales de vinculación, consulta de saldo y presentación en USD.}


\tablaRequisito
{RF-37}
{Solicitud de depósito de fondos}
{Inversionista / Proveedor financiero}
{Alta}
{El sistema deberá permitir solicitar depósitos a la cuenta de inversión mediante los mecanismos habilitados por el proveedor financiero. La operación deberá registrar el importe y la moneda de origen y permitir consultar el estado confirmado por el proveedor. Los fondos depositados se destinarán a su expresión y, cuando corresponda, conversión a USD.}
{Pruebas funcionales de depósitos y pruebas de integración de estados y confirmaciones.}


\tablaRequisito
{RF-38}
{Conversión de fondos depositados a USD}
{Inversionista / Proveedor financiero}
{Alta}
{El sistema deberá gestionar mediante el proveedor autorizado la conversión a dólares estadounidenses (USD) de los fondos depositados en otras monedas, registrando el importe original, moneda, tipo de cambio aplicado, posibles cargos e importe resultante. La moneda principal de presentación de saldos y operaciones financieras será USD. El proveedor y las condiciones de conversión requieren validación.}
{Pruebas de integración de conversión y verificación de importes, tipo de cambio y registro de cargos.}


\tablaRequisito
{RF-39}
{Solicitud de retiro de fondos}
{Inversionista / Proveedor financiero}
{Alta}
{El sistema deberá permitir solicitar el retiro de fondos conforme al saldo disponible, a los mecanismos autorizados y a las restricciones del proveedor financiero. Deberá registrar la solicitud y presentar su estado y resultado en USD, sin asumir que el destino bancario necesariamente recibe la misma moneda.}
{Pruebas funcionales de retiro con fondos suficientes e insuficientes y pruebas de integración del estado real.}


\tablaRequisito
{RF-40}
{Consulta y registro de movimientos de fondos}
{Inversionista / Proveedor financiero}
{Alta}
{El sistema deberá registrar y permitir consultar los movimientos asociados a la cuenta del inversionista, incluidos depósitos, retiros y conversiones de moneda cuando estén disponibles. Deberá conservar el tipo de movimiento, importe, fecha, estado y, cuando aplique, moneda de origen e importe equivalente en USD.}
{Pruebas funcionales de historial de movimientos y verificación de consistencia con el proveedor.}


\subsection{Gestión de Portafolios y Reportes}

\tablaRequisito
{RF-41}
{Consulta y seguimiento del portafolio e historial de operaciones}
{Inversionista / Comisionista}
{Alta}
{El sistema deberá permitir al inversionista consultar las posiciones de su portafolio, instrumentos, cantidades, valoraciones e historial de operaciones, presentando los valores monetarios en USD. Los comisionistas solo podrán acceder a las posiciones e historial de los clientes con los que mantengan una autorización vigente.}
{Pruebas funcionales de consulta del portafolio, historial de operaciones, presentación monetaria y autorización por cliente.}


\tablaRequisito
{RF-42}
{Generación y entrega de reportes de inversión}
{Inversionista / Comisionista}
{Media}
{El sistema deberá generar reportes por períodos sobre operaciones, posiciones, ganancias y pérdidas y las comisiones efectivamente aplicadas, desglosadas entre Acciones UD y el comisionista cuando corresponda. Para posiciones abiertas, calculará la ganancia o pérdida no realizada a partir de la diferencia entre precio actual y costo de adquisición, considerando los costos de compra que correspondan. Para operaciones cerradas, calculará la ganancia o pérdida realizada utilizando los valores de compra y venta y sus comisiones efectivas. Los importes se presentarán en USD, con acceso limitado a la información autorizada. Cuando el inversionista active el resumen diario, el sistema lo generará a partir de las operaciones del día. Si durante el registro o la gestión de perfil se habilita el reporte mensual, al cierre de cada mes el sistema lo generará automáticamente, lo pondrá a disposición dentro de la plataforma y notificará su disponibilidad.}
{Pruebas funcionales de reportes y resúmenes diarios, cálculos de ganancias y pérdidas realizadas y no realizadas, costos de transacción y generación del reporte mensual según preferencia.}


\subsection{Gestión de Comisiones}

\tablaRequisito
{RF-43}
{Cálculo de comisiones}
{Administrador de la plataforma / Inversionista / Comisionista}
{Alta}
{El sistema deberá calcular las comisiones de las operaciones utilizando el porcentaje configurado por la administración para el mercado correspondiente. Antes de confirmar la operación deberá mostrar al inversionista el porcentaje aplicable, el importe estimado de la comisión y el costo total estimado, expresados en USD. Cuando intervenga un comisionista, deberá calcular la distribución de la comisión entre este y Acciones UD según las reglas acordadas. Los importes definitivos se registrarán cuando existan datos confirmados, sin asumir la liquidación o transferencia de pagos al comisionista.}
{Pruebas de cálculo por mercado, desglose de costos antes de confirmar y verificación de distribución de comisiones.}


\subsection{Seguridad y Auditoría de Operaciones}

\tablaRequisito
{RF-44}
{Registro y consulta de eventos de auditoría}
{Analista de Seguridad y Auditoría / Administrador de la plataforma}
{Alta}
{El sistema deberá registrar los accesos y eventos relevantes de seguridad, los cambios de estado de órdenes y los eventos de propuestas de inversión, incluidas modificaciones, aceptaciones, retiros, autorizaciones y solicitudes de cancelación. El Analista de Seguridad y Auditoría y el Administrador de la plataforma, en el ámbito de sus permisos, deberán poder consultar estos registros con la trazabilidad necesaria para auditoría.}
{Pruebas funcionales de registro de eventos, consulta restringida y conservación de historial de cambios.}


\tablaRequisito
{RF-45}
{Análisis de operaciones atípicas y gestión de restricciones}
{Analista de Seguridad y Auditoría}
{Alta}
{El sistema deberá proporcionar al Analista de Seguridad y Auditoría funciones de consulta y análisis de operaciones posiblemente atípicas y permitir aplicar restricciones sobre cuentas cuando corresponda y cuente con autorización. Los criterios de identificación y los procedimientos para imponer o levantar restricciones quedan sujetos a definición con el cliente.}
{Pruebas funcionales de consulta, aplicación de restricciones y validación de autorizaciones.}


\subsection{Administración y Analítica del Negocio}

\tablaRequisito
{RF-46}
{Configuración de mercados, horarios y parámetros operativos}
{Administrador de la plataforma}
{Alta}
{El sistema deberá permitir a los administradores configurar los mercados habilitados, sus horarios operativos, calendarios de cierre, restricciones parametrizables y el porcentaje de comisión correspondiente a cada mercado. La configuración deberá respetar las condiciones reales de negociación, así como las capacidades de los proveedores externos.}
{Pruebas funcionales de modificación de horarios, cierres, mercados y porcentajes de comisión con control de permisos.}


\tablaRequisito
{RF-47}
{Consulta de indicadores de negocio}
{Directivo del negocio}
{Media}
{El sistema deberá proporcionar indicadores de usuarios registrados, ingresos por comisiones, volumen de operaciones y evolución de la actividad de la plataforma, calculados a partir de los datos registrados y disponibles.}
{Pruebas funcionales de consulta de indicadores y verificación con datos de prueba.}


\subsection{Cuestiones Pendientes de Recabación}

La especificación funcional constituye una versión inicial sujeta a la
validación con el cliente. Las preguntas abiertas se referencian
mediante la numeración vigente de los requisitos:

\begin{itemize}

\item \textbf{RF-01}: Confirmar si la plataforma debe solicitar
  información adicional de identificación o verificación y qué datos
  son obligatorios para los diferentes tipos de usuarios.

\item \textbf{RF-02, RF-03}: Especificar los mecanismos de MFA,
  la recuperación de acceso, la asignación de roles y las autorizaciones
  concretas de Inversionista, Comisionista, Administrador de la plataforma,
  Analista de Seguridad y Auditoría y Directivo del negocio.

\item \textbf{RF-04}: Confirmar si el perfil de actividad se basa
  exclusivamente en transacciones o si también se desea un cuestionario
  inicial para conocer objetivos, experiencia y tolerancia al riesgo.
  Definir cuándo existe historial suficiente y cómo se actualiza
  la clasificación. El comportamiento transaccional no equivale,
  por sí solo, al perfil de tolerancia al riesgo.

\item \textbf{RF-05}: Precisar el período predeterminado y si
  los porcentajes de actividad se calculan sobre cantidad de operaciones,
  volumen negociado u otra base. Confirmar la cobertura de clasificación
  sectorial disponible en el proveedor de datos.

\item \textbf{RF-07, RF-08, RF-09}: Definir cuáles serán los canales
  de notificación habilitados, qué peso se considera bajo al ocultar
  activos, cómo se cuenta el límite diario de operaciones, qué valor
  constituye una orden grande y cómo se realizará la autorización adicional.

\item \textbf{RF-10, RF-11, RF-12, RF-26}: Confirmar las reglas
  de renovación y cancelación efectiva de las suscripciones premium,
  pagos fallidos y reembolsos, así como el alcance exacto de las alertas
  y listas de observación. Las tarifas son USD \$12 mensuales
  y USD \$120 anuales; estas cifras no están pendientes.

\item \textbf{RF-13, RF-14, RF-15}: Definir la información pública
  que se mostrará de los comisionistas, cómo se verifican sus
  certificaciones, cuál es el criterio de sugerencia inicial,
  si la vinculación tiene una vigencia y qué condiciones limitan
  la desvinculación y la nueva selección.

\item \textbf{RF-16, RF-17, RF-18}: Confirmar los campos obligatorios
  y reglas de una propuesta, su vencimiento, la posibilidad de modificarla
  o retirarla y el tratamiento de precios y cambios comunicados al
  inversionista antes de su aceptación.

\item \textbf{RF-19, RF-20}: Precisar el significado técnico y jurídico
  de la firma o segunda aceptación del comisionista, los registros
  de evidencia exigidos y las condiciones que obligan a volver
  a solicitar autorización del inversionista.

\item \textbf{RF-21}: Validar con el cliente la inclusión del requisito
  \emph{propuesto} de consulta de historial de propuestas, sus filtros,
  su cobertura respecto de comisionistas anteriores y la duración
  de la conservación del historial.

\item \textbf{RF-22, RF-23}: Definir la escala de calificación,
  frecuencia permitida de reseñas, moderación de comentarios,
  calificaciones agregadas y opiniones tras finalizar una vinculación.

\item \textbf{RF-24, RF-25, RF-27, RF-28}: Precisar el alcance
  de mercados e instrumentos, cobertura y retraso de datos de
  cotización (incluidos \textit{bid}, \textit{ask} y \textit{spread}),
  y si el proveedor permite realmente elegir el mercado de ejecución.

\item \textbf{RF-09, RF-29}: Confirmar qué tipos de órdenes y vigencias
  se deben implementar, incluidas DAY, GTC, IOC, FOK, OPG y CLS,
  así como las condiciones aplicables a \textit{trailing stop},
  \textit{stop-loss} y \textit{take-profit} y su disponibilidad real
  en el proveedor.

\item \textbf{RF-31}: Precisar qué variación de precio exige
  una nueva aprobación, si se compara precio cotizado o precio ejecutable
  y cuándo debe rechazarse una orden que ha cambiado de condiciones.

\item \textbf{RF-32, RF-33, RF-34}: Definir desde qué momento corren
  los cinco segundos y qué respuesta del proveedor satisface el plazo
  (recepción, aceptación o ejecución). Acordar cómo tratar respuestas
  tardías, estados desconocidos, cancelaciones pendientes y operaciones
  ejecutadas de manera parcial.

\item \textbf{RF-27, RF-35, RF-46}: Confirmar cómo se consideran
  feriados, cambios de horario y colas de órdenes pendientes para la
  siguiente sesión, y qué datos deben revalidarse antes del envío.

\item \textbf{RF-36, RF-37, RF-38, RF-39, RF-40}: Confirmar el
  proveedor de fondos, medios de depósito y retiro, divisas de origen,
  conversión efectiva a USD, tipos de cambio, cargos e inconsistencias
  entre los registros internos y las API externas.

\item \textbf{RF-41, RF-42}: Acordar el método de costo de adquisición
  cuando se compren posiciones en diferentes momentos, la fuente de
  cotizaciones para valores no realizados, el tratamiento de cargos
  adicionales y el período y contenido del resumen diario. El reporte
  mensual se genera al cierre de mes, cuando la preferencia está activa,
  y queda disponible dentro de la plataforma.

\item \textbf{RF-43, RF-46}: Establecer los porcentajes concretos de
  comisión por mercado, la regla de reparto entre plataforma y
  comisionista y si se incluirá posteriormente la liquidación o
  transferencia efectiva de estos importes.

\item \textbf{RF-44, RF-45}: Definir el plazo de conservación de
  registros de auditoría, los criterios de detección de operaciones
  atípicas, las acciones permitidas al Analista de Seguridad y
  Auditoría y las condiciones para imponer o levantar restricciones.

\item \textbf{RF-46, RF-47}: Precisar qué parámetros puede modificar
  el Administrador de la plataforma y cuáles son los indicadores y
  niveles de detalle permitidos al Directivo del negocio.

\end{itemize}


\subsection{Cuestiones Pendientes de Recabación}

La especificación funcional presentada constituye una versión inicial sujeta a validación con el cliente. Las siguientes decisiones permanecen abiertas y se relacionan con los requisitos funcionales indicados:


\begin{itemize}

    \item \textbf{RF-01, RF-04}: Confirmar los datos obligatorios de registro y verificación de identidad, las preguntas exactas para determinar experiencia y riesgo, los criterios de puntuación y la posibilidad de actualizar la clasificación.

    \item \textbf{RF-02, RF-03}: Definir los mecanismos concretos de MFA, la recuperación del acceso y los permisos detallados de inversionistas, comisionistas, administradores, cumplimiento y responsables del negocio.

    \item \textbf{RF-05}: Definir qué interacciones del usuario se analizarán, cómo se calcularán los porcentajes por mercado e instrumento y si se mostrarán recomendaciones personalizadas o solo estadísticas descriptivas.

    \item \textbf{RF-07, RF-08, RF-09}: Confirmar los beneficios específicos de la suscripción premium, las reglas de renovación automática, cancelación efectiva, pagos fallidos y reembolsos. Los precios de USD \$12 al mes y USD \$120 al año ya están definidos.

    \item \textbf{RF-10, RF-11, RF-12}: Determinar qué información pública se mostrará de cada comisionista, cómo funcionará la sugerencia inicial, si la vinculación tiene vigencia y qué restricciones aplican para cambiarlo o desvincularlo.

    \item \textbf{RF-13, RF-14, RF-15}: Definir los campos obligatorios y reglas de una propuesta, su posible vencimiento, cuándo puede modificarse o retirarse y si los precios mostrados permanecen como referencia fija o se actualizan. Precisar cómo se informan los cambios al inversionista.

    \item \textbf{RF-16, RF-17}: Precisar si la aceptación del inversionista y la confirmación del comisionista constituyen una firma electrónica suficiente, qué evidencia se conserva y cómo se procede cuando cambia el precio antes del envío. (Yo diria que si que no haria falta nada mas)

    \item \textbf{RF-18}: Validar con el cliente la inclusión del requisito propuesto de consulta del historial completo de propuestas por inversionista, así como su alcance y período de conservación.

    \item \textbf{RF-19, RF-20}: Definir escala y criterios de calificación de comisionistas, condiciones para opinar, frecuencia permitida, moderación de comentarios y tratamiento de reseñas de asociaciones terminadas.

    \item \textbf{RF-21, RF-22, RF-24, RF-25}: Precisar mercados e instrumentos habilitados, cobertura y retraso de los datos bursátiles, qué cotizaciones podrán compararse, cuál será la fuente de datos (IEX, SIP u otra) y si Alpaca u otro proveedor permitirá seleccionar explícitamente el mercado de ejecución.

    \item \textbf{RF-26}: Acordar qué tipos de órdenes y condiciones de vigencia deben implementarse obligatoriamente, incluidos DAY, GTC, IOC, FOK, OPG y CLS, así como las modalidades de apertura o cierre y las combinaciones admitidas por la API.

    \item \textbf{RF-28}: Definir cuál es el cambio de precio que invalida la autorización anterior, si se solicita una nueva confirmación o se descarta la orden y si la regla aplica por igual a órdenes \textit{market}, \textit{limit} y propuestas de comisionistas.

    \item \textbf{RF-29, RF-30, RF-31}: Definir desde qué evento se cuentan los cinco segundos y cuál respuesta los satisface: recepción, aceptación del proveedor o ejecución. Establecer el manejo de confirmaciones tardías, ejecuciones parciales, cancelaciones pendientes e inconsistencias entre el registro interno y el estado de la API.

    \item \textbf{RF-24, RF-32, RF-43}: Acordar qué sucede al cambiar horarios o al dejar órdenes en cola para la siguiente sesión, cómo se determina el siguiente día hábil y qué validaciones requieren nueva autorización antes del envío.

    \item \textbf{RF-33, RF-34, RF-35, RF-36, RF-37}: Confirmar el proveedor de fondos, las formas de depósito y retiro, las monedas de origen y destino, el mecanismo efectivo de conversión a USD, los tipos de cambio, cargos y manejo de inconsistencias entre Acciones UD y las API externas.

    \item \textbf{RF-40, RF-43}: Establecer los porcentajes de comisión para cada mercado, la regla de reparto con los comisionistas y si habrá posteriormente un mecanismo de liquidación o pago a estos.

    \item \textbf{RF-41, RF-42}: Definir el período de conservación de auditoría, los criterios de operaciones atípicas, las acciones permitidas al personal de cumplimiento y el alcance de las restricciones sobre cuentas.

    \item Los mercados e instrumentos financieros que estarán
    disponibles y las capacidades reales de los proveedores
    seleccionados.

    \item Los tipos de órdenes admitidos, sus condiciones de
    vigencia y las modalidades de ejecución que deberán soportarse.

    \item Las reglas de asociación entre inversionistas y
    comisionistas, así como los procedimientos de aprobación,
    autorización y firma de operaciones.

    \item El significado operativo de confirmar una transacción,
    diferenciando la recepción de la solicitud, su validación,
    la aceptación del proveedor y la ejecución efectiva.

    \item Los mecanismos de depósito y retiro de fondos,
    incluyendo las operaciones que serán gestionadas directamente
    por la plataforma o delegadas a proveedores externos.

    \item Las monedas admitidas y las reglas aplicables
    a posibles conversiones de divisas.

    \item El tratamiento de inconsistencias entre el estado
    registrado por Acciones UD y el informado por las API
    de los proveedores externos.

\end{itemize}

\subsection{Cuestiones Pendientes de Recabación}
Antes de cerrar la especificación se debe acordar, entre otros asuntos, el alcance real de los mercados por proveedor; los tipos de órdenes, sus vigencias y modalidades; la política de aprobación y firma por comisionista; el significado de ``confirmar una transacción''; la gestión de depósitos y retiros; el manejo de divisas; y el comportamiento ante inconsistencias entre el sistema y las API externas.

% ============================================================================
\section{Motivadores Arquitecturales}
\subsection{Motivadores de Negocio}
El proyecto AccionesUD contempla los siguientes objetivos de negocio, los cuales orientan las decisiones arquitecturales y permiten establecer las prioridades funcionales, técnicas y económicas de la solución.

\textbf{Objetivo de negocio preliminar.} Ofrecer en Colombia una plataforma web de servicios bursátiles digitales que facilite a inversionistas principiantes y experimentados la consulta de información financiera y la gestión de operaciones de compra y venta de acciones, mediante la integración con la API de Alpaca. La solución contará con apoyo opcional de comisionistas, mecanismos de monetización mediante comisiones y suscripciones, y controles de seguridad y trazabilidad que promuevan la confianza y transparencia en las operaciones.
\begin{enumerate}
    \item \textbf{Ampliar el acceso a la inversión bursátil:} Proporcionar una plataforma web accesible para usuarios con diferentes niveles de experiencia, facilitando la consulta de información financiera y el acceso a los instrumentos bursátiles soportados por Alpaca.
    \item \textbf{Generar confianza y transparencia:}  Presentar información clara y verificable sobre las operaciones realizadas, sus estados, resultados y cargos asociados, permitiendo que los usuarios comprendan y supervisen sus actividades financieras.
    \item \textbf{Apoyar una operación segura y controlada:} Implementar mecanismos de autenticación, autorización, gestión de permisos y trazabilidad que permitan proteger la información de los usuarios y controlar las actividades de inversionistas, comisionistas y administradores.
    \item \textbf{Sostener el crecimiento del negocio:} Diseñar una solución capaz de atender la concurrencia prevista, facilitar el seguimiento de usuarios y transacciones y permitir su evolución funcional, considerando las capacidades y restricciones de los servicios externos integrados.
    \item \textbf{Habilitar ingresos por servicios:} Soportar mecanismos de monetización mediante comisiones parametrizables y suscripciones premium, garantizando la transparencia de los cargos y la diferenciación de los servicios ofrecidos.
    \item \textbf{Garantizar la viabilidad económica y operativa de la solución:} Optimizar los recursos de desarrollo e integración mediante la utilización de un único proveedor de servicios bursátiles, reduciendo la complejidad técnica y los costos asociados a múltiples integraciones, en concordancia con las restricciones presupuestales y el cronograma académico del proyecto.
\end{enumerate}

Como posibles indicadores de seguimiento se consideran el número de usuarios activos, las transacciones procesadas, el volumen de operaciones, los ingresos por comisiones, la adopción de suscripciones premium, los eventos de seguridad y la disponibilidad de los servicios bursátiles integrados. Sus metas cuantitativas se encuentran pendientes de definición.


\begingroup
\footnotesize
\setlength{\tabcolsep}{6pt}
\setlength{\anchoTres}{\dimexpr\textwidth-6\tabcolsep-4\arrayrulewidth\relax}
\renewcommand{\arraystretch}{1.16}

\begin{table}[H]
\caption{Motivador de negocio: Viabilidad económica y operativa de AccionesUD.}

\begin{tabular}{|
>{\RaggedRight\arraybackslash}m{.40\anchoTres}|
>{\RaggedRight\arraybackslash}m{.28\anchoTres}|
>{\RaggedRight\arraybackslash}m{.32\anchoTres}|}
\hline

\cabeceraNegra{\scriptsize Nombre del Motivador de Negocio}
&
\multicolumn{2}{
>{\columncolor{negroTabla}\color{white}\RaggedRight\arraybackslash}
p{\dimexpr.60\anchoTres+2\tabcolsep+\arrayrulewidth\relax}|
}{
\bfseries Descripción del Motivador de Negocio
}\\
\hline

\textbf{Garantizar la viabilidad económica y operativa}
&
\multicolumn{2}{p{\dimexpr.60\anchoTres+2\tabcolsep+\arrayrulewidth\relax}|}{
Desarrollar una plataforma bursátil funcional y sostenible,
optimizando costos y tiempos de implementación mediante
una arquitectura de Monolito Modular con Arquitectura
Hexagonal y la integración exclusiva con la API de Alpaca
para los servicios bursátiles soportados.
}\\
\hline

\multicolumn{3}{
|>{\columncolor{negroTabla}\color{white}\RaggedRight\arraybackslash}l|
}{
\bfseries\rule{0pt}{.72cm}Medida del Impacto
}\\
\hline

\multicolumn{3}{|p{\dimexpr\textwidth-2\tabcolsep-2\arrayrulewidth\relax}|}{
Porcentaje de hitos técnicos completados dentro
del presupuesto y cronograma aprobados.

Indicador = (Hitos cumplidos dentro de ambas
restricciones / Total de hitos evaluados) $\times$ 100.
}\\
\hline

\rowcolor{negroTabla}
\bfseries\color{white}Rangos
&
\bfseries\color{white}Cota Mínima
&
\bfseries\color{white}Cota Máxima\\
\hline

\cellcolor{rangoNulo}Ninguno
& 0\% & 19\%\\
\hline

\cellcolor{rangoBajo}Bajo
& 20\% & 39\%\\
\hline

\cellcolor{rangoModerado}Moderado
& 40\% & 59\%\\
\hline

\cellcolor{rangoFuerte}\color{white}Fuerte
& 60\% & 79\%\\
\hline

\cellcolor{rangoMuyFuerte}\color{white}\bfseries Muy Fuerte
& 80\% & 100\%\\
\hline

\cellcolor{negroTabla}
\parbox[c][1.70cm][c]{\linewidth}{
\RaggedRight\color{white}\bfseries
Asociación del Motivador con el Negocio
}
&
\cellcolor{negroTabla}
\parbox[c][2.20cm][c]{\linewidth}{
    \RaggedRight\color{white}\bfseries
    Definido por:\\[0.20cm]
    Ejecutado por:
}
&
\parbox[c][2.20cm][c]{\linewidth}{
    \RaggedRight
    Cliente y equipo del proyecto\\[0.20cm]
    Equipo de desarrollo Grupo 4
}
\\\hline

\end{tabular}
\end{table}
\endgroup



\subsection{Restricciones de Tecnología}

El diseño e implementación de AccionesUD estará condicionado
por las restricciones tecnológicas, económicas, de seguridad
y de infraestructura establecidas en el Acta de Constitución,
así como por las decisiones arquitectónicas adoptadas
posteriormente por el equipo de desarrollo.

La solución deberá ser accesible mediante un entorno web
desplegado y contar con una base de datos relacional
normalizada, como mínimo, hasta la tercera forma normal
(3FN).

Se adopta una arquitectura de \textbf{Monolito Modular
con Arquitectura Hexagonal}, aprobada por el cliente,
con el propósito de favorecer la separación de
responsabilidades, la modularidad y el desacoplamiento
entre la lógica de negocio y las integraciones externas.

El Acta de Constitución contempla inicialmente Python
con FastAPI, React y/o Astro, PostgreSQL, Docker y Render
como tecnologías propuestas. Posteriormente, el equipo
definió \textbf{Next.js para el frontend, FastAPI (Python)
para el backend y PostgreSQL para la persistencia}.
La utilización de Docker y Render deberá confirmarse
mediante las actividades de configuración y validación
de infraestructura.

Para la consulta de cotizaciones, datos de mercado y
gestión de operaciones bursátiles habilitadas, se
seleccionó \textbf{Alpaca como proveedor único de
servicios bursátiles}. La integración estará sujeta
a sus permisos, límites de consumo, cobertura de
instrumentos, disponibilidad, costos y condiciones
de uso. No se garantizará el acceso a mercados o
funcionalidades que no estén soportados por el proveedor
y por el entorno de integración implementado.

La plataforma deberá incorporar mecanismos de
autenticación multifactor (MFA), autorización basada
en roles, protección de credenciales, trazabilidad
de operaciones, registros de auditoría y monitoreo.
Asimismo, el tratamiento de datos personales y
financieros deberá considerar las obligaciones
legales aplicables, incluida la Ley 1581 de 2012
en Colombia, cuando corresponda.

El despliegue deberá ajustarse a los recursos técnicos
y económicos disponibles. La utilización de servicios
de infraestructura adicionales, como AWS, estará
condicionada a una necesidad justificada, la evaluación
de costos y la viabilidad de implementación.

Finalmente, los atributos de calidad relacionados
con rendimiento, disponibilidad, concurrencia y
seguridad deberán verificarse mediante pruebas
documentadas. Los resultados obtenidos en entornos
de desarrollo o pruebas no constituirán, por sí solos,
evidencia de cumplimiento en producción.


\subsection{Restricciones de Negocio}

El alcance de AccionesUD estará limitado por el
calendario académico, las fechas de entrega y los
productos establecidos para el proyecto. Cualquier
modificación de requerimientos o incorporación de
funcionalidades deberá evaluarse según su impacto
en el cronograma, los costos y los objetivos acordados.

Las operaciones bursátiles deberán respetar los
horarios de negociación, días no hábiles, instrumentos
disponibles y estados de las órdenes, de acuerdo con
las reglas de negocio y las capacidades de Alpaca.

Las compras requerirán la validación de fondos
suficientes, mientras que las ventas exigirán
verificar la disponibilidad de las posiciones
correspondientes. Los saldos, movimientos y
autorizaciones deberán mantener coherencia con
los registros internos y con los servicios externos,
según corresponda.

Los comisionistas únicamente podrán consultar
información de inversionistas con quienes mantengan
una relación autorizada. Las operaciones deberán
conservar registros que permitan verificar sus
estados, movimientos y cargos asociados.

Las comisiones deberán ser transparentes y
parametrizables. Como referencia preliminar, se
contempla una comisión del \textbf{2\%} sobre el
valor de la transacción, distribuida en un
\textbf{60\% para AccionesUD} y un
\textbf{40\% para el comisionista}, sujeta a
confirmación comercial.

Asimismo, se contemplan suscripciones premium
con tarifas preliminares de \textbf{USD 12 mensuales}
o \textbf{USD 120 anuales}, cuya implementación
y mecanismo de cobro deberán validarse.

La operación y evolución del sistema deberán
considerar los costos de infraestructura, dominio,
soporte y mantenimiento, evitando comprometer
funcionalidades o niveles de servicio que excedan
los recursos aprobados.

El cumplimiento de los objetivos funcionales
y atributos de calidad deberá sustentarse mediante
pruebas, evidencias y criterios de aceptación
definidos para cada entrega.


\section{Registro de Restricciones de Tecnología y Negocio}
\label{an:restricciones}

Las restricciones siguientes se derivan del Acta de Constitución,
los requerimientos del proyecto y las decisiones documentadas
posteriormente. Las alternativas describen opciones consideradas
o estrategias de cumplimiento, sin invalidar las restricciones.

\begingroup
\small
\setlength{\tabcolsep}{7pt}
\renewcommand{\arraystretch}{1.30}
\newlength{\anchoRegistro}
\setlength{\anchoRegistro}{\dimexpr\textwidth-4\tabcolsep-3\arrayrulewidth\relax}
\begin{table}[H]
\centering
\caption{Restricción tecnológica RT-01: Arquitectura de software seleccionada.}
\label{tab:restriccion-rt-01}
\begin{tabular}{|
>{\RaggedRight\arraybackslash}p{.28\anchoRegistro}|
>{\RaggedRight\arraybackslash}p{.72\anchoRegistro}|}
\hline
\rowcolor{negroTabla}
\multicolumn{2}{|l|}{\color{white}\bfseries Registro de Restricción Tecnológica}\\\hline
\cellcolor{azulTabla}\bfseries ID Restricción
& RT-01\\\hline
\cellcolor{azulTabla}\bfseries Tipo
& Tecnología\\\hline
\cellcolor{azulTabla}\bfseries Nombre
& Arquitectura de software seleccionada\\\hline
\cellcolor{azulTabla}\bfseries Descripción
& El sistema se organizará como Monolito Modular con Arquitectura Hexagonal, separando responsabilidades y aislando la lógica de negocio de los adaptadores de persistencia, interfaz e integración externa.\\\hline
\cellcolor{azulTabla}\bfseries Establecida por
& Cliente y equipo de desarrollo; decisión documentada en el acta del 31 de octubre de 2026.\\\hline
\cellcolor{azulTabla}\bfseries Alternativas
& Microservicios y arquitectura serverless, consideradas durante la evaluación arquitectónica.\\\hline
\cellcolor{azulTabla}\bfseries Observaciones
& La decisión deberá reflejarse en los diagramas de componentes, despliegue y organización modular del SAD.\\\hline
\end{tabular}
\end{table}
\endgroup

\begingroup
\small
\setlength{\tabcolsep}{7pt}
\renewcommand{\arraystretch}{1.30}
\setlength{\anchoRegistro}{\dimexpr\textwidth-4\tabcolsep-3\arrayrulewidth\relax}
\begin{table}[H]
\centering
\caption{Restricción tecnológica RT-02: Tecnologías de desarrollo.}
\label{tab:restriccion-rt-02}
\begin{tabular}{|
>{\RaggedRight\arraybackslash}p{.28\anchoRegistro}|
>{\RaggedRight\arraybackslash}p{.72\anchoRegistro}|}
\hline
\rowcolor{negroTabla}
\multicolumn{2}{|l|}{\color{white}\bfseries Registro de Restricción Tecnológica}\\\hline
\cellcolor{azulTabla}\bfseries ID Restricción
& RT-02\\\hline
\cellcolor{azulTabla}\bfseries Tipo
& Tecnología\\\hline
\cellcolor{azulTabla}\bfseries Nombre
& Tecnologías de desarrollo\\\hline
\cellcolor{azulTabla}\bfseries Descripción
& La solución empleará Next.js para el frontend, FastAPI (Python) para el backend y PostgreSQL para la persistencia, de acuerdo con la selección técnica posterior a la propuesta inicial del Acta de Constitución.\\\hline
\cellcolor{azulTabla}\bfseries Establecida por
& Equipo de desarrollo; propuesta inicial del Acta de Constitución y selección técnica posterior.\\\hline
\cellcolor{azulTabla}\bfseries Alternativas
& La propuesta preliminar contemplaba React y/o Astro para frontend; las tecnologías definitivas se documentan conforme a las decisiones del equipo.\\\hline
\cellcolor{azulTabla}\bfseries Observaciones
& Distinguir la propuesta preliminar de la selección posterior. Docker y Render no deben presentarse como aprobados sin confirmación adicional.\\\hline
\end{tabular}
\end{table}
\endgroup

\begingroup
\small
\setlength{\tabcolsep}{7pt}
\renewcommand{\arraystretch}{1.30}
\setlength{\anchoRegistro}{\dimexpr\textwidth-4\tabcolsep-3\arrayrulewidth\relax}
\begin{table}[H]
\centering
\caption{Restricción tecnológica RT-03: Proveedor único de servicios bursátiles.}
\label{tab:restriccion-rt-03}
\begin{tabular}{|
>{\RaggedRight\arraybackslash}p{.28\anchoRegistro}|
>{\RaggedRight\arraybackslash}p{.72\anchoRegistro}|}
\hline
\rowcolor{negroTabla}
\multicolumn{2}{|l|}{\color{white}\bfseries Registro de Restricción Tecnológica}\\\hline
\cellcolor{azulTabla}\bfseries ID Restricción
& RT-03\\\hline
\cellcolor{azulTabla}\bfseries Tipo
& Tecnología\\\hline
\cellcolor{azulTabla}\bfseries Nombre
& Proveedor único de servicios bursátiles\\\hline
\cellcolor{azulTabla}\bfseries Descripción
& La consulta de cotizaciones y datos de mercado, así como la gestión de operaciones bursátiles habilitadas, se integrarán mediante Alpaca como proveedor único de servicios bursátiles.\\\hline
\cellcolor{azulTabla}\bfseries Establecida por
& Equipo de desarrollo; decisión documentada en el acta del 5 de noviembre de 2026.\\\hline
\cellcolor{azulTabla}\bfseries Alternativas
& Integración con varios proveedores bursátiles o de datos financieros, evaluada frente a costos y complejidad.\\\hline
\cellcolor{azulTabla}\bfseries Observaciones
& Verificar cobertura de instrumentos y mercados, permisos, límites de consumo, costos, condiciones de uso, disponibilidad y capacidades del entorno de pruebas antes de comprometer funcionalidades.\\\hline
\end{tabular}
\end{table}
\endgroup

\begingroup
\small
\setlength{\tabcolsep}{7pt}
\renewcommand{\arraystretch}{1.30}
\setlength{\anchoRegistro}{\dimexpr\textwidth-4\tabcolsep-3\arrayrulewidth\relax}
\begin{table}[H]
\centering
\caption{Restricción tecnológica RT-04: Persistencia y normalización de datos.}
\label{tab:restriccion-rt-04}
\begin{tabular}{|
>{\RaggedRight\arraybackslash}p{.28\anchoRegistro}|
>{\RaggedRight\arraybackslash}p{.72\anchoRegistro}|}
\hline
\rowcolor{negroTabla}
\multicolumn{2}{|l|}{\color{white}\bfseries Registro de Restricción Tecnológica}\\\hline
\cellcolor{azulTabla}\bfseries ID Restricción
& RT-04\\\hline
\cellcolor{azulTabla}\bfseries Tipo
& Tecnología\\\hline
\cellcolor{azulTabla}\bfseries Nombre
& Persistencia y normalización de datos\\\hline
\cellcolor{azulTabla}\bfseries Descripción
& La información persistente del sistema se almacenará en una base de datos relacional cuyo esquema deberá cumplir, como mínimo, la tercera forma normal (3FN).\\\hline
\cellcolor{azulTabla}\bfseries Establecida por
& Requerimientos del proyecto recogidos en el Acta de Constitución y documentación del SAD.\\\hline
\cellcolor{azulTabla}\bfseries Alternativas
& Distintas estrategias de modelado relacional y optimización, siempre que respeten el nivel mínimo de normalización exigido.\\\hline
\cellcolor{azulTabla}\bfseries Observaciones
& PostgreSQL es el gestor seleccionado. Documentar el modelo entidad-relación, las restricciones de integridad y la justificación de posibles decisiones de rendimiento.\\\hline
\end{tabular}
\end{table}
\endgroup

\begingroup
\small
\setlength{\tabcolsep}{7pt}
\renewcommand{\arraystretch}{1.30}
\setlength{\anchoRegistro}{\dimexpr\textwidth-4\tabcolsep-3\arrayrulewidth\relax}
\begin{table}[H]
\centering
\caption{Restricción tecnológica RT-05: Seguridad, privacidad y auditoría.}
\label{tab:restriccion-rt-05}
\begin{tabular}{|
>{\RaggedRight\arraybackslash}p{.28\anchoRegistro}|
>{\RaggedRight\arraybackslash}p{.72\anchoRegistro}|}
\hline
\rowcolor{negroTabla}
\multicolumn{2}{|l|}{\color{white}\bfseries Registro de Restricción Tecnológica}\\\hline
\cellcolor{azulTabla}\bfseries ID Restricción
& RT-05\\\hline
\cellcolor{azulTabla}\bfseries Tipo
& Tecnología\\\hline
\cellcolor{azulTabla}\bfseries Nombre
& Seguridad, privacidad y auditoría\\\hline
\cellcolor{azulTabla}\bfseries Descripción
& El sistema deberá proteger credenciales, datos personales e información financiera mediante controles de acceso por roles, autenticación multifactor, mecanismos de protección de datos y registros de auditoría.\\\hline
\cellcolor{azulTabla}\bfseries Establecida por
& Acta de Constitución y requerimientos de seguridad del proyecto.\\\hline
\cellcolor{azulTabla}\bfseries Alternativas
& Mecanismos de autenticación, gestión de secretos, centralización de registros y controles de protección compatibles con la arquitectura definida.\\\hline
\cellcolor{azulTabla}\bfseries Observaciones
& Considerar las obligaciones legales aplicables, incluida la Ley 1581 de 2012 en Colombia cuando corresponda. La selección de controles deberá validarse mediante pruebas.\\\hline
\end{tabular}
\end{table}
\endgroup

\begingroup
\small
\setlength{\tabcolsep}{7pt}
\renewcommand{\arraystretch}{1.30}
\setlength{\anchoRegistro}{\dimexpr\textwidth-4\tabcolsep-3\arrayrulewidth\relax}
\begin{table}[H]
\centering
\caption{Restricción tecnológica RT-06: Infraestructura y despliegue web.}
\label{tab:restriccion-rt-06}
\begin{tabular}{|
>{\RaggedRight\arraybackslash}p{.28\anchoRegistro}|
>{\RaggedRight\arraybackslash}p{.72\anchoRegistro}|}
\hline
\rowcolor{negroTabla}
\multicolumn{2}{|l|}{\color{white}\bfseries Registro de Restricción Tecnológica}\\\hline
\cellcolor{azulTabla}\bfseries ID Restricción
& RT-06\\\hline
\cellcolor{azulTabla}\bfseries Tipo
& Tecnología\\\hline
\cellcolor{azulTabla}\bfseries Nombre
& Infraestructura y despliegue web\\\hline
\cellcolor{azulTabla}\bfseries Descripción
& La plataforma deberá estar disponible en un entorno web desplegado y no depender exclusivamente de localhost. La infraestructura se seleccionará de acuerdo con los recursos técnicos y económicos disponibles.\\\hline
\cellcolor{azulTabla}\bfseries Establecida por
& Acta de Constitución y decisiones de despliegue del equipo.\\\hline
\cellcolor{azulTabla}\bfseries Alternativas
& Docker y Render como opciones preliminares; AWS u otra infraestructura solo ante una necesidad justificada y recursos suficientes.\\\hline
\cellcolor{azulTabla}\bfseries Observaciones
& No asumir una migración a AWS ni un proveedor definitivo de alojamiento. Evaluar costos, configuración, mantenimiento y capacidad de crecimiento antes de aprobar el despliegue.\\\hline
\end{tabular}
\end{table}
\endgroup

\begingroup
\small
\setlength{\tabcolsep}{7pt}
\renewcommand{\arraystretch}{1.30}
\setlength{\anchoRegistro}{\dimexpr\textwidth-4\tabcolsep-3\arrayrulewidth\relax}
\begin{table}[H]
\centering
\caption{Restricción tecnológica RT-07: Validación de atributos de calidad.}
\label{tab:restriccion-rt-07}
\begin{tabular}{|
>{\RaggedRight\arraybackslash}p{.28\anchoRegistro}|
>{\RaggedRight\arraybackslash}p{.72\anchoRegistro}|}
\hline
\rowcolor{negroTabla}
\multicolumn{2}{|l|}{\color{white}\bfseries Registro de Restricción Tecnológica}\\\hline
\cellcolor{azulTabla}\bfseries ID Restricción
& RT-07\\\hline
\cellcolor{azulTabla}\bfseries Tipo
& Tecnología\\\hline
\cellcolor{azulTabla}\bfseries Nombre
& Validación de atributos de calidad\\\hline
\cellcolor{azulTabla}\bfseries Descripción
& El cumplimiento de desempeño, disponibilidad, concurrencia y seguridad deberá sustentarse en pruebas y evidencias verificables, con criterios de aceptación definidos para cada atributo.\\\hline
\cellcolor{azulTabla}\bfseries Establecida por
& Acta de Constitución del proyecto AccionesUD.\\\hline
\cellcolor{azulTabla}\bfseries Alternativas
& Pruebas de carga, concurrencia, integración, seguridad y monitoreo, seleccionadas según el atributo evaluado.\\\hline
\cellcolor{azulTabla}\bfseries Observaciones
& Los resultados de desarrollo o pruebas no garantizan el mismo comportamiento en producción. No declarar metas cumplidas sin mediciones y evidencias.\\\hline
\end{tabular}
\end{table}
\endgroup

\begingroup
\small
\setlength{\tabcolsep}{7pt}
\renewcommand{\arraystretch}{1.30}
\setlength{\anchoRegistro}{\dimexpr\textwidth-4\tabcolsep-3\arrayrulewidth\relax}
\begin{table}[H]
\centering
\caption{Restricción de negocio RN-01: Tiempo, alcance y control de cambios.}
\label{tab:restriccion-rn-01}
\begin{tabular}{|
>{\RaggedRight\arraybackslash}p{.28\anchoRegistro}|
>{\RaggedRight\arraybackslash}p{.72\anchoRegistro}|}
\hline
\rowcolor{negroTabla}
\multicolumn{2}{|l|}{\color{white}\bfseries Registro de Restricción de Negocio}\\\hline
\cellcolor{azulTabla}\bfseries ID Restricción
& RN-01\\\hline
\cellcolor{azulTabla}\bfseries Tipo
& Negocio\\\hline
\cellcolor{azulTabla}\bfseries Nombre
& Tiempo, alcance y control de cambios\\\hline
\cellcolor{azulTabla}\bfseries Descripción
& El desarrollo deberá respetar el calendario académico, los hitos de entrega y el alcance aprobado. Los cambios de requerimientos y las funcionalidades adicionales estarán sujetos a evaluación de impacto.\\\hline
\cellcolor{azulTabla}\bfseries Establecida por
& Acta de Constitución y planificación del proyecto.\\\hline
\cellcolor{azulTabla}\bfseries Alternativas
& Priorización del backlog, implementación incremental, reprogramación de funciones no críticas y control formal de cambios.\\\hline
\cellcolor{azulTabla}\bfseries Observaciones
& Relacionar la restricción con la planificación por sprints y con las decisiones que reducen complejidad técnica.\\\hline
\end{tabular}
\end{table}
\endgroup

\begingroup
\small
\setlength{\tabcolsep}{7pt}
\renewcommand{\arraystretch}{1.30}
\setlength{\anchoRegistro}{\dimexpr\textwidth-4\tabcolsep-3\arrayrulewidth\relax}
\begin{table}[H]
\centering
\caption{Restricción de negocio RN-02: Condiciones de negociación bursátil.}
\label{tab:restriccion-rn-02}
\begin{tabular}{|
>{\RaggedRight\arraybackslash}p{.28\anchoRegistro}|
>{\RaggedRight\arraybackslash}p{.72\anchoRegistro}|}
\hline
\rowcolor{negroTabla}
\multicolumn{2}{|l|}{\color{white}\bfseries Registro de Restricción de Negocio}\\\hline
\cellcolor{azulTabla}\bfseries ID Restricción
& RN-02\\\hline
\cellcolor{azulTabla}\bfseries Tipo
& Negocio\\\hline
\cellcolor{azulTabla}\bfseries Nombre
& Condiciones de negociación bursátil\\\hline
\cellcolor{azulTabla}\bfseries Descripción
& La consulta y eventual gestión de órdenes deberá respetar horarios de mercado, días no hábiles, instrumentos disponibles, estados de órdenes y demás reglas soportadas por el proveedor externo.\\\hline
\cellcolor{azulTabla}\bfseries Establecida por
& Acta de Constitución, requerimientos de negocio y decisión de integración con Alpaca.\\\hline
\cellcolor{azulTabla}\bfseries Alternativas
& Limitar las operaciones a los mercados e instrumentos efectivamente habilitados; emplear simulación cuando el alcance de pruebas así lo establezca.\\\hline
\cellcolor{azulTabla}\bfseries Observaciones
& La ejecución real depende de permisos, condiciones y modalidades habilitadas por Alpaca. No prometer mercados internacionales sin validación de cobertura.\\\hline
\end{tabular}
\end{table}
\endgroup

\begingroup
\small
\setlength{\tabcolsep}{7pt}
\renewcommand{\arraystretch}{1.30}
\setlength{\anchoRegistro}{\dimexpr\textwidth-4\tabcolsep-3\arrayrulewidth\relax}
\begin{table}[H]
\centering
\caption{Restricción de negocio RN-03: Validación de fondos, posiciones y órdenes.}
\label{tab:restriccion-rn-03}
\begin{tabular}{|
>{\RaggedRight\arraybackslash}p{.28\anchoRegistro}|
>{\RaggedRight\arraybackslash}p{.72\anchoRegistro}|}
\hline
\rowcolor{negroTabla}
\multicolumn{2}{|l|}{\color{white}\bfseries Registro de Restricción de Negocio}\\\hline
\cellcolor{azulTabla}\bfseries ID Restricción
& RN-03\\\hline
\cellcolor{azulTabla}\bfseries Tipo
& Negocio\\\hline
\cellcolor{azulTabla}\bfseries Nombre
& Validación de fondos, posiciones y órdenes\\\hline
\cellcolor{azulTabla}\bfseries Descripción
& Una compra requerirá comprobar fondos suficientes y una venta exigirá verificar la cantidad disponible del instrumento. Las autorizaciones, saldos y estados deberán mantenerse coherentes con las reglas de negocio y los servicios externos.\\\hline
\cellcolor{azulTabla}\bfseries Establecida por
& Acta de Constitución y requerimientos funcionales de operación financiera.\\\hline
\cellcolor{azulTabla}\bfseries Alternativas
& Validaciones previas en la aplicación y conciliación con la información disponible del proveedor, según las capacidades habilitadas.\\\hline
\cellcolor{azulTabla}\bfseries Observaciones
& Definir el origen de verdad de saldos y posiciones, el manejo de rechazos y la sincronización de estados antes de implementar operaciones.\\\hline
\end{tabular}
\end{table}
\endgroup

\begingroup
\small
\setlength{\tabcolsep}{7pt}
\renewcommand{\arraystretch}{1.30}
\setlength{\anchoRegistro}{\dimexpr\textwidth-4\tabcolsep-3\arrayrulewidth\relax}
\begin{table}[H]
\centering
\caption{Restricción de negocio RN-04: Acceso autorizado de comisionistas.}
\label{tab:restriccion-rn-04}
\begin{tabular}{|
>{\RaggedRight\arraybackslash}p{.28\anchoRegistro}|
>{\RaggedRight\arraybackslash}p{.72\anchoRegistro}|}
\hline
\rowcolor{negroTabla}
\multicolumn{2}{|l|}{\color{white}\bfseries Registro de Restricción de Negocio}\\\hline
\cellcolor{azulTabla}\bfseries ID Restricción
& RN-04\\\hline
\cellcolor{azulTabla}\bfseries Tipo
& Negocio\\\hline
\cellcolor{azulTabla}\bfseries Nombre
& Acceso autorizado de comisionistas\\\hline
\cellcolor{azulTabla}\bfseries Descripción
& Cada comisionista solo podrá consultar la información y las operaciones de inversionistas con quienes mantenga una relación autorizada, conforme a los permisos asignados.\\\hline
\cellcolor{azulTabla}\bfseries Establecida por
& Requerimientos de negocio y de control de acceso de AccionesUD.\\\hline
\cellcolor{azulTabla}\bfseries Alternativas
& Autorización por roles complementada con validaciones de asociación entre comisionista e inversionista.\\\hline
\cellcolor{azulTabla}\bfseries Observaciones
& Verificar restricciones de acceso tanto en los servicios del backend como en las consultas de datos; conservar trazabilidad de accesos relevantes.\\\hline
\end{tabular}
\end{table}
\endgroup

\begingroup
\small
\setlength{\tabcolsep}{7pt}
\renewcommand{\arraystretch}{1.30}
\setlength{\anchoRegistro}{\dimexpr\textwidth-4\tabcolsep-3\arrayrulewidth\relax}
\begin{table}[H]
\centering
\caption{Restricción de negocio RN-05: Comisiones y suscripciones parametrizables.}
\label{tab:restriccion-rn-05}
\begin{tabular}{|
>{\RaggedRight\arraybackslash}p{.28\anchoRegistro}|
>{\RaggedRight\arraybackslash}p{.72\anchoRegistro}|}
\hline
\rowcolor{negroTabla}
\multicolumn{2}{|l|}{\color{white}\bfseries Registro de Restricción de Negocio}\\\hline
\cellcolor{azulTabla}\bfseries ID Restricción
& RN-05\\\hline
\cellcolor{azulTabla}\bfseries Tipo
& Negocio\\\hline
\cellcolor{azulTabla}\bfseries Nombre
& Comisiones y suscripciones parametrizables\\\hline
\cellcolor{azulTabla}\bfseries Descripción
& La plataforma deberá calcular y presentar los cargos de manera transparente y permitir parametrizar comisiones y planes premium. Se contemplan como referencias una comisión del 2\%, distribución 60\%/40\% y planes de USD 12 mensuales o USD 120 anuales.\\\hline
\cellcolor{azulTabla}\bfseries Establecida por
& Enunciado y condiciones comerciales preliminares del proyecto; sujetas a confirmación del cliente.\\\hline
\cellcolor{azulTabla}\bfseries Alternativas
& Configuración administrativa de tarifas, porcentajes y planes; integración de cobro de suscripciones en un entorno de pruebas si se valida Stripe.\\\hline
\cellcolor{azulTabla}\bfseries Observaciones
& Los valores indicados no son tarifas invariables ni deben considerarse aprobados definitivamente. Informar cargos antes de operar y conservar su desglose.\\\hline
\end{tabular}
\end{table}
\endgroup

\begingroup
\small
\setlength{\tabcolsep}{7pt}
\renewcommand{\arraystretch}{1.30}
\setlength{\anchoRegistro}{\dimexpr\textwidth-4\tabcolsep-3\arrayrulewidth\relax}
\begin{table}[H]
\centering
\caption{Restricción de negocio RN-06: Presupuesto y sostenibilidad operativa.}
\label{tab:restriccion-rn-06}
\begin{tabular}{|
>{\RaggedRight\arraybackslash}p{.28\anchoRegistro}|
>{\RaggedRight\arraybackslash}p{.72\anchoRegistro}|}
\hline
\rowcolor{negroTabla}
\multicolumn{2}{|l|}{\color{white}\bfseries Registro de Restricción de Negocio}\\\hline
\cellcolor{azulTabla}\bfseries ID Restricción
& RN-06\\\hline
\cellcolor{azulTabla}\bfseries Tipo
& Negocio\\\hline
\cellcolor{azulTabla}\bfseries Nombre
& Presupuesto y sostenibilidad operativa\\\hline
\cellcolor{azulTabla}\bfseries Descripción
& La implementación y operación deberán ajustarse a los recursos aprobados y considerar costos de infraestructura, integraciones, dominio, soporte y mantenimiento.\\\hline
\cellcolor{azulTabla}\bfseries Establecida por
& Acta de Constitución y evaluación de viabilidad económica del proyecto.\\\hline
\cellcolor{azulTabla}\bfseries Alternativas
& Optimización de recursos, escalamiento controlado y contratación de servicios adicionales solo cuando exista justificación técnica y presupuestal.\\\hline
\cellcolor{azulTabla}\bfseries Observaciones
& Evaluar el costo total de operación antes de ampliar infraestructura o integrar nuevos servicios; evitar comprometer niveles de servicio no financiados.\\\hline
\end{tabular}
\end{table}
\endgroup


\subsection{Atributos de Calidad}
\subsubsection{Recabación y Justificación de Requerimientos No Funcionales}
\begin{itemize}
    \item \textbf{RNF-01. Rendimiento:} el enunciado plantea una capacidad de hasta \textbf{5.000 usuarios} y un objetivo de atención de \textbf{1.500 concurrentes} con tiempos de respuesta que no excedan \textbf{2 segundos}. Se debe acordar cuáles transacciones, percentiles y condiciones de carga se medirán.
    \item \textbf{RNF-02. Tiempo de confirmación:} para el escenario de compra se menciona un plazo inferior a \textbf{5 segundos}. Falta definir si se refiere a recepción, validación, aceptación del proveedor o ejecución bursátil final; no deben considerarse equivalentes.
    \item \textbf{RNF-03. Disponibilidad:} el servicio debe priorizar continuidad durante los horarios de los mercados. No se proporciona un porcentaje contractual de disponibilidad ni una política de recuperación.
    \item \textbf{RNF-04. Seguridad:} emplear MFA, separación de privilegios y protección de sesiones y datos financieros. La estrategia concreta de cifrado, rotación y gestión de secretos está por diseñarse.
    \item \textbf{RNF-05. Auditabilidad:} conservar quién origina, revisa, firma, envía, cancela y recibe respuesta de cada orden, con tiempos y cambios de estado verificables.
    \item \textbf{RNF-06. Interoperabilidad:} aislar la integración con fuentes de mercado, ejecución, fondos, pasarela de pagos y canales de notificación para manejar sus diferentes contratos y fallos.
    \item \textbf{RNF-07. Confiabilidad e integridad:} evitar inconsistencias entre saldos, posiciones, órdenes y respuestas del tercero, sin presentar una confirmación local como ejecución final.
    \item \textbf{RNF-08. Observabilidad:} centralizar registros y facilitar monitoreo, alertamiento y análisis de fallas y eventos sospechosos.
    \item \textbf{RNF-09. Usabilidad:} mantener una interfaz web responsiva con información comprensible para usuarios principiantes y experimentados.
    \item \textbf{RNF-10. Mantenibilidad:} organizar el diseño para modificar reglas, mercados, proveedores y módulos sin afectar innecesariamente funciones no relacionadas.
\end{itemize}

\subsubsection{Escenarios de Calidad}
Los siguientes escenarios concretan los requerimientos no funcionales mediante fuente, estímulo, artefacto, ambiente, respuesta y medida de la respuesta. Los umbrales aún no aprobados se identifican como pendientes de validación; ningún escenario se considera satisfecho sin evidencias de prueba. La disponibilidad de Alpaca y de otros servicios externos condiciona los resultados de integración.

\begingroup
\small
\setlength{\tabcolsep}{5pt}
\renewcommand{\arraystretch}{1.25}
\newlength{\anchoEscenario}
\setlength{\anchoEscenario}{\dimexpr\textwidth-4\tabcolsep-3\arrayrulewidth\relax}
\newcommand{\campoEscenario}[2]{\cellcolor{azulTabla}\bfseries #1 & #2\\\hline}
\newcommand{\inicioEscenario}[3]{%
\begin{table}[H]\centering
\caption{Escenario de calidad #1: #2.}\label{tab:escenario-#1}
\begin{tabular}{|>{\RaggedRight\arraybackslash}p{.28\anchoEscenario}|>{\RaggedRight\arraybackslash}p{.72\anchoEscenario}|}\hline
\rowcolor{negroTabla}\multicolumn{2}{|l|}{\color{white}\bfseries Registro de Escenario de Calidad}\\\hline
\campoEscenario{Escenario de Calidad \#}{#1}
\campoEscenario{Stakeholder}{#3}
}
\newcommand{\finEscenario}{\end{tabular}\end{table}}

\inicioEscenario{EC-01}{Rendimiento bajo concurrencia}{Inversionistas, comisionistas y equipo de operaciones}
\campoEscenario{Atributo de Calidad}{Rendimiento (RNF-01).}
\campoEscenario{Justificación}{El enunciado plantea hasta 5.000 usuarios y una meta de 1.500 concurrentes con respuestas de hasta 2 segundos. La capacidad debe comprobarse con pruebas de carga representativas.}
\campoEscenario{Fuente}{Usuarios concurrentes y herramienta de pruebas de carga.}
\campoEscenario{Estímulo}{1.500 sesiones concurrentes solicitan cotizaciones y consultan información de su portafolio.}
\campoEscenario{Artefacto}{Frontend Next.js, API FastAPI, módulos de consulta y persistencia PostgreSQL, y adaptador de datos de Alpaca.}
\campoEscenario{Ambiente}{Entorno de pruebas con infraestructura y datos representativos; dependencias externas simuladas o controladas cuando sea necesario.}
\campoEscenario{Respuesta}{El sistema atiende solicitudes, controla saturación y errores, y registra métricas de latencia y capacidad.}
\campoEscenario{Medida de la Respuesta}{Objetivo del enunciado: tiempo de respuesta no superior a 2 s para las operaciones y el percentil que se acuerden, bajo 1.500 concurrentes. Definir antes de probar percentil, tasa de error, mezcla de transacciones y alcance de las llamadas externas.}
\finEscenario

\inicioEscenario{EC-02}{Confirmación de solicitud de compra}{Inversionista y equipo de operaciones}
\campoEscenario{Atributo de Calidad}{Tiempo de confirmación (RNF-02).}
\campoEscenario{Justificación}{El enunciado menciona un tiempo inferior a 5 segundos, pero no precisa la etapa de confirmación.}
\campoEscenario{Fuente}{Inversionista autenticado.}
\campoEscenario{Estímulo}{El usuario envía una solicitud de compra con instrumento, cantidad y parámetros válidos.}
\campoEscenario{Artefacto}{Módulo de órdenes, validación de fondos, adaptador de Alpaca y registro de estados.}
\campoEscenario{Ambiente}{Sesión válida, mercado y cuenta habilitados, con proveedor disponible y latencia medida.}
\campoEscenario{Respuesta}{Se valida la solicitud, se registra su identificador y se muestra un estado inequívoco: recibida, rechazada, aceptada por proveedor o ejecutada, según evidencia disponible.}
\campoEscenario{Medida de la Respuesta}{Objetivo preliminar: menos de 5 s para la etapa de confirmación que apruebe el cliente. Medir por separado recepción local, validación, aceptación del proveedor y ejecución final; esta última no se presupone dentro de los 5 s.}
\finEscenario

\inicioEscenario{EC-03}{Interrupción del proveedor bursátil}{Inversionistas, comisionistas y administrador}
\campoEscenario{Atributo de Calidad}{Disponibilidad y resiliencia (RNF-03).}
\campoEscenario{Justificación}{La continuidad durante el horario bursátil depende también de la disponibilidad de Alpaca; no existe aún un SLA aprobado.}
\campoEscenario{Fuente}{Fallo, demora o indisponibilidad de la API de Alpaca.}
\campoEscenario{Estímulo}{La API externa deja de responder durante la consulta de datos o el envío de una orden.}
\campoEscenario{Artefacto}{Adaptador de Alpaca, módulo de órdenes, interfaz y monitoreo.}
\campoEscenario{Ambiente}{Operación normal en horario de mercado con dependencia externa degradada.}
\campoEscenario{Respuesta}{El sistema identifica la falla, informa la indisponibilidad y conserva un estado verificable; evita indicar que una orden fue ejecutada sin confirmación del proveedor.}
\campoEscenario{Medida de la Respuesta}{Registrar tasa de errores, tiempo de detección, tiempo de recuperación y porcentaje de solicitudes con estado inequívoco. Umbrales y política de reintentos pendientes de aprobación.}
\finEscenario

\inicioEscenario{EC-04}{Acceso no autorizado a información financiera}{Inversionista, comisionista y responsable de seguridad}
\campoEscenario{Atributo de Calidad}{Seguridad y privacidad (RNF-04).}
\campoEscenario{Justificación}{La plataforma maneja información personal y financiera; debe restringir el acceso por rol y por relación autorizada con el cliente.}
\campoEscenario{Fuente}{Comisionista autenticado sin asociación con el inversionista consultado.}
\campoEscenario{Estímulo}{El comisionista intenta consultar, mediante URL o API, el portafolio de un inversionista no asociado.}
\campoEscenario{Artefacto}{Autorización en FastAPI, control de acceso a recursos y auditoría.}
\campoEscenario{Ambiente}{Sistema en operación con usuario autenticado y permisos limitados.}
\campoEscenario{Respuesta}{La solicitud se deniega sin revelar datos sensibles y el intento queda registrado para auditoría.}
\campoEscenario{Medida de la Respuesta}{El 100\% de los intentos no autorizados del conjunto de pruebas debe ser rechazado y no exponer datos; verificar registros de auditoría. Esta es una condición de aceptación propuesta, sujeta a validación.}
\finEscenario

\inicioEscenario{EC-05}{Reconstrucción de la historia de una orden}{Administrador, auditor y equipo de soporte}
\campoEscenario{Atributo de Calidad}{Auditabilidad (RNF-05).}
\campoEscenario{Justificación}{Las operaciones deben ser rastreables para explicar cambios de estado, decisiones y cargos.}
\campoEscenario{Fuente}{Administrador o auditor autorizado.}
\campoEscenario{Estímulo}{Se solicita investigar una orden disputada o con estados aparentemente inconsistentes.}
\campoEscenario{Artefacto}{Módulo de órdenes, bitácora de auditoría y base de datos.}
\campoEscenario{Ambiente}{Consulta posterior a la operación con registros históricos disponibles.}
\campoEscenario{Respuesta}{Se reconstruye la secuencia de eventos, actor, marcas de tiempo, cambios de estado y referencias de respuestas externas, protegiendo datos sensibles.}
\campoEscenario{Medida de la Respuesta}{Para los casos de prueba, trazabilidad completa de los eventos definidos como obligatorios y correlación verificable con el identificador de orden; plazo de retención y tiempo máximo de consulta pendientes de definir.}
\finEscenario

\inicioEscenario{EC-06}{Cambio o fallo de una integración externa}{Equipo de desarrollo e integración}
\campoEscenario{Atributo de Calidad}{Interoperabilidad (RNF-06).}
\campoEscenario{Justificación}{La integración bursátil con Alpaca y las integraciones independientes de pagos y notificaciones deben aislar contratos y errores.}
\campoEscenario{Fuente}{API externa o proveedor de servicios.}
\campoEscenario{Estímulo}{Alpaca cambia una respuesta compatible con su versionado o devuelve un error de límite de consumo.}
\campoEscenario{Artefacto}{Puerto de integración y adaptador de Alpaca de la arquitectura hexagonal.}
\campoEscenario{Ambiente}{Entorno de integración con respuestas simuladas y pruebas de contrato.}
\campoEscenario{Respuesta}{El adaptador interpreta la respuesta o produce un error controlado sin propagar formatos externos a la lógica de negocio.}
\campoEscenario{Medida de la Respuesta}{Las pruebas de contrato e integración definidas deben detectar incompatibilidades y comprobar que los módulos no dependientes siguen funcionando. Tiempo máximo de adaptación pendiente de acordar.}
\finEscenario

\inicioEscenario{EC-07}{Prevención de órdenes duplicadas}{Inversionista y responsable de operaciones}
\campoEscenario{Atributo de Calidad}{Confiabilidad e integridad (RNF-07).}
\campoEscenario{Justificación}{Las pérdidas de conexión y reintentos no deben producir dobles registros o confirmaciones financieras incorrectas.}
\campoEscenario{Fuente}{Usuario, red o cliente API que reintenta una solicitud.}
\campoEscenario{Estímulo}{La misma solicitud de compra se envía nuevamente tras un tiempo de espera sin respuesta visible.}
\campoEscenario{Artefacto}{Módulo de órdenes, persistencia, identificadores de solicitud y adaptador de Alpaca.}
\campoEscenario{Ambiente}{Fallo intermitente de red durante el procesamiento de una orden.}
\campoEscenario{Respuesta}{El sistema correlaciona las solicitudes, consulta el estado externo cuando proceda y evita duplicaciones dentro de las garantías disponibles; no presume ejecución sin evidencia.}
\campoEscenario{Medida de la Respuesta}{En los casos de prueba de reintento, cero órdenes duplicadas atribuibles al sistema y consistencia entre estado interno y respuesta del proveedor; validar las capacidades reales de idempotencia de la API.}
\finEscenario

\inicioEscenario{EC-08}{Detección y diagnóstico de una falla}{Administrador y equipo de soporte}
\campoEscenario{Atributo de Calidad}{Observabilidad (RNF-08).}
\campoEscenario{Justificación}{La operación necesita registros centralizados para detectar degradaciones y eventos sospechosos.}
\campoEscenario{Fuente}{Fallo de aplicación, dependencia o evento de seguridad.}
\campoEscenario{Estímulo}{Aumentan los errores de una ruta de órdenes o falla una dependencia.}
\campoEscenario{Artefacto}{API FastAPI, registro de eventos, métricas, alertas y monitoreo.}
\campoEscenario{Ambiente}{Sistema desplegado con instrumentación activa.}
\campoEscenario{Respuesta}{Se registran eventos correlacionables, se identifica el componente afectado y se notifica al responsable según reglas configuradas.}
\campoEscenario{Medida de la Respuesta}{Comprobar emisión de registros y métricas para los fallos previstos, correlación por identificador y generación de alertas; umbrales y tiempo de detección pendientes de definición.}
\finEscenario

\inicioEscenario{EC-09}{Consulta financiera desde dispositivo móvil}{Inversionista principiante o experimentado}
\campoEscenario{Atributo de Calidad}{Usabilidad (RNF-09).}
\campoEscenario{Justificación}{La plataforma web debe ofrecer información comprensible y accesible desde diferentes tamaños de pantalla.}
\campoEscenario{Fuente}{Usuario que accede desde navegador móvil.}
\campoEscenario{Estímulo}{Consulta cotizaciones, detalle de instrumentos y estados de órdenes en una pantalla pequeña.}
\campoEscenario{Artefacto}{Interfaz Next.js y componentes de presentación.}
\campoEscenario{Ambiente}{Navegador móvil compatible y conexión de red representativa.}
\campoEscenario{Respuesta}{La interfaz se adapta al tamaño de pantalla, mantiene legibles los estados y cargos, y permite completar el recorrido sin pérdida de funciones esenciales.}
\campoEscenario{Medida de la Respuesta}{Validar tareas críticas mediante pruebas de usabilidad y compatibilidad responsive; definir dispositivos, tareas, tasa de éxito y umbrales de satisfacción antes de la aceptación.}
\finEscenario

\inicioEscenario{EC-10}{Actualización de reglas de negociación}{Equipo de desarrollo y responsable funcional}
\campoEscenario{Atributo de Calidad}{Mantenibilidad (RNF-10).}
\campoEscenario{Justificación}{La arquitectura modular debe permitir ajustar reglas y calendarios sin afectar módulos no relacionados.}
\campoEscenario{Fuente}{Cambio funcional aprobado o actualización de calendario bursátil.}
\campoEscenario{Estímulo}{Se incorpora un día no hábil o se modifica una regla de validación de operaciones.}
\campoEscenario{Artefacto}{Módulo de órdenes, configuración de calendarios y pruebas automatizadas.}
\campoEscenario{Ambiente}{Entorno de desarrollo e integración con control de versiones.}
\campoEscenario{Respuesta}{Se realiza el cambio en el módulo responsable y se verifican sus efectos mediante pruebas unitarias, de integración y regresión.}
\campoEscenario{Medida de la Respuesta}{Registrar módulos modificados, tiempo de implementación y pruebas de regresión aprobadas. Definir umbrales de esfuerzo y acoplamiento según la línea base del proyecto.}
\finEscenario

\endgroup

% ============================================================================
\section{Contexto}
\subsection{Escenarios Operacionales}
El contexto operacional parte de los escenarios del enunciado. A nivel inicial se reconocen los siguientes flujos principales:
\begin{itemize}
    \item \textbf{Compra o venta por un inversionista:} el usuario consulta un activo, establece los datos de la orden, recibe las validaciones y confirmaciones que correspondan, y observa el resultado en el historial y el portafolio.
    \item \textbf{Asesoría mediante comisionista:} el inversionista se vincula con un especialista, revisa una propuesta y realiza las autorizaciones previstas antes de solicitar el envío de órdenes.
    \item \textbf{Operación fuera del horario del mercado:} el sistema identifica la restricción, informa al solicitante y, si existe consentimiento y soporte del tercero, programa el tratamiento posterior de la orden.
    \item \textbf{Supervisión y auditoría:} el personal autorizado consulta eventos, identifica comportamientos que requieren revisión y genera la evidencia correspondiente.
    \item \textbf{Gestión comercial:} un usuario contrata una suscripción premium y la plataforma actualiza sus derechos cuando recibe la confirmación del sistema de pagos.
\end{itemize}

\begingroup
\footnotesize
\setlength{\tabcolsep}{5pt}
\renewcommand{\arraystretch}{1.20}
\newlength{\anchoOperacional}
\setlength{\anchoOperacional}{\dimexpr\linewidth-8\tabcolsep-5\arrayrulewidth\relax}

\begin{table}[H]
\centering
\caption{Escenario operacional EO-01: Compra o venta de instrumentos bursátiles.}
\label{tab:escenario-eo-01}
\begin{tabular}{|
>{\RaggedRight\arraybackslash}p{.26\anchoOperacional}|
>{\RaggedRight\arraybackslash}p{.25\anchoOperacional}|
>{\RaggedRight\arraybackslash}p{.25\anchoOperacional}|
>{\RaggedRight\arraybackslash}p{.24\anchoOperacional}|}
\hline
\multicolumn{4}{|>{\columncolor{azulTabla}\RaggedRight\arraybackslash}p{\dimexpr\linewidth-2\tabcolsep-2\arrayrulewidth\relax}|}{\bfseries Título del Escenario Operacional} \\
\hline
\multicolumn{4}{|>{\RaggedRight\arraybackslash}p{\dimexpr\linewidth-2\tabcolsep-2\arrayrulewidth\relax}|}{\textbf{Compra o venta de instrumentos bursátiles}} \\
\hline
\cellcolor{azulTabla}\textbf{Stakeholder} & \multicolumn{1}{>{\RaggedRight\arraybackslash}p{.25\anchoOperacional}|}{Inversionista} & \cellcolor{azulTabla}\textbf{ID} & EO-01 \\
\hline

\cellcolor{azulTabla}\textbf{Descripción general de la funcionalidad} & \multicolumn{3}{>{\RaggedRight\arraybackslash}p{\dimexpr.74\anchoOperacional+4\tabcolsep+2\arrayrulewidth\relax}|}{Permitir que un inversionista consulte instrumentos disponibles y solicite operaciones de compra o venta mediante la plataforma.} \\
\hline

\cellcolor{azulTabla}\textbf{Describa lo que el stakeholder hace ahora o le gustaría poder hacer} & \multicolumn{3}{>{\RaggedRight\arraybackslash}p{\dimexpr.74\anchoOperacional+4\tabcolsep+2\arrayrulewidth\relax}|}{El inversionista desea consultar una cotización, seleccionar un instrumento, indicar cantidad y tipo de operación, revisar costos y enviar la solicitud; posteriormente consulta su estado.} \\
\hline

\cellcolor{azulTabla}\textbf{Describa cualquier entrada provista o disponible al momento del inicio} & \multicolumn{3}{>{\RaggedRight\arraybackslash}p{\dimexpr.74\anchoOperacional+4\tabcolsep+2\arrayrulewidth\relax}|}{Sesión autenticada; instrumento seleccionado; cotización disponible; tipo de orden; cantidad; fondos o posiciones; condiciones y cargos aplicables.} \\
\hline

\cellcolor{azulTabla}\textbf{Describa el contexto de la operación} & \multicolumn{3}{>{\RaggedRight\arraybackslash}p{\dimexpr.74\anchoOperacional+4\tabcolsep+2\arrayrulewidth\relax}|}{Interfaz web de AccionesUD durante un horario y mercado soportados por Alpaca. La información y ejecución dependen de los permisos y servicios habilitados.} \\
\hline

\cellcolor{azulTabla}\textbf{Describa cómo el sistema debe responder} & \multicolumn{3}{>{\RaggedRight\arraybackslash}p{\dimexpr.74\anchoOperacional+4\tabcolsep+2\arrayrulewidth\relax}|}{Validar autenticación, permisos, datos de la orden, fondos o posiciones y condiciones del mercado; mostrar cargos antes de confirmar; registrar la solicitud y enviarla a Alpaca cuando proceda; actualizar estados conforme a las respuestas recibidas.} \\
\hline

\cellcolor{azulTabla}\textbf{Describa las salidas que el sistema produce como resultado de la acción} & \multicolumn{3}{>{\RaggedRight\arraybackslash}p{\dimexpr.74\anchoOperacional+4\tabcolsep+2\arrayrulewidth\relax}|}{Resumen de costos; identificador de solicitud; resultado de validación; estado de recepción, aceptación, rechazo o ejecución, según corresponda; historial y portafolio actualizado únicamente con eventos confirmados.} \\
\hline

\cellcolor{azulTabla}\textbf{Describa quién o qué usa la salida y para qué es utilizada} & \multicolumn{3}{>{\RaggedRight\arraybackslash}p{\dimexpr.74\anchoOperacional+4\tabcolsep+2\arrayrulewidth\relax}|}{El inversionista consulta el seguimiento y sus posiciones; el administrador autorizado utiliza los registros para soporte y auditoría.} \\
\hline

\end{tabular}
\end{table}

\begin{table}[H]
\centering
\caption{Escenario operacional EO-02: Asesoría y autorización mediante comisionista.}
\label{tab:escenario-eo-02}
\begin{tabular}{|
>{\RaggedRight\arraybackslash}p{.26\anchoOperacional}|
>{\RaggedRight\arraybackslash}p{.25\anchoOperacional}|
>{\RaggedRight\arraybackslash}p{.25\anchoOperacional}|
>{\RaggedRight\arraybackslash}p{.24\anchoOperacional}|}
\hline
\multicolumn{4}{|>{\columncolor{azulTabla}\RaggedRight\arraybackslash}p{\dimexpr\linewidth-2\tabcolsep-2\arrayrulewidth\relax}|}{\bfseries Título del Escenario Operacional} \\
\hline
\multicolumn{4}{|>{\RaggedRight\arraybackslash}p{\dimexpr\linewidth-2\tabcolsep-2\arrayrulewidth\relax}|}{\textbf{Asesoría y autorización mediante comisionista}} \\
\hline
\cellcolor{azulTabla}\textbf{Stakeholder} & \multicolumn{1}{>{\RaggedRight\arraybackslash}p{.25\anchoOperacional}|}{Inversionista y comisionista} & \cellcolor{azulTabla}\textbf{ID} & EO-02 \\
\hline

\cellcolor{azulTabla}\textbf{Descripción general de la funcionalidad} & \multicolumn{3}{>{\RaggedRight\arraybackslash}p{\dimexpr.74\anchoOperacional+4\tabcolsep+2\arrayrulewidth\relax}|}{Facilitar la relación autorizada entre inversionista y comisionista para consultar información pertinente, formular propuestas y gestionar autorizaciones previstas.} \\
\hline

\cellcolor{azulTabla}\textbf{Describa lo que el stakeholder hace ahora o le gustaría poder hacer} & \multicolumn{3}{>{\RaggedRight\arraybackslash}p{\dimexpr.74\anchoOperacional+4\tabcolsep+2\arrayrulewidth\relax}|}{El inversionista vincula o autoriza a un comisionista, revisa una propuesta y decide aceptarla o rechazarla; el comisionista consulta exclusivamente los datos habilitados y presenta recomendaciones.} \\
\hline

\cellcolor{azulTabla}\textbf{Describa cualquier entrada provista o disponible al momento del inicio} & \multicolumn{3}{>{\RaggedRight\arraybackslash}p{\dimexpr.74\anchoOperacional+4\tabcolsep+2\arrayrulewidth\relax}|}{Identidad y rol de ambos actores; vínculo o autorización vigente; información de instrumentos y portafolio permitida; propuesta de asesoría; decisión del inversionista.} \\
\hline

\cellcolor{azulTabla}\textbf{Describa el contexto de la operación} & \multicolumn{3}{>{\RaggedRight\arraybackslash}p{\dimexpr.74\anchoOperacional+4\tabcolsep+2\arrayrulewidth\relax}|}{Sesiones web autenticadas y control de acceso por roles y relación autorizada. La capacidad de enviar órdenes en nombre de terceros depende del alcance funcional y de los permisos externos validados.} \\
\hline

\cellcolor{azulTabla}\textbf{Describa cómo el sistema debe responder} & \multicolumn{3}{>{\RaggedRight\arraybackslash}p{\dimexpr.74\anchoOperacional+4\tabcolsep+2\arrayrulewidth\relax}|}{Verificar vínculo y permisos; restringir la consulta a clientes asociados; registrar propuesta y consentimiento; impedir envíos sin autorización válida; aplicar las mismas validaciones financieras y externas que a una operación directa.} \\
\hline

\cellcolor{azulTabla}\textbf{Describa las salidas que el sistema produce como resultado de la acción} & \multicolumn{3}{>{\RaggedRight\arraybackslash}p{\dimexpr.74\anchoOperacional+4\tabcolsep+2\arrayrulewidth\relax}|}{Propuesta registrada; estado de autorización; notificaciones pertinentes; trazabilidad de accesos y decisiones; eventual referencia de solicitud de orden si la funcionalidad está habilitada.} \\
\hline

\cellcolor{azulTabla}\textbf{Describa quién o qué usa la salida y para qué es utilizada} & \multicolumn{3}{>{\RaggedRight\arraybackslash}p{\dimexpr.74\anchoOperacional+4\tabcolsep+2\arrayrulewidth\relax}|}{El inversionista revisa y autoriza; el comisionista da seguimiento a sus clientes asociados; auditoría utiliza el historial para verificar actuaciones.} \\
\hline

\end{tabular}
\end{table}

\begin{table}[H]
\centering
\caption{Escenario operacional EO-03: Solicitud de operación fuera del horario del mercado.}
\label{tab:escenario-eo-03}
\begin{tabular}{|
>{\RaggedRight\arraybackslash}p{.26\anchoOperacional}|
>{\RaggedRight\arraybackslash}p{.25\anchoOperacional}|
>{\RaggedRight\arraybackslash}p{.25\anchoOperacional}|
>{\RaggedRight\arraybackslash}p{.24\anchoOperacional}|}
\hline
\multicolumn{4}{|>{\columncolor{azulTabla}\RaggedRight\arraybackslash}p{\dimexpr\linewidth-2\tabcolsep-2\arrayrulewidth\relax}|}{\bfseries Título del Escenario Operacional} \\
\hline
\multicolumn{4}{|>{\RaggedRight\arraybackslash}p{\dimexpr\linewidth-2\tabcolsep-2\arrayrulewidth\relax}|}{\textbf{Solicitud de operación fuera del horario del mercado}} \\
\hline
\cellcolor{azulTabla}\textbf{Stakeholder} & \multicolumn{1}{>{\RaggedRight\arraybackslash}p{.25\anchoOperacional}|}{Inversionista} & \cellcolor{azulTabla}\textbf{ID} & EO-03 \\
\hline

\cellcolor{azulTabla}\textbf{Descripción general de la funcionalidad} & \multicolumn{3}{>{\RaggedRight\arraybackslash}p{\dimexpr.74\anchoOperacional+4\tabcolsep+2\arrayrulewidth\relax}|}{Informar y gestionar de forma controlada las solicitudes cuando el mercado o instrumento no admita negociación inmediata.} \\
\hline

\cellcolor{azulTabla}\textbf{Describa lo que el stakeholder hace ahora o le gustaría poder hacer} & \multicolumn{3}{>{\RaggedRight\arraybackslash}p{\dimexpr.74\anchoOperacional+4\tabcolsep+2\arrayrulewidth\relax}|}{El inversionista intenta comprar o vender fuera del horario habilitado y necesita conocer si puede conservar una solicitud para tratamiento posterior.} \\
\hline

\cellcolor{azulTabla}\textbf{Describa cualquier entrada provista o disponible al momento del inicio} & \multicolumn{3}{>{\RaggedRight\arraybackslash}p{\dimexpr.74\anchoOperacional+4\tabcolsep+2\arrayrulewidth\relax}|}{Instrumento; mercado; hora de la solicitud; calendario y estado de negociación; tipo de orden; consentimiento explícito para tratamiento posterior, si aplica.} \\
\hline

\cellcolor{azulTabla}\textbf{Describa el contexto de la operación} & \multicolumn{3}{>{\RaggedRight\arraybackslash}p{\dimexpr.74\anchoOperacional+4\tabcolsep+2\arrayrulewidth\relax}|}{Mercado cerrado, feriado o instrumento sin sesión habilitada; disponibilidad y reglas obtenidas o verificadas según los servicios de Alpaca.} \\
\hline

\cellcolor{azulTabla}\textbf{Describa cómo el sistema debe responder} & \multicolumn{3}{>{\RaggedRight\arraybackslash}p{\dimexpr.74\anchoOperacional+4\tabcolsep+2\arrayrulewidth\relax}|}{Detectar la condición; explicar por qué no puede ejecutarse inmediatamente; rechazar o registrar una solicitud pendiente solo si la modalidad está soportada y el inversionista la autoriza; volver a validar condiciones antes de cualquier envío posterior.} \\
\hline

\cellcolor{azulTabla}\textbf{Describa las salidas que el sistema produce como resultado de la acción} & \multicolumn{3}{>{\RaggedRight\arraybackslash}p{\dimexpr.74\anchoOperacional+4\tabcolsep+2\arrayrulewidth\relax}|}{Aviso de horario y causa; opciones permitidas; constancia de consentimiento; estado pendiente o rechazo; notificación de cambios posteriores cuando corresponda.} \\
\hline

\cellcolor{azulTabla}\textbf{Describa quién o qué usa la salida y para qué es utilizada} & \multicolumn{3}{>{\RaggedRight\arraybackslash}p{\dimexpr.74\anchoOperacional+4\tabcolsep+2\arrayrulewidth\relax}|}{El inversionista decide si continúa y consulta el estado; el equipo de soporte y auditoría examina eventos y autorizaciones.} \\
\hline

\end{tabular}
\end{table}

\begin{table}[H]
\centering
\caption{Escenario operacional EO-04: Supervisión, monitoreo y auditoría de operaciones.}
\label{tab:escenario-eo-04}
\begin{tabular}{|
>{\RaggedRight\arraybackslash}p{.26\anchoOperacional}|
>{\RaggedRight\arraybackslash}p{.25\anchoOperacional}|
>{\RaggedRight\arraybackslash}p{.25\anchoOperacional}|
>{\RaggedRight\arraybackslash}p{.24\anchoOperacional}|}
\hline
\multicolumn{4}{|>{\columncolor{azulTabla}\RaggedRight\arraybackslash}p{\dimexpr\linewidth-2\tabcolsep-2\arrayrulewidth\relax}|}{\bfseries Título del Escenario Operacional} \\
\hline
\multicolumn{4}{|>{\RaggedRight\arraybackslash}p{\dimexpr\linewidth-2\tabcolsep-2\arrayrulewidth\relax}|}{\textbf{Supervisión, monitoreo y auditoría de operaciones}} \\
\hline
\cellcolor{azulTabla}\textbf{Stakeholder} & \multicolumn{1}{>{\RaggedRight\arraybackslash}p{.25\anchoOperacional}|}{Administrador autorizado} & \cellcolor{azulTabla}\textbf{ID} & EO-04 \\
\hline

\cellcolor{azulTabla}\textbf{Descripción general de la funcionalidad} & \multicolumn{3}{>{\RaggedRight\arraybackslash}p{\dimexpr.74\anchoOperacional+4\tabcolsep+2\arrayrulewidth\relax}|}{Permitir el seguimiento de eventos operativos, accesos, cambios de estado de órdenes e incidentes para facilitar control y trazabilidad.} \\
\hline

\cellcolor{azulTabla}\textbf{Describa lo que el stakeholder hace ahora o le gustaría poder hacer} & \multicolumn{3}{>{\RaggedRight\arraybackslash}p{\dimexpr.74\anchoOperacional+4\tabcolsep+2\arrayrulewidth\relax}|}{El administrador revisa actividad relevante, busca una orden o evento, identifica anomalías y consulta evidencias para una revisión autorizada.} \\
\hline

\cellcolor{azulTabla}\textbf{Describa cualquier entrada provista o disponible al momento del inicio} & \multicolumn{3}{>{\RaggedRight\arraybackslash}p{\dimexpr.74\anchoOperacional+4\tabcolsep+2\arrayrulewidth\relax}|}{Credenciales y permisos administrativos; filtros por fechas, usuarios o identificadores; registros de acceso, órdenes, integraciones y errores disponibles.} \\
\hline

\cellcolor{azulTabla}\textbf{Describa el contexto de la operación} & \multicolumn{3}{>{\RaggedRight\arraybackslash}p{\dimexpr.74\anchoOperacional+4\tabcolsep+2\arrayrulewidth\relax}|}{Entorno administrativo protegido; operación normal o investigación de incidente; los datos sensibles se presentan según privilegios y reglas de protección.} \\
\hline

\cellcolor{azulTabla}\textbf{Describa cómo el sistema debe responder} & \multicolumn{3}{>{\RaggedRight\arraybackslash}p{\dimexpr.74\anchoOperacional+4\tabcolsep+2\arrayrulewidth\relax}|}{Validar privilegios; consultar registros íntegros y correlacionados; mostrar secuencia temporal y origen de eventos; permitir filtros y exportación autorizada sin alterar la evidencia original; registrar la propia consulta administrativa.} \\
\hline

\cellcolor{azulTabla}\textbf{Describa las salidas que el sistema produce como resultado de la acción} & \multicolumn{3}{>{\RaggedRight\arraybackslash}p{\dimexpr.74\anchoOperacional+4\tabcolsep+2\arrayrulewidth\relax}|}{Listado de eventos; trazabilidad cronológica de órdenes; indicadores o alertas disponibles; reporte o evidencia de auditoría con identificadores y marcas de tiempo.} \\
\hline

\cellcolor{azulTabla}\textbf{Describa quién o qué usa la salida y para qué es utilizada} & \multicolumn{3}{>{\RaggedRight\arraybackslash}p{\dimexpr.74\anchoOperacional+4\tabcolsep+2\arrayrulewidth\relax}|}{Administradores y responsables autorizados investigan incidentes, verifican operaciones y documentan hallazgos.} \\
\hline

\end{tabular}
\end{table}

\begin{table}[H]
\centering
\caption{Escenario operacional EO-05: Contratación y activación de suscripción premium.}
\label{tab:escenario-eo-05}
\begin{tabular}{|
>{\RaggedRight\arraybackslash}p{.26\anchoOperacional}|
>{\RaggedRight\arraybackslash}p{.25\anchoOperacional}|
>{\RaggedRight\arraybackslash}p{.25\anchoOperacional}|
>{\RaggedRight\arraybackslash}p{.24\anchoOperacional}|}
\hline
\multicolumn{4}{|>{\columncolor{azulTabla}\RaggedRight\arraybackslash}p{\dimexpr\linewidth-2\tabcolsep-2\arrayrulewidth\relax}|}{\bfseries Título del Escenario Operacional} \\
\hline
\multicolumn{4}{|>{\RaggedRight\arraybackslash}p{\dimexpr\linewidth-2\tabcolsep-2\arrayrulewidth\relax}|}{\textbf{Contratación y activación de suscripción premium}} \\
\hline
\cellcolor{azulTabla}\textbf{Stakeholder} & \multicolumn{1}{>{\RaggedRight\arraybackslash}p{.25\anchoOperacional}|}{Inversionista} & \cellcolor{azulTabla}\textbf{ID} & EO-05 \\
\hline

\cellcolor{azulTabla}\textbf{Descripción general de la funcionalidad} & \multicolumn{3}{>{\RaggedRight\arraybackslash}p{\dimexpr.74\anchoOperacional+4\tabcolsep+2\arrayrulewidth\relax}|}{Gestionar la solicitud de un plan premium y la actualización de beneficios una vez verificado el resultado del cobro.} \\
\hline

\cellcolor{azulTabla}\textbf{Describa lo que el stakeholder hace ahora o le gustaría poder hacer} & \multicolumn{3}{>{\RaggedRight\arraybackslash}p{\dimexpr.74\anchoOperacional+4\tabcolsep+2\arrayrulewidth\relax}|}{El inversionista compara planes, selecciona una modalidad, inicia el proceso de pago y consulta la vigencia y beneficios de su suscripción.} \\
\hline

\cellcolor{azulTabla}\textbf{Describa cualquier entrada provista o disponible al momento del inicio} & \multicolumn{3}{>{\RaggedRight\arraybackslash}p{\dimexpr.74\anchoOperacional+4\tabcolsep+2\arrayrulewidth\relax}|}{Sesión autenticada; plan elegido; tarifa vigente y condiciones comerciales; identificador de la transacción; respuesta verificable del proveedor de pagos.} \\
\hline

\cellcolor{azulTabla}\textbf{Describa el contexto de la operación} & \multicolumn{3}{>{\RaggedRight\arraybackslash}p{\dimexpr.74\anchoOperacional+4\tabcolsep+2\arrayrulewidth\relax}|}{Interfaz web y servicio de pagos externo. Stripe se contempla para integración en entorno de pruebas; tarifas y puesta en producción requieren confirmación.} \\
\hline

\cellcolor{azulTabla}\textbf{Describa cómo el sistema debe responder} & \multicolumn{3}{>{\RaggedRight\arraybackslash}p{\dimexpr.74\anchoOperacional+4\tabcolsep+2\arrayrulewidth\relax}|}{Presentar precio y condiciones; iniciar el flujo de pago autorizado; verificar la confirmación recibida; evitar activaciones duplicadas; actualizar estado y vigencia del plan solo tras confirmación válida; gestionar rechazos y errores.} \\
\hline

\cellcolor{azulTabla}\textbf{Describa las salidas que el sistema produce como resultado de la acción} & \multicolumn{3}{>{\RaggedRight\arraybackslash}p{\dimexpr.74\anchoOperacional+4\tabcolsep+2\arrayrulewidth\relax}|}{Resumen del plan; estado del pago; comprobante o referencia; fecha de vigencia; beneficios habilitados o motivo de no activación.} \\
\hline

\cellcolor{azulTabla}\textbf{Describa quién o qué usa la salida y para qué es utilizada} & \multicolumn{3}{>{\RaggedRight\arraybackslash}p{\dimexpr.74\anchoOperacional+4\tabcolsep+2\arrayrulewidth\relax}|}{El inversionista consulta sus beneficios y facturación; administración verifica suscripciones y atiende incidencias comerciales.} \\
\hline

\end{tabular}
\end{table}

\endgroup

\subsection{Casos de Uso: Descripción y Modelo}

\begingroup
\footnotesize
\setlength{\tabcolsep}{5pt}
\renewcommand{\arraystretch}{1.22}
\setlength{\LTleft}{0pt}
\setlength{\LTright}{0pt}
\newlength{\anchoCUAcciones}
\setlength{\anchoCUAcciones}{\dimexpr\textwidth-8\tabcolsep-5\arrayrulewidth\relax}

% CU01
\begin{longtable}{|>{\RaggedRight\arraybackslash}p{.15\anchoCUAcciones}|>{\RaggedRight\arraybackslash}p{.43\anchoCUAcciones}|>{\RaggedRight\arraybackslash}p{.26\anchoCUAcciones}|>{\RaggedRight\arraybackslash}p{.16\anchoCUAcciones}|}
\caption{Caso de uso CU01: Gestionar registro, autenticación y control de acceso.}\label{tab:cu-cu01}\\
\hline
\cellcolor{negroTabla}\color{white}\bfseries Titulo del caso de uso & \multicolumn{1}{>{\RaggedRight\arraybackslash}p{.43\anchoCUAcciones}|}{Gestionar registro, autenticación y control de acceso} & \cellcolor{negroTabla}\color{white}\bfseries ID & CU01\\\hline
\multicolumn{4}{|>{\columncolor{negroTabla}\color{white}\bfseries\RaggedRight\arraybackslash}p{\dimexpr\textwidth-2\tabcolsep-2\arrayrulewidth\relax}|}{Descripcion general del caso de uso}\\\hline
\multicolumn{4}{|>{\RaggedRight\arraybackslash}p{\dimexpr\textwidth-2\tabcolsep-2\arrayrulewidth\relax}|}{Registrar usuarios, autenticar credenciales y establecer sesiones según rol; incluye CU01.1. RF-01 a RF-03; MFA obligatorio RF-04.}\\\hline
\multicolumn{4}{|>{\columncolor{negroTabla}\color{white}\bfseries\RaggedRight\arraybackslash}p{\dimexpr\textwidth-2\tabcolsep-2\arrayrulewidth\relax}|}{Entidades involucradas}\\\hline
\multicolumn{4}{|>{\RaggedRight\arraybackslash}p{\dimexpr\textwidth-2\tabcolsep-2\arrayrulewidth\relax}|}{Inversionista, Comisionista, Administrador; servicio de autenticación}\\\hline
\multicolumn{4}{|>{\columncolor{negroTabla}\color{white}\bfseries\RaggedRight\arraybackslash}p{\dimexpr\textwidth-2\tabcolsep-2\arrayrulewidth\relax}|}{Precondiciones}\\\hline
\multicolumn{4}{|>{\RaggedRight\arraybackslash}p{\dimexpr\textwidth-2\tabcolsep-2\arrayrulewidth\relax}|}{Cuenta válida para inicio de sesión; conexión segura. Para registro, datos exigidos y aceptación de condiciones.}\\\hline
\multicolumn{4}{|>{\columncolor{negroTabla}\color{white}\bfseries\RaggedRight\arraybackslash}p{\dimexpr\textwidth-2\tabcolsep-2\arrayrulewidth\relax}|}{Flujo normal de eventos}\\\hline
\cellcolor{gray!18}\bfseries Paso & \cellcolor{gray!18}\bfseries Usuario / actor & \multicolumn{2}{>{\columncolor{gray!18}\RaggedRight\arraybackslash}p{\dimexpr.42\anchoCUAcciones+2\tabcolsep+\arrayrulewidth\relax}|}{\bfseries Sistema}\\\hline
1 & Solicita registro o inicio de sesión. & \multicolumn{2}{>{\RaggedRight\arraybackslash}p{\dimexpr.42\anchoCUAcciones+2\tabcolsep+\arrayrulewidth\relax}|}{Valida datos y credenciales; inicia CU01.1 y determina permisos RBAC.}\\\hline
2 & Introduce el código MFA requerido. & \multicolumn{2}{>{\RaggedRight\arraybackslash}p{\dimexpr.42\anchoCUAcciones+2\tabcolsep+\arrayrulewidth\relax}|}{Verifica segundo factor y crea sesión segura si es válido.}\\\hline
3 & Consulta su perfil o solicita cerrar sesión. & \multicolumn{2}{>{\RaggedRight\arraybackslash}p{\dimexpr.42\anchoCUAcciones+2\tabcolsep+\arrayrulewidth\relax}|}{Presenta opciones autorizadas y revoca sesión al cerrar.}\\\hline
\multicolumn{4}{|>{\columncolor{negroTabla}\color{white}\bfseries\RaggedRight\arraybackslash}p{\dimexpr\textwidth-2\tabcolsep-2\arrayrulewidth\relax}|}{Postcondiciones principales del caso de uso}\\\hline
\multicolumn{4}{|>{\RaggedRight\arraybackslash}p{\dimexpr\textwidth-2\tabcolsep-2\arrayrulewidth\relax}|}{Cuenta creada o sesión autenticada con privilegios correspondientes; evento auditado.}\\\hline
\multicolumn{4}{|>{\columncolor{negroTabla}\color{white}\bfseries\RaggedRight\arraybackslash}p{\dimexpr\textwidth-2\tabcolsep-2\arrayrulewidth\relax}|}{Caminos de excepcion}\\\hline
\multicolumn{4}{|>{\RaggedRight\arraybackslash}p{\dimexpr\textwidth-2\tabcolsep-2\arrayrulewidth\relax}|}{Correo duplicado, credenciales inválidas, MFA fallido o expirado, bloqueo por intentos y servicio indisponible: rechazar acceso sin revelar datos sensibles.}\\\hline
\end{longtable}

% CU01.1
\begin{longtable}{|>{\RaggedRight\arraybackslash}p{.15\anchoCUAcciones}|>{\RaggedRight\arraybackslash}p{.43\anchoCUAcciones}|>{\RaggedRight\arraybackslash}p{.26\anchoCUAcciones}|>{\RaggedRight\arraybackslash}p{.16\anchoCUAcciones}|}
\caption{Caso de uso CU01.1: Validar segundo factor de autenticación (MFA).}\label{tab:cu-cu01-1}\\
\hline
\cellcolor{negroTabla}\color{white}\bfseries Titulo del caso de uso & \multicolumn{1}{>{\RaggedRight\arraybackslash}p{.43\anchoCUAcciones}|}{Validar segundo factor de autenticación (MFA)} & \cellcolor{negroTabla}\color{white}\bfseries ID & CU01.1\\\hline
\multicolumn{4}{|>{\columncolor{negroTabla}\color{white}\bfseries\RaggedRight\arraybackslash}p{\dimexpr\textwidth-2\tabcolsep-2\arrayrulewidth\relax}|}{Descripcion general del caso de uso}\\\hline
\multicolumn{4}{|>{\RaggedRight\arraybackslash}p{\dimexpr\textwidth-2\tabcolsep-2\arrayrulewidth\relax}|}{Verificar segundo factor durante inicio de sesión y acciones críticas; incluido por CU01. RF-04; mecanismo TOTP/SMS sujeto a configuración.}\\\hline
\multicolumn{4}{|>{\columncolor{negroTabla}\color{white}\bfseries\RaggedRight\arraybackslash}p{\dimexpr\textwidth-2\tabcolsep-2\arrayrulewidth\relax}|}{Entidades involucradas}\\\hline
\multicolumn{4}{|>{\RaggedRight\arraybackslash}p{\dimexpr\textwidth-2\tabcolsep-2\arrayrulewidth\relax}|}{Usuario autenticándose; servicio MFA}\\\hline
\multicolumn{4}{|>{\columncolor{negroTabla}\color{white}\bfseries\RaggedRight\arraybackslash}p{\dimexpr\textwidth-2\tabcolsep-2\arrayrulewidth\relax}|}{Precondiciones}\\\hline
\multicolumn{4}{|>{\RaggedRight\arraybackslash}p{\dimexpr\textwidth-2\tabcolsep-2\arrayrulewidth\relax}|}{Primer factor validado; mecanismo MFA configurado y desafío vigente.}\\\hline
\multicolumn{4}{|>{\columncolor{negroTabla}\color{white}\bfseries\RaggedRight\arraybackslash}p{\dimexpr\textwidth-2\tabcolsep-2\arrayrulewidth\relax}|}{Flujo normal de eventos}\\\hline
\cellcolor{gray!18}\bfseries Paso & \cellcolor{gray!18}\bfseries Usuario / actor & \multicolumn{2}{>{\columncolor{gray!18}\RaggedRight\arraybackslash}p{\dimexpr.42\anchoCUAcciones+2\tabcolsep+\arrayrulewidth\relax}|}{\bfseries Sistema}\\\hline
1 & Recibe solicitud de código y proporciona TOTP o código habilitado. & \multicolumn{2}{>{\RaggedRight\arraybackslash}p{\dimexpr.42\anchoCUAcciones+2\tabcolsep+\arrayrulewidth\relax}|}{Comprueba validez, vigencia y asociación del código con la cuenta.}\\\hline
2 & Solicita reenvío o nuevo desafío cuando aplique. & \multicolumn{2}{>{\RaggedRight\arraybackslash}p{\dimexpr.42\anchoCUAcciones+2\tabcolsep+\arrayrulewidth\relax}|}{Gestiona el nuevo desafío conforme a controles de seguridad.}\\\hline
\multicolumn{4}{|>{\columncolor{negroTabla}\color{white}\bfseries\RaggedRight\arraybackslash}p{\dimexpr\textwidth-2\tabcolsep-2\arrayrulewidth\relax}|}{Postcondiciones principales del caso de uso}\\\hline
\multicolumn{4}{|>{\RaggedRight\arraybackslash}p{\dimexpr\textwidth-2\tabcolsep-2\arrayrulewidth\relax}|}{Acceso o acción crítica autorizada únicamente con MFA válido; intento registrado.}\\\hline
\multicolumn{4}{|>{\columncolor{negroTabla}\color{white}\bfseries\RaggedRight\arraybackslash}p{\dimexpr\textwidth-2\tabcolsep-2\arrayrulewidth\relax}|}{Caminos de excepcion}\\\hline
\multicolumn{4}{|>{\RaggedRight\arraybackslash}p{\dimexpr\textwidth-2\tabcolsep-2\arrayrulewidth\relax}|}{Código erróneo, vencido, reutilizado, límite de intentos o proveedor SMS no disponible: denegar y permitir recuperación controlada.}\\\hline
\end{longtable}

% CU02
\begin{longtable}{|>{\RaggedRight\arraybackslash}p{.15\anchoCUAcciones}|>{\RaggedRight\arraybackslash}p{.43\anchoCUAcciones}|>{\RaggedRight\arraybackslash}p{.26\anchoCUAcciones}|>{\RaggedRight\arraybackslash}p{.16\anchoCUAcciones}|}
\caption{Caso de uso CU02: Administrar perfil de inversionista y preferencias.}\label{tab:cu-cu02}\\
\hline
\cellcolor{negroTabla}\color{white}\bfseries Titulo del caso de uso & \multicolumn{1}{>{\RaggedRight\arraybackslash}p{.43\anchoCUAcciones}|}{Administrar perfil de inversionista y preferencias} & \cellcolor{negroTabla}\color{white}\bfseries ID & CU02\\\hline
\multicolumn{4}{|>{\columncolor{negroTabla}\color{white}\bfseries\RaggedRight\arraybackslash}p{\dimexpr\textwidth-2\tabcolsep-2\arrayrulewidth\relax}|}{Descripcion general del caso de uso}\\\hline
\multicolumn{4}{|>{\RaggedRight\arraybackslash}p{\dimexpr\textwidth-2\tabcolsep-2\arrayrulewidth\relax}|}{Consultar y modificar datos personales, perfil de riesgo y preferencias de notificación. RF-05 y RF-06; tratamiento de datos personales.}\\\hline
\multicolumn{4}{|>{\columncolor{negroTabla}\color{white}\bfseries\RaggedRight\arraybackslash}p{\dimexpr\textwidth-2\tabcolsep-2\arrayrulewidth\relax}|}{Entidades involucradas}\\\hline
\multicolumn{4}{|>{\RaggedRight\arraybackslash}p{\dimexpr\textwidth-2\tabcolsep-2\arrayrulewidth\relax}|}{Inversionista}\\\hline
\multicolumn{4}{|>{\columncolor{negroTabla}\color{white}\bfseries\RaggedRight\arraybackslash}p{\dimexpr\textwidth-2\tabcolsep-2\arrayrulewidth\relax}|}{Precondiciones}\\\hline
\multicolumn{4}{|>{\RaggedRight\arraybackslash}p{\dimexpr\textwidth-2\tabcolsep-2\arrayrulewidth\relax}|}{Sesión activa y permisos sobre el perfil propio.}\\\hline
\multicolumn{4}{|>{\columncolor{negroTabla}\color{white}\bfseries\RaggedRight\arraybackslash}p{\dimexpr\textwidth-2\tabcolsep-2\arrayrulewidth\relax}|}{Flujo normal de eventos}\\\hline
\cellcolor{gray!18}\bfseries Paso & \cellcolor{gray!18}\bfseries Usuario / actor & \multicolumn{2}{>{\columncolor{gray!18}\RaggedRight\arraybackslash}p{\dimexpr.42\anchoCUAcciones+2\tabcolsep+\arrayrulewidth\relax}|}{\bfseries Sistema}\\\hline
1 & Abre su perfil y edita los campos permitidos. & \multicolumn{2}{>{\RaggedRight\arraybackslash}p{\dimexpr.42\anchoCUAcciones+2\tabcolsep+\arrayrulewidth\relax}|}{Muestra datos vigentes y valida formatos y reglas.}\\\hline
2 & Confirma los cambios. & \multicolumn{2}{>{\RaggedRight\arraybackslash}p{\dimexpr.42\anchoCUAcciones+2\tabcolsep+\arrayrulewidth\relax}|}{Guarda modificaciones y registra cambios sensibles.}\\\hline
\multicolumn{4}{|>{\columncolor{negroTabla}\color{white}\bfseries\RaggedRight\arraybackslash}p{\dimexpr\textwidth-2\tabcolsep-2\arrayrulewidth\relax}|}{Postcondiciones principales del caso de uso}\\\hline
\multicolumn{4}{|>{\RaggedRight\arraybackslash}p{\dimexpr\textwidth-2\tabcolsep-2\arrayrulewidth\relax}|}{Perfil y preferencias actualizados y visibles para su titular.}\\\hline
\multicolumn{4}{|>{\columncolor{negroTabla}\color{white}\bfseries\RaggedRight\arraybackslash}p{\dimexpr\textwidth-2\tabcolsep-2\arrayrulewidth\relax}|}{Caminos de excepcion}\\\hline
\multicolumn{4}{|>{\RaggedRight\arraybackslash}p{\dimexpr\textwidth-2\tabcolsep-2\arrayrulewidth\relax}|}{Datos inválidos, conflicto de actualización o sesión vencida: no aplicar cambios parciales y mostrar causa.}\\\hline
\end{longtable}

% CU03
\begin{longtable}{|>{\RaggedRight\arraybackslash}p{.15\anchoCUAcciones}|>{\RaggedRight\arraybackslash}p{.43\anchoCUAcciones}|>{\RaggedRight\arraybackslash}p{.26\anchoCUAcciones}|>{\RaggedRight\arraybackslash}p{.16\anchoCUAcciones}|}
\caption{Caso de uso CU03: Gestionar vinculación y calificaciones de comisionistas.}\label{tab:cu-cu03}\\
\hline
\cellcolor{negroTabla}\color{white}\bfseries Titulo del caso de uso & \multicolumn{1}{>{\RaggedRight\arraybackslash}p{.43\anchoCUAcciones}|}{Gestionar vinculación y calificaciones de comisionistas} & \cellcolor{negroTabla}\color{white}\bfseries ID & CU03\\\hline
\multicolumn{4}{|>{\columncolor{negroTabla}\color{white}\bfseries\RaggedRight\arraybackslash}p{\dimexpr\textwidth-2\tabcolsep-2\arrayrulewidth\relax}|}{Descripcion general del caso de uso}\\\hline
\multicolumn{4}{|>{\RaggedRight\arraybackslash}p{\dimexpr\textwidth-2\tabcolsep-2\arrayrulewidth\relax}|}{Buscar comisionistas, solicitar vinculación y registrar calificaciones conforme al vínculo y reglas de negocio. RF-11 a RF-15; verificar reglas contractuales.}\\\hline
\multicolumn{4}{|>{\columncolor{negroTabla}\color{white}\bfseries\RaggedRight\arraybackslash}p{\dimexpr\textwidth-2\tabcolsep-2\arrayrulewidth\relax}|}{Entidades involucradas}\\\hline
\multicolumn{4}{|>{\RaggedRight\arraybackslash}p{\dimexpr\textwidth-2\tabcolsep-2\arrayrulewidth\relax}|}{Inversionista, Comisionista}\\\hline
\multicolumn{4}{|>{\columncolor{negroTabla}\color{white}\bfseries\RaggedRight\arraybackslash}p{\dimexpr\textwidth-2\tabcolsep-2\arrayrulewidth\relax}|}{Precondiciones}\\\hline
\multicolumn{4}{|>{\RaggedRight\arraybackslash}p{\dimexpr\textwidth-2\tabcolsep-2\arrayrulewidth\relax}|}{Inversionista autenticado; comisionista habilitado y disponible.}\\\hline
\multicolumn{4}{|>{\columncolor{negroTabla}\color{white}\bfseries\RaggedRight\arraybackslash}p{\dimexpr\textwidth-2\tabcolsep-2\arrayrulewidth\relax}|}{Flujo normal de eventos}\\\hline
\cellcolor{gray!18}\bfseries Paso & \cellcolor{gray!18}\bfseries Usuario / actor & \multicolumn{2}{>{\columncolor{gray!18}\RaggedRight\arraybackslash}p{\dimexpr.42\anchoCUAcciones+2\tabcolsep+\arrayrulewidth\relax}|}{\bfseries Sistema}\\\hline
1 & Consulta perfiles y solicita vinculación. & \multicolumn{2}{>{\RaggedRight\arraybackslash}p{\dimexpr.42\anchoCUAcciones+2\tabcolsep+\arrayrulewidth\relax}|}{Presenta información autorizada y registra solicitud pendiente.}\\\hline
2 & Comisionista acepta o rechaza la solicitud; inversionista califica cuando corresponda. & \multicolumn{2}{>{\RaggedRight\arraybackslash}p{\dimexpr.42\anchoCUAcciones+2\tabcolsep+\arrayrulewidth\relax}|}{Actualiza vínculo, habilita acceso autorizado y registra valoración válida.}\\\hline
\multicolumn{4}{|>{\columncolor{negroTabla}\color{white}\bfseries\RaggedRight\arraybackslash}p{\dimexpr\textwidth-2\tabcolsep-2\arrayrulewidth\relax}|}{Postcondiciones principales del caso de uso}\\\hline
\multicolumn{4}{|>{\RaggedRight\arraybackslash}p{\dimexpr\textwidth-2\tabcolsep-2\arrayrulewidth\relax}|}{Relación registrada con estado verificable; calificación asociada cuando procede.}\\\hline
\multicolumn{4}{|>{\columncolor{negroTabla}\color{white}\bfseries\RaggedRight\arraybackslash}p{\dimexpr\textwidth-2\tabcolsep-2\arrayrulewidth\relax}|}{Caminos de excepcion}\\\hline
\multicolumn{4}{|>{\RaggedRight\arraybackslash}p{\dimexpr\textwidth-2\tabcolsep-2\arrayrulewidth\relax}|}{Comisionista no habilitado, solicitud duplicada, vínculo rechazado o valoración no permitida: impedir operación y notificar.}\\\hline
\end{longtable}

% CU04
\begin{longtable}{|>{\RaggedRight\arraybackslash}p{.15\anchoCUAcciones}|>{\RaggedRight\arraybackslash}p{.43\anchoCUAcciones}|>{\RaggedRight\arraybackslash}p{.26\anchoCUAcciones}|>{\RaggedRight\arraybackslash}p{.16\anchoCUAcciones}|}
\caption{Caso de uso CU04: Formular y tramitar propuestas de inversión.}\label{tab:cu-cu04}\\
\hline
\cellcolor{negroTabla}\color{white}\bfseries Titulo del caso de uso & \multicolumn{1}{>{\RaggedRight\arraybackslash}p{.43\anchoCUAcciones}|}{Formular y tramitar propuestas de inversión} & \cellcolor{negroTabla}\color{white}\bfseries ID & CU04\\\hline
\multicolumn{4}{|>{\columncolor{negroTabla}\color{white}\bfseries\RaggedRight\arraybackslash}p{\dimexpr\textwidth-2\tabcolsep-2\arrayrulewidth\relax}|}{Descripcion general del caso de uso}\\\hline
\multicolumn{4}{|>{\RaggedRight\arraybackslash}p{\dimexpr\textwidth-2\tabcolsep-2\arrayrulewidth\relax}|}{Crear propuestas de cartera y permitir su revisión, aprobación o rechazo por el inversionista; puede extenderse mediante CU04.1. RF-16 a RF-19; aprobación no equivale a ejecución.}\\\hline
\multicolumn{4}{|>{\columncolor{negroTabla}\color{white}\bfseries\RaggedRight\arraybackslash}p{\dimexpr\textwidth-2\tabcolsep-2\arrayrulewidth\relax}|}{Entidades involucradas}\\\hline
\multicolumn{4}{|>{\RaggedRight\arraybackslash}p{\dimexpr\textwidth-2\tabcolsep-2\arrayrulewidth\relax}|}{Comisionista, Inversionista}\\\hline
\multicolumn{4}{|>{\columncolor{negroTabla}\color{white}\bfseries\RaggedRight\arraybackslash}p{\dimexpr\textwidth-2\tabcolsep-2\arrayrulewidth\relax}|}{Precondiciones}\\\hline
\multicolumn{4}{|>{\RaggedRight\arraybackslash}p{\dimexpr\textwidth-2\tabcolsep-2\arrayrulewidth\relax}|}{Vínculo activo y autorizado entre ambas partes; instrumentos disponibles para consulta.}\\\hline
\multicolumn{4}{|>{\columncolor{negroTabla}\color{white}\bfseries\RaggedRight\arraybackslash}p{\dimexpr\textwidth-2\tabcolsep-2\arrayrulewidth\relax}|}{Flujo normal de eventos}\\\hline
\cellcolor{gray!18}\bfseries Paso & \cellcolor{gray!18}\bfseries Usuario / actor & \multicolumn{2}{>{\columncolor{gray!18}\RaggedRight\arraybackslash}p{\dimexpr.42\anchoCUAcciones+2\tabcolsep+\arrayrulewidth\relax}|}{\bfseries Sistema}\\\hline
1 & Comisionista define activos, cantidades y justificación de la propuesta. & \multicolumn{2}{>{\RaggedRight\arraybackslash}p{\dimexpr.42\anchoCUAcciones+2\tabcolsep+\arrayrulewidth\relax}|}{Valida datos, registra versión y la presenta al inversionista.}\\\hline
2 & Inversionista revisa y aprueba o rechaza. & \multicolumn{2}{>{\RaggedRight\arraybackslash}p{\dimexpr.42\anchoCUAcciones+2\tabcolsep+\arrayrulewidth\relax}|}{Actualiza estado; si fue aprobada, habilita proceso de firma lógica CU04.1.}\\\hline
\multicolumn{4}{|>{\columncolor{negroTabla}\color{white}\bfseries\RaggedRight\arraybackslash}p{\dimexpr\textwidth-2\tabcolsep-2\arrayrulewidth\relax}|}{Postcondiciones principales del caso de uso}\\\hline
\multicolumn{4}{|>{\RaggedRight\arraybackslash}p{\dimexpr\textwidth-2\tabcolsep-2\arrayrulewidth\relax}|}{Propuesta versionada con decisión del inversionista y trazabilidad.}\\\hline
\multicolumn{4}{|>{\columncolor{negroTabla}\color{white}\bfseries\RaggedRight\arraybackslash}p{\dimexpr\textwidth-2\tabcolsep-2\arrayrulewidth\relax}|}{Caminos de excepcion}\\\hline
\multicolumn{4}{|>{\RaggedRight\arraybackslash}p{\dimexpr\textwidth-2\tabcolsep-2\arrayrulewidth\relax}|}{Vínculo revocado, propuesta vencida, datos incompletos o cambio de condiciones: suspender y solicitar revisión.}\\\hline
\end{longtable}

% CU04.1
\begin{longtable}{|>{\RaggedRight\arraybackslash}p{.15\anchoCUAcciones}|>{\RaggedRight\arraybackslash}p{.43\anchoCUAcciones}|>{\RaggedRight\arraybackslash}p{.26\anchoCUAcciones}|>{\RaggedRight\arraybackslash}p{.16\anchoCUAcciones}|}
\caption{Caso de uso CU04.1: Firmar lógicamente y autorizar propuesta.}\label{tab:cu-cu04-1}\\
\hline
\cellcolor{negroTabla}\color{white}\bfseries Titulo del caso de uso & \multicolumn{1}{>{\RaggedRight\arraybackslash}p{.43\anchoCUAcciones}|}{Firmar lógicamente y autorizar propuesta} & \cellcolor{negroTabla}\color{white}\bfseries ID & CU04.1\\\hline
\multicolumn{4}{|>{\columncolor{negroTabla}\color{white}\bfseries\RaggedRight\arraybackslash}p{\dimexpr\textwidth-2\tabcolsep-2\arrayrulewidth\relax}|}{Descripcion general del caso de uso}\\\hline
\multicolumn{4}{|>{\RaggedRight\arraybackslash}p{\dimexpr\textwidth-2\tabcolsep-2\arrayrulewidth\relax}|}{Extensión condicional de CU04 posterior a la aceptación del inversionista. RF-20; la firma lógica no implica necesariamente firma digital certificada.}\\\hline
\multicolumn{4}{|>{\columncolor{negroTabla}\color{white}\bfseries\RaggedRight\arraybackslash}p{\dimexpr\textwidth-2\tabcolsep-2\arrayrulewidth\relax}|}{Entidades involucradas}\\\hline
\multicolumn{4}{|>{\RaggedRight\arraybackslash}p{\dimexpr\textwidth-2\tabcolsep-2\arrayrulewidth\relax}|}{Comisionista}\\\hline
\multicolumn{4}{|>{\columncolor{negroTabla}\color{white}\bfseries\RaggedRight\arraybackslash}p{\dimexpr\textwidth-2\tabcolsep-2\arrayrulewidth\relax}|}{Precondiciones}\\\hline
\multicolumn{4}{|>{\RaggedRight\arraybackslash}p{\dimexpr\textwidth-2\tabcolsep-2\arrayrulewidth\relax}|}{Propuesta aprobada, vigente y asociada a comisionista autorizado; autenticación reforzada cuando aplique.}\\\hline
\multicolumn{4}{|>{\columncolor{negroTabla}\color{white}\bfseries\RaggedRight\arraybackslash}p{\dimexpr\textwidth-2\tabcolsep-2\arrayrulewidth\relax}|}{Flujo normal de eventos}\\\hline
\cellcolor{gray!18}\bfseries Paso & \cellcolor{gray!18}\bfseries Usuario / actor & \multicolumn{2}{>{\columncolor{gray!18}\RaggedRight\arraybackslash}p{\dimexpr.42\anchoCUAcciones+2\tabcolsep+\arrayrulewidth\relax}|}{\bfseries Sistema}\\\hline
1 & Comisionista revisa la propuesta aceptada y solicita firma. & \multicolumn{2}{>{\RaggedRight\arraybackslash}p{\dimexpr.42\anchoCUAcciones+2\tabcolsep+\arrayrulewidth\relax}|}{Presenta versión íntegra y solicita confirmación verificable.}\\\hline
2 & Confirma la firma lógica. & \multicolumn{2}{>{\RaggedRight\arraybackslash}p{\dimexpr.42\anchoCUAcciones+2\tabcolsep+\arrayrulewidth\relax}|}{Registra autor, instante, versión e integridad; habilita tramitación posterior.}\\\hline
\multicolumn{4}{|>{\columncolor{negroTabla}\color{white}\bfseries\RaggedRight\arraybackslash}p{\dimexpr\textwidth-2\tabcolsep-2\arrayrulewidth\relax}|}{Postcondiciones principales del caso de uso}\\\hline
\multicolumn{4}{|>{\RaggedRight\arraybackslash}p{\dimexpr\textwidth-2\tabcolsep-2\arrayrulewidth\relax}|}{Propuesta firmada y autorizada para siguiente etapa, sin afirmar ejecución bursátil.}\\\hline
\multicolumn{4}{|>{\columncolor{negroTabla}\color{white}\bfseries\RaggedRight\arraybackslash}p{\dimexpr\textwidth-2\tabcolsep-2\arrayrulewidth\relax}|}{Caminos de excepcion}\\\hline
\multicolumn{4}{|>{\RaggedRight\arraybackslash}p{\dimexpr\textwidth-2\tabcolsep-2\arrayrulewidth\relax}|}{Propuesta modificada, aprobación revocada, identidad no validada o firma fallida: no autorizar envío.}\\\hline
\end{longtable}

% CU05
\begin{longtable}{|>{\RaggedRight\arraybackslash}p{.15\anchoCUAcciones}|>{\RaggedRight\arraybackslash}p{.43\anchoCUAcciones}|>{\RaggedRight\arraybackslash}p{.26\anchoCUAcciones}|>{\RaggedRight\arraybackslash}p{.16\anchoCUAcciones}|}
\caption{Caso de uso CU05: Consultar y monitorear instrumentos y cotizaciones.}\label{tab:cu-cu05}\\
\hline
\cellcolor{negroTabla}\color{white}\bfseries Titulo del caso de uso & \multicolumn{1}{>{\RaggedRight\arraybackslash}p{.43\anchoCUAcciones}|}{Consultar y monitorear instrumentos y cotizaciones} & \cellcolor{negroTabla}\color{white}\bfseries ID & CU05\\\hline
\multicolumn{4}{|>{\columncolor{negroTabla}\color{white}\bfseries\RaggedRight\arraybackslash}p{\dimexpr\textwidth-2\tabcolsep-2\arrayrulewidth\relax}|}{Descripcion general del caso de uso}\\\hline
\multicolumn{4}{|>{\RaggedRight\arraybackslash}p{\dimexpr\textwidth-2\tabcolsep-2\arrayrulewidth\relax}|}{Buscar instrumentos y consultar precios e históricos disponibles por integración externa. RF-07, RF-08 y RF-10; tiempo real depende del plan.}\\\hline
\multicolumn{4}{|>{\columncolor{negroTabla}\color{white}\bfseries\RaggedRight\arraybackslash}p{\dimexpr\textwidth-2\tabcolsep-2\arrayrulewidth\relax}|}{Entidades involucradas}\\\hline
\multicolumn{4}{|>{\RaggedRight\arraybackslash}p{\dimexpr\textwidth-2\tabcolsep-2\arrayrulewidth\relax}|}{Inversionista, Comisionista; Alpaca (datos bursátiles)}\\\hline
\multicolumn{4}{|>{\columncolor{negroTabla}\color{white}\bfseries\RaggedRight\arraybackslash}p{\dimexpr\textwidth-2\tabcolsep-2\arrayrulewidth\relax}|}{Precondiciones}\\\hline
\multicolumn{4}{|>{\RaggedRight\arraybackslash}p{\dimexpr\textwidth-2\tabcolsep-2\arrayrulewidth\relax}|}{Servicio web disponible; permisos de datos y acceso a Alpaca vigentes.}\\\hline
\multicolumn{4}{|>{\columncolor{negroTabla}\color{white}\bfseries\RaggedRight\arraybackslash}p{\dimexpr\textwidth-2\tabcolsep-2\arrayrulewidth\relax}|}{Flujo normal de eventos}\\\hline
\cellcolor{gray!18}\bfseries Paso & \cellcolor{gray!18}\bfseries Usuario / actor & \multicolumn{2}{>{\columncolor{gray!18}\RaggedRight\arraybackslash}p{\dimexpr.42\anchoCUAcciones+2\tabcolsep+\arrayrulewidth\relax}|}{\bfseries Sistema}\\\hline
1 & Busca un instrumento o abre su ficha. & \multicolumn{2}{>{\RaggedRight\arraybackslash}p{\dimexpr.42\anchoCUAcciones+2\tabcolsep+\arrayrulewidth\relax}|}{Consulta datos permitidos y muestra precio, marca temporal y procedencia.}\\\hline
2 & Selecciona período o detalle. & \multicolumn{2}{>{\RaggedRight\arraybackslash}p{\dimexpr.42\anchoCUAcciones+2\tabcolsep+\arrayrulewidth\relax}|}{Devuelve serie histórica disponible y comunica posibles retrasos.}\\\hline
\multicolumn{4}{|>{\columncolor{negroTabla}\color{white}\bfseries\RaggedRight\arraybackslash}p{\dimexpr\textwidth-2\tabcolsep-2\arrayrulewidth\relax}|}{Postcondiciones principales del caso de uso}\\\hline
\multicolumn{4}{|>{\RaggedRight\arraybackslash}p{\dimexpr\textwidth-2\tabcolsep-2\arrayrulewidth\relax}|}{Información bursátil presentada con contexto temporal y cobertura del proveedor.}\\\hline
\multicolumn{4}{|>{\columncolor{negroTabla}\color{white}\bfseries\RaggedRight\arraybackslash}p{\dimexpr\textwidth-2\tabcolsep-2\arrayrulewidth\relax}|}{Caminos de excepcion}\\\hline
\multicolumn{4}{|>{\RaggedRight\arraybackslash}p{\dimexpr\textwidth-2\tabcolsep-2\arrayrulewidth\relax}|}{Símbolo inexistente, feed retrasado, límite de API o indisponibilidad: mostrar estado y evitar datos como actuales.}\\\hline
\end{longtable}

% CU06
\begin{longtable}{|>{\RaggedRight\arraybackslash}p{.15\anchoCUAcciones}|>{\RaggedRight\arraybackslash}p{.43\anchoCUAcciones}|>{\RaggedRight\arraybackslash}p{.26\anchoCUAcciones}|>{\RaggedRight\arraybackslash}p{.16\anchoCUAcciones}|}
\caption{Caso de uso CU06: Administrar listas de observación y alertas de precios.}\label{tab:cu-cu06}\\
\hline
\cellcolor{negroTabla}\color{white}\bfseries Titulo del caso de uso & \multicolumn{1}{>{\RaggedRight\arraybackslash}p{.43\anchoCUAcciones}|}{Administrar listas de observación y alertas de precios} & \cellcolor{negroTabla}\color{white}\bfseries ID & CU06\\\hline
\multicolumn{4}{|>{\columncolor{negroTabla}\color{white}\bfseries\RaggedRight\arraybackslash}p{\dimexpr\textwidth-2\tabcolsep-2\arrayrulewidth\relax}|}{Descripcion general del caso de uso}\\\hline
\multicolumn{4}{|>{\RaggedRight\arraybackslash}p{\dimexpr\textwidth-2\tabcolsep-2\arrayrulewidth\relax}|}{Crear listas y establecer umbrales de precios; CU06.1 ocurre si se cumple una condición. RF-09; disparo en CU06.1.}\\\hline
\multicolumn{4}{|>{\columncolor{negroTabla}\color{white}\bfseries\RaggedRight\arraybackslash}p{\dimexpr\textwidth-2\tabcolsep-2\arrayrulewidth\relax}|}{Entidades involucradas}\\\hline
\multicolumn{4}{|>{\RaggedRight\arraybackslash}p{\dimexpr\textwidth-2\tabcolsep-2\arrayrulewidth\relax}|}{Inversionista}\\\hline
\multicolumn{4}{|>{\columncolor{negroTabla}\color{white}\bfseries\RaggedRight\arraybackslash}p{\dimexpr\textwidth-2\tabcolsep-2\arrayrulewidth\relax}|}{Precondiciones}\\\hline
\multicolumn{4}{|>{\RaggedRight\arraybackslash}p{\dimexpr\textwidth-2\tabcolsep-2\arrayrulewidth\relax}|}{Usuario autenticado; instrumento identificable y notificaciones configuradas.}\\\hline
\multicolumn{4}{|>{\columncolor{negroTabla}\color{white}\bfseries\RaggedRight\arraybackslash}p{\dimexpr\textwidth-2\tabcolsep-2\arrayrulewidth\relax}|}{Flujo normal de eventos}\\\hline
\cellcolor{gray!18}\bfseries Paso & \cellcolor{gray!18}\bfseries Usuario / actor & \multicolumn{2}{>{\columncolor{gray!18}\RaggedRight\arraybackslash}p{\dimexpr.42\anchoCUAcciones+2\tabcolsep+\arrayrulewidth\relax}|}{\bfseries Sistema}\\\hline
1 & Agrega instrumentos y define límites superior/inferior. & \multicolumn{2}{>{\RaggedRight\arraybackslash}p{\dimexpr.42\anchoCUAcciones+2\tabcolsep+\arrayrulewidth\relax}|}{Valida parámetros y guarda reglas de seguimiento.}\\\hline
2 & Edita, desactiva o elimina alertas. & \multicolumn{2}{>{\RaggedRight\arraybackslash}p{\dimexpr.42\anchoCUAcciones+2\tabcolsep+\arrayrulewidth\relax}|}{Actualiza sus estados y programa evaluación con datos disponibles.}\\\hline
\multicolumn{4}{|>{\columncolor{negroTabla}\color{white}\bfseries\RaggedRight\arraybackslash}p{\dimexpr\textwidth-2\tabcolsep-2\arrayrulewidth\relax}|}{Postcondiciones principales del caso de uso}\\\hline
\multicolumn{4}{|>{\RaggedRight\arraybackslash}p{\dimexpr\textwidth-2\tabcolsep-2\arrayrulewidth\relax}|}{Listas y reglas persistidas; alertas activas identificadas.}\\\hline
\multicolumn{4}{|>{\columncolor{negroTabla}\color{white}\bfseries\RaggedRight\arraybackslash}p{\dimexpr\textwidth-2\tabcolsep-2\arrayrulewidth\relax}|}{Caminos de excepcion}\\\hline
\multicolumn{4}{|>{\RaggedRight\arraybackslash}p{\dimexpr\textwidth-2\tabcolsep-2\arrayrulewidth\relax}|}{Umbral inválido, instrumento no soportado o cuota excedida: rechazar configuración y explicar.}\\\hline
\end{longtable}

% CU06.1
\begin{longtable}{|>{\RaggedRight\arraybackslash}p{.15\anchoCUAcciones}|>{\RaggedRight\arraybackslash}p{.43\anchoCUAcciones}|>{\RaggedRight\arraybackslash}p{.26\anchoCUAcciones}|>{\RaggedRight\arraybackslash}p{.16\anchoCUAcciones}|}
\caption{Caso de uso CU06.1: Disparar notificación de alerta por umbral.}\label{tab:cu-cu06-1}\\
\hline
\cellcolor{negroTabla}\color{white}\bfseries Titulo del caso de uso & \multicolumn{1}{>{\RaggedRight\arraybackslash}p{.43\anchoCUAcciones}|}{Disparar notificación de alerta por umbral} & \cellcolor{negroTabla}\color{white}\bfseries ID & CU06.1\\\hline
\multicolumn{4}{|>{\columncolor{negroTabla}\color{white}\bfseries\RaggedRight\arraybackslash}p{\dimexpr\textwidth-2\tabcolsep-2\arrayrulewidth\relax}|}{Descripcion general del caso de uso}\\\hline
\multicolumn{4}{|>{\RaggedRight\arraybackslash}p{\dimexpr\textwidth-2\tabcolsep-2\arrayrulewidth\relax}|}{Extensión de CU06 activada cuando una cotización cruza el umbral definido. RF-09.1; actor disparador externo.}\\\hline
\multicolumn{4}{|>{\columncolor{negroTabla}\color{white}\bfseries\RaggedRight\arraybackslash}p{\dimexpr\textwidth-2\tabcolsep-2\arrayrulewidth\relax}|}{Entidades involucradas}\\\hline
\multicolumn{4}{|>{\RaggedRight\arraybackslash}p{\dimexpr\textwidth-2\tabcolsep-2\arrayrulewidth\relax}|}{Alpaca (datos bursátiles), Inversionista; servicio de notificaciones}\\\hline
\multicolumn{4}{|>{\columncolor{negroTabla}\color{white}\bfseries\RaggedRight\arraybackslash}p{\dimexpr\textwidth-2\tabcolsep-2\arrayrulewidth\relax}|}{Precondiciones}\\\hline
\multicolumn{4}{|>{\RaggedRight\arraybackslash}p{\dimexpr\textwidth-2\tabcolsep-2\arrayrulewidth\relax}|}{Regla activa, cotización válida y canal de notificación configurado.}\\\hline
\multicolumn{4}{|>{\columncolor{negroTabla}\color{white}\bfseries\RaggedRight\arraybackslash}p{\dimexpr\textwidth-2\tabcolsep-2\arrayrulewidth\relax}|}{Flujo normal de eventos}\\\hline
\cellcolor{gray!18}\bfseries Paso & \cellcolor{gray!18}\bfseries Usuario / actor & \multicolumn{2}{>{\columncolor{gray!18}\RaggedRight\arraybackslash}p{\dimexpr.42\anchoCUAcciones+2\tabcolsep+\arrayrulewidth\relax}|}{\bfseries Sistema}\\\hline
1 & No requiere acción directa; mantiene una alerta configurada. & \multicolumn{2}{>{\RaggedRight\arraybackslash}p{\dimexpr.42\anchoCUAcciones+2\tabcolsep+\arrayrulewidth\relax}|}{Recibe cotización y compara contra umbral y estado previo.}\\\hline
2 & Consulta notificación cuando llega. & \multicolumn{2}{>{\RaggedRight\arraybackslash}p{\dimexpr.42\anchoCUAcciones+2\tabcolsep+\arrayrulewidth\relax}|}{Emite alerta, registra evento y aplica política antirrepetición.}\\\hline
\multicolumn{4}{|>{\columncolor{negroTabla}\color{white}\bfseries\RaggedRight\arraybackslash}p{\dimexpr\textwidth-2\tabcolsep-2\arrayrulewidth\relax}|}{Postcondiciones principales del caso de uso}\\\hline
\multicolumn{4}{|>{\RaggedRight\arraybackslash}p{\dimexpr\textwidth-2\tabcolsep-2\arrayrulewidth\relax}|}{Notificación con instrumento, valor y hora; evento trazable.}\\\hline
\multicolumn{4}{|>{\columncolor{negroTabla}\color{white}\bfseries\RaggedRight\arraybackslash}p{\dimexpr\textwidth-2\tabcolsep-2\arrayrulewidth\relax}|}{Caminos de excepcion}\\\hline
\multicolumn{4}{|>{\RaggedRight\arraybackslash}p{\dimexpr\textwidth-2\tabcolsep-2\arrayrulewidth\relax}|}{Dato atrasado, canal caído o duplicidad: marcar retraso, reintentar controladamente o evitar alertas redundantes.}\\\hline
\end{longtable}

% CU07
\begin{longtable}{|>{\RaggedRight\arraybackslash}p{.15\anchoCUAcciones}|>{\RaggedRight\arraybackslash}p{.43\anchoCUAcciones}|>{\RaggedRight\arraybackslash}p{.26\anchoCUAcciones}|>{\RaggedRight\arraybackslash}p{.16\anchoCUAcciones}|}
\caption{Caso de uso CU07: Crear y validar órdenes bursátiles.}\label{tab:cu-cu07}\\
\hline
\cellcolor{negroTabla}\color{white}\bfseries Titulo del caso de uso & \multicolumn{1}{>{\RaggedRight\arraybackslash}p{.43\anchoCUAcciones}|}{Crear y validar órdenes bursátiles} & \cellcolor{negroTabla}\color{white}\bfseries ID & CU07\\\hline
\multicolumn{4}{|>{\columncolor{negroTabla}\color{white}\bfseries\RaggedRight\arraybackslash}p{\dimexpr\textwidth-2\tabcolsep-2\arrayrulewidth\relax}|}{Descripcion general del caso de uso}\\\hline
\multicolumn{4}{|>{\RaggedRight\arraybackslash}p{\dimexpr\textwidth-2\tabcolsep-2\arrayrulewidth\relax}|}{Capturar orden de mercado, límite o stop y validar requisitos; incluye CU08 conforme al modelo, con extensión CU07.1 cuando aplique. RF-21, RF-22; CU07.1 RF-23; CU08 RF-24 a RF-27.}\\\hline
\multicolumn{4}{|>{\columncolor{negroTabla}\color{white}\bfseries\RaggedRight\arraybackslash}p{\dimexpr\textwidth-2\tabcolsep-2\arrayrulewidth\relax}|}{Entidades involucradas}\\\hline
\multicolumn{4}{|>{\RaggedRight\arraybackslash}p{\dimexpr\textwidth-2\tabcolsep-2\arrayrulewidth\relax}|}{Inversionista; servicio de fondos/posiciones}\\\hline
\multicolumn{4}{|>{\columncolor{negroTabla}\color{white}\bfseries\RaggedRight\arraybackslash}p{\dimexpr\textwidth-2\tabcolsep-2\arrayrulewidth\relax}|}{Precondiciones}\\\hline
\multicolumn{4}{|>{\RaggedRight\arraybackslash}p{\dimexpr\textwidth-2\tabcolsep-2\arrayrulewidth\relax}|}{Sesión autorizada y MFA para acción crítica; instrumento operable y datos mínimos disponibles.}\\\hline
\multicolumn{4}{|>{\columncolor{negroTabla}\color{white}\bfseries\RaggedRight\arraybackslash}p{\dimexpr\textwidth-2\tabcolsep-2\arrayrulewidth\relax}|}{Flujo normal de eventos}\\\hline
\cellcolor{gray!18}\bfseries Paso & \cellcolor{gray!18}\bfseries Usuario / actor & \multicolumn{2}{>{\columncolor{gray!18}\RaggedRight\arraybackslash}p{\dimexpr.42\anchoCUAcciones+2\tabcolsep+\arrayrulewidth\relax}|}{\bfseries Sistema}\\\hline
1 & Selecciona activo, tipo, sentido y cantidad/precio. & \multicolumn{2}{>{\RaggedRight\arraybackslash}p{\dimexpr.42\anchoCUAcciones+2\tabcolsep+\arrayrulewidth\relax}|}{Valida formato, fondos o posiciones, límites y condiciones de mercado.}\\\hline
2 & Revisa resumen de costos y confirma solicitud. & \multicolumn{2}{>{\RaggedRight\arraybackslash}p{\dimexpr.42\anchoCUAcciones+2\tabcolsep+\arrayrulewidth\relax}|}{Registra orden y la deriva al trámite CU08; informa estado inicial.}\\\hline
\multicolumn{4}{|>{\columncolor{negroTabla}\color{white}\bfseries\RaggedRight\arraybackslash}p{\dimexpr\textwidth-2\tabcolsep-2\arrayrulewidth\relax}|}{Postcondiciones principales del caso de uso}\\\hline
\multicolumn{4}{|>{\RaggedRight\arraybackslash}p{\dimexpr\textwidth-2\tabcolsep-2\arrayrulewidth\relax}|}{Solicitud validada y registrada; no se presupone ejecución final.}\\\hline
\multicolumn{4}{|>{\columncolor{negroTabla}\color{white}\bfseries\RaggedRight\arraybackslash}p{\dimexpr\textwidth-2\tabcolsep-2\arrayrulewidth\relax}|}{Caminos de excepcion}\\\hline
\multicolumn{4}{|>{\RaggedRight\arraybackslash}p{\dimexpr\textwidth-2\tabcolsep-2\arrayrulewidth\relax}|}{Saldo insuficiente, posición inexistente, mercado restringido, orden inválida o variación crítica: rechazar o solicitar reconfirmación.}\\\hline
\end{longtable}

% CU07.1
\begin{longtable}{|>{\RaggedRight\arraybackslash}p{.15\anchoCUAcciones}|>{\RaggedRight\arraybackslash}p{.43\anchoCUAcciones}|>{\RaggedRight\arraybackslash}p{.26\anchoCUAcciones}|>{\RaggedRight\arraybackslash}p{.16\anchoCUAcciones}|}
\caption{Caso de uso CU07.1: Reautorizar orden por variación crítica de precio.}\label{tab:cu-cu07-1}\\
\hline
\cellcolor{negroTabla}\color{white}\bfseries Titulo del caso de uso & \multicolumn{1}{>{\RaggedRight\arraybackslash}p{.43\anchoCUAcciones}|}{Reautorizar orden por variación crítica de precio} & \cellcolor{negroTabla}\color{white}\bfseries ID & CU07.1\\\hline
\multicolumn{4}{|>{\columncolor{negroTabla}\color{white}\bfseries\RaggedRight\arraybackslash}p{\dimexpr\textwidth-2\tabcolsep-2\arrayrulewidth\relax}|}{Descripcion general del caso de uso}\\\hline
\multicolumn{4}{|>{\RaggedRight\arraybackslash}p{\dimexpr\textwidth-2\tabcolsep-2\arrayrulewidth\relax}|}{Extensión condicional de CU07 cuando la variación del precio supera 3 \% según regla definida. RF-23; verificar viabilidad temporal y capacidad de control de Alpaca.}\\\hline
\multicolumn{4}{|>{\columncolor{negroTabla}\color{white}\bfseries\RaggedRight\arraybackslash}p{\dimexpr\textwidth-2\tabcolsep-2\arrayrulewidth\relax}|}{Entidades involucradas}\\\hline
\multicolumn{4}{|>{\RaggedRight\arraybackslash}p{\dimexpr\textwidth-2\tabcolsep-2\arrayrulewidth\relax}|}{Inversionista; proveedor de datos bursátiles}\\\hline
\multicolumn{4}{|>{\columncolor{negroTabla}\color{white}\bfseries\RaggedRight\arraybackslash}p{\dimexpr\textwidth-2\tabcolsep-2\arrayrulewidth\relax}|}{Precondiciones}\\\hline
\multicolumn{4}{|>{\RaggedRight\arraybackslash}p{\dimexpr\textwidth-2\tabcolsep-2\arrayrulewidth\relax}|}{Orden pendiente de envío o modificación autorizable; referencia de precio verificable.}\\\hline
\multicolumn{4}{|>{\columncolor{negroTabla}\color{white}\bfseries\RaggedRight\arraybackslash}p{\dimexpr\textwidth-2\tabcolsep-2\arrayrulewidth\relax}|}{Flujo normal de eventos}\\\hline
\cellcolor{gray!18}\bfseries Paso & \cellcolor{gray!18}\bfseries Usuario / actor & \multicolumn{2}{>{\columncolor{gray!18}\RaggedRight\arraybackslash}p{\dimexpr.42\anchoCUAcciones+2\tabcolsep+\arrayrulewidth\relax}|}{\bfseries Sistema}\\\hline
1 & Recibe aviso de variación y revisa nuevas condiciones. & \multicolumn{2}{>{\RaggedRight\arraybackslash}p{\dimexpr.42\anchoCUAcciones+2\tabcolsep+\arrayrulewidth\relax}|}{Suspende trámite y presenta precio de referencia, variación y cargos estimados.}\\\hline
2 & Confirma o rechaza la nueva condición. & \multicolumn{2}{>{\RaggedRight\arraybackslash}p{\dimexpr.42\anchoCUAcciones+2\tabcolsep+\arrayrulewidth\relax}|}{Registra decisión y continúa o cancela solicitud según estado real.}\\\hline
\multicolumn{4}{|>{\columncolor{negroTabla}\color{white}\bfseries\RaggedRight\arraybackslash}p{\dimexpr\textwidth-2\tabcolsep-2\arrayrulewidth\relax}|}{Postcondiciones principales del caso de uso}\\\hline
\multicolumn{4}{|>{\RaggedRight\arraybackslash}p{\dimexpr\textwidth-2\tabcolsep-2\arrayrulewidth\relax}|}{Reconfirmación trazada o solicitud detenida sin envío indebido.}\\\hline
\multicolumn{4}{|>{\columncolor{negroTabla}\color{white}\bfseries\RaggedRight\arraybackslash}p{\dimexpr\textwidth-2\tabcolsep-2\arrayrulewidth\relax}|}{Caminos de excepcion}\\\hline
\multicolumn{4}{|>{\RaggedRight\arraybackslash}p{\dimexpr\textwidth-2\tabcolsep-2\arrayrulewidth\relax}|}{Tiempo de confirmación agotado, precio vuelve a variar u orden ya ejecutada: no prometer reversión y conciliar estado.}\\\hline
\end{longtable}

% CU08
\begin{longtable}{|>{\RaggedRight\arraybackslash}p{.15\anchoCUAcciones}|>{\RaggedRight\arraybackslash}p{.43\anchoCUAcciones}|>{\RaggedRight\arraybackslash}p{.26\anchoCUAcciones}|>{\RaggedRight\arraybackslash}p{.16\anchoCUAcciones}|}
\caption{Caso de uso CU08: Tramitar, encolar y ejecutar órdenes en mercado.}\label{tab:cu-cu08}\\
\hline
\cellcolor{negroTabla}\color{white}\bfseries Titulo del caso de uso & \multicolumn{1}{>{\RaggedRight\arraybackslash}p{.43\anchoCUAcciones}|}{Tramitar, encolar y ejecutar órdenes en mercado} & \cellcolor{negroTabla}\color{white}\bfseries ID & CU08\\\hline
\multicolumn{4}{|>{\columncolor{negroTabla}\color{white}\bfseries\RaggedRight\arraybackslash}p{\dimexpr\textwidth-2\tabcolsep-2\arrayrulewidth\relax}|}{Descripcion general del caso de uso}\\\hline
\multicolumn{4}{|>{\RaggedRight\arraybackslash}p{\dimexpr\textwidth-2\tabcolsep-2\arrayrulewidth\relax}|}{Transmitir órdenes validadas al proveedor y reflejar sus estados asíncronos. RF-24 a RF-27; ejecución real condicionada al servicio autorizado.}\\\hline
\multicolumn{4}{|>{\columncolor{negroTabla}\color{white}\bfseries\RaggedRight\arraybackslash}p{\dimexpr\textwidth-2\tabcolsep-2\arrayrulewidth\relax}|}{Entidades involucradas}\\\hline
\multicolumn{4}{|>{\RaggedRight\arraybackslash}p{\dimexpr\textwidth-2\tabcolsep-2\arrayrulewidth\relax}|}{Inversionista; Alpaca (ejecución)}\\\hline
\multicolumn{4}{|>{\columncolor{negroTabla}\color{white}\bfseries\RaggedRight\arraybackslash}p{\dimexpr\textwidth-2\tabcolsep-2\arrayrulewidth\relax}|}{Precondiciones}\\\hline
\multicolumn{4}{|>{\RaggedRight\arraybackslash}p{\dimexpr\textwidth-2\tabcolsep-2\arrayrulewidth\relax}|}{Orden validada; credenciales y permisos de ejecución habilitados.}\\\hline
\multicolumn{4}{|>{\columncolor{negroTabla}\color{white}\bfseries\RaggedRight\arraybackslash}p{\dimexpr\textwidth-2\tabcolsep-2\arrayrulewidth\relax}|}{Flujo normal de eventos}\\\hline
\cellcolor{gray!18}\bfseries Paso & \cellcolor{gray!18}\bfseries Usuario / actor & \multicolumn{2}{>{\columncolor{gray!18}\RaggedRight\arraybackslash}p{\dimexpr.42\anchoCUAcciones+2\tabcolsep+\arrayrulewidth\relax}|}{\bfseries Sistema}\\\hline
1 & Confirma envío de orden desde CU07. & \multicolumn{2}{>{\RaggedRight\arraybackslash}p{\dimexpr.42\anchoCUAcciones+2\tabcolsep+\arrayrulewidth\relax}|}{Registra identificador idempotente, transmite y conserva acuse o rechazo.}\\\hline
2 & Consulta estado o espera notificación. & \multicolumn{2}{>{\RaggedRight\arraybackslash}p{\dimexpr.42\anchoCUAcciones+2\tabcolsep+\arrayrulewidth\relax}|}{Procesa eventos de aceptación, ejecución parcial, ejecución total o rechazo; concilia saldos y posiciones.}\\\hline
\multicolumn{4}{|>{\columncolor{negroTabla}\color{white}\bfseries\RaggedRight\arraybackslash}p{\dimexpr\textwidth-2\tabcolsep-2\arrayrulewidth\relax}|}{Postcondiciones principales del caso de uso}\\\hline
\multicolumn{4}{|>{\RaggedRight\arraybackslash}p{\dimexpr\textwidth-2\tabcolsep-2\arrayrulewidth\relax}|}{Orden con historial y estado confirmado por proveedor; ejecución solo si se recibe evidencia.}\\\hline
\multicolumn{4}{|>{\columncolor{negroTabla}\color{white}\bfseries\RaggedRight\arraybackslash}p{\dimexpr\textwidth-2\tabcolsep-2\arrayrulewidth\relax}|}{Caminos de excepcion}\\\hline
\multicolumn{4}{|>{\RaggedRight\arraybackslash}p{\dimexpr\textwidth-2\tabcolsep-2\arrayrulewidth\relax}|}{Timeout, duplicación, rechazo, mercado cerrado o estado desconocido: consultar y conciliar antes de reintentar.}\\\hline
\end{longtable}

% CU09
\begin{longtable}{|>{\RaggedRight\arraybackslash}p{.15\anchoCUAcciones}|>{\RaggedRight\arraybackslash}p{.43\anchoCUAcciones}|>{\RaggedRight\arraybackslash}p{.26\anchoCUAcciones}|>{\RaggedRight\arraybackslash}p{.16\anchoCUAcciones}|}
\caption{Caso de uso CU09: Cancelar órdenes bursátiles pendientes.}\label{tab:cu-cu09}\\
\hline
\cellcolor{negroTabla}\color{white}\bfseries Titulo del caso de uso & \multicolumn{1}{>{\RaggedRight\arraybackslash}p{.43\anchoCUAcciones}|}{Cancelar órdenes bursátiles pendientes} & \cellcolor{negroTabla}\color{white}\bfseries ID & CU09\\\hline
\multicolumn{4}{|>{\columncolor{negroTabla}\color{white}\bfseries\RaggedRight\arraybackslash}p{\dimexpr\textwidth-2\tabcolsep-2\arrayrulewidth\relax}|}{Descripcion general del caso de uso}\\\hline
\multicolumn{4}{|>{\RaggedRight\arraybackslash}p{\dimexpr\textwidth-2\tabcolsep-2\arrayrulewidth\relax}|}{Solicitar cancelación de una orden aún cancelable. RF-28; cancelación depende del proveedor.}\\\hline
\multicolumn{4}{|>{\columncolor{negroTabla}\color{white}\bfseries\RaggedRight\arraybackslash}p{\dimexpr\textwidth-2\tabcolsep-2\arrayrulewidth\relax}|}{Entidades involucradas}\\\hline
\multicolumn{4}{|>{\RaggedRight\arraybackslash}p{\dimexpr\textwidth-2\tabcolsep-2\arrayrulewidth\relax}|}{Inversionista; Alpaca (ejecución)}\\\hline
\multicolumn{4}{|>{\columncolor{negroTabla}\color{white}\bfseries\RaggedRight\arraybackslash}p{\dimexpr\textwidth-2\tabcolsep-2\arrayrulewidth\relax}|}{Precondiciones}\\\hline
\multicolumn{4}{|>{\RaggedRight\arraybackslash}p{\dimexpr\textwidth-2\tabcolsep-2\arrayrulewidth\relax}|}{Orden propia identificada y estado susceptible de cancelación.}\\\hline
\multicolumn{4}{|>{\columncolor{negroTabla}\color{white}\bfseries\RaggedRight\arraybackslash}p{\dimexpr\textwidth-2\tabcolsep-2\arrayrulewidth\relax}|}{Flujo normal de eventos}\\\hline
\cellcolor{gray!18}\bfseries Paso & \cellcolor{gray!18}\bfseries Usuario / actor & \multicolumn{2}{>{\columncolor{gray!18}\RaggedRight\arraybackslash}p{\dimexpr.42\anchoCUAcciones+2\tabcolsep+\arrayrulewidth\relax}|}{\bfseries Sistema}\\\hline
1 & Selecciona orden y solicita cancelación. & \multicolumn{2}{>{\RaggedRight\arraybackslash}p{\dimexpr.42\anchoCUAcciones+2\tabcolsep+\arrayrulewidth\relax}|}{Verifica titularidad y estado; transmite solicitud al proveedor.}\\\hline
2 & Consulta resultado. & \multicolumn{2}{>{\RaggedRight\arraybackslash}p{\dimexpr.42\anchoCUAcciones+2\tabcolsep+\arrayrulewidth\relax}|}{Actualiza estado con respuesta del proveedor y conserva trazabilidad.}\\\hline
\multicolumn{4}{|>{\columncolor{negroTabla}\color{white}\bfseries\RaggedRight\arraybackslash}p{\dimexpr\textwidth-2\tabcolsep-2\arrayrulewidth\relax}|}{Postcondiciones principales del caso de uso}\\\hline
\multicolumn{4}{|>{\RaggedRight\arraybackslash}p{\dimexpr\textwidth-2\tabcolsep-2\arrayrulewidth\relax}|}{Cancelación confirmada o estado vigente informado, sin suponer éxito inmediato.}\\\hline
\multicolumn{4}{|>{\columncolor{negroTabla}\color{white}\bfseries\RaggedRight\arraybackslash}p{\dimexpr\textwidth-2\tabcolsep-2\arrayrulewidth\relax}|}{Caminos de excepcion}\\\hline
\multicolumn{4}{|>{\RaggedRight\arraybackslash}p{\dimexpr\textwidth-2\tabcolsep-2\arrayrulewidth\relax}|}{Orden ejecutada, parcialmente ejecutada o cancelación rechazada: informar ejecución remanente y evitar estados falsos.}\\\hline
\end{longtable}

% CU10
\begin{longtable}{|>{\RaggedRight\arraybackslash}p{.15\anchoCUAcciones}|>{\RaggedRight\arraybackslash}p{.43\anchoCUAcciones}|>{\RaggedRight\arraybackslash}p{.26\anchoCUAcciones}|>{\RaggedRight\arraybackslash}p{.16\anchoCUAcciones}|}
\caption{Caso de uso CU10: Gestionar fondos, retiros y conversión a USD.}\label{tab:cu-cu10}\\
\hline
\cellcolor{negroTabla}\color{white}\bfseries Titulo del caso de uso & \multicolumn{1}{>{\RaggedRight\arraybackslash}p{.43\anchoCUAcciones}|}{Gestionar fondos, retiros y conversión a USD} & \cellcolor{negroTabla}\color{white}\bfseries ID & CU10\\\hline
\multicolumn{4}{|>{\columncolor{negroTabla}\color{white}\bfseries\RaggedRight\arraybackslash}p{\dimexpr\textwidth-2\tabcolsep-2\arrayrulewidth\relax}|}{Descripcion general del caso de uso}\\\hline
\multicolumn{4}{|>{\RaggedRight\arraybackslash}p{\dimexpr\textwidth-2\tabcolsep-2\arrayrulewidth\relax}|}{Consultar saldos y tramitar depósitos, retiros, conversión COP/USD y reservas según capacidades autorizadas. RF-29 a RF-33; TRM y custodia requieren validación de proveedor.}\\\hline
\multicolumn{4}{|>{\columncolor{negroTabla}\color{white}\bfseries\RaggedRight\arraybackslash}p{\dimexpr\textwidth-2\tabcolsep-2\arrayrulewidth\relax}|}{Entidades involucradas}\\\hline
\multicolumn{4}{|>{\RaggedRight\arraybackslash}p{\dimexpr\textwidth-2\tabcolsep-2\arrayrulewidth\relax}|}{Inversionista; proveedor financiero habilitado}\\\hline
\multicolumn{4}{|>{\columncolor{negroTabla}\color{white}\bfseries\RaggedRight\arraybackslash}p{\dimexpr\textwidth-2\tabcolsep-2\arrayrulewidth\relax}|}{Precondiciones}\\\hline
\multicolumn{4}{|>{\RaggedRight\arraybackslash}p{\dimexpr\textwidth-2\tabcolsep-2\arrayrulewidth\relax}|}{Cuenta verificada, permisos aplicables y servicio financiero realmente disponible.}\\\hline
\multicolumn{4}{|>{\columncolor{negroTabla}\color{white}\bfseries\RaggedRight\arraybackslash}p{\dimexpr\textwidth-2\tabcolsep-2\arrayrulewidth\relax}|}{Flujo normal de eventos}\\\hline
\cellcolor{gray!18}\bfseries Paso & \cellcolor{gray!18}\bfseries Usuario / actor & \multicolumn{2}{>{\columncolor{gray!18}\RaggedRight\arraybackslash}p{\dimexpr.42\anchoCUAcciones+2\tabcolsep+\arrayrulewidth\relax}|}{\bfseries Sistema}\\\hline
1 & Solicita depósito, retiro o conversión e introduce importe. & \multicolumn{2}{>{\RaggedRight\arraybackslash}p{\dimexpr.42\anchoCUAcciones+2\tabcolsep+\arrayrulewidth\relax}|}{Verifica límites, fondos, tasas de referencia, comisiones y soporte de operación.}\\\hline
2 & Revisa detalle y autoriza. & \multicolumn{2}{>{\RaggedRight\arraybackslash}p{\dimexpr.42\anchoCUAcciones+2\tabcolsep+\arrayrulewidth\relax}|}{Registra solicitud y actualiza saldos solo con confirmación válida; bloquea reservas cuando corresponda.}\\\hline
\multicolumn{4}{|>{\columncolor{negroTabla}\color{white}\bfseries\RaggedRight\arraybackslash}p{\dimexpr\textwidth-2\tabcolsep-2\arrayrulewidth\relax}|}{Postcondiciones principales del caso de uso}\\\hline
\multicolumn{4}{|>{\RaggedRight\arraybackslash}p{\dimexpr\textwidth-2\tabcolsep-2\arrayrulewidth\relax}|}{Movimiento con estado, importe, moneda, tasa aplicada y comprobante disponible.}\\\hline
\multicolumn{4}{|>{\columncolor{negroTabla}\color{white}\bfseries\RaggedRight\arraybackslash}p{\dimexpr\textwidth-2\tabcolsep-2\arrayrulewidth\relax}|}{Caminos de excepcion}\\\hline
\multicolumn{4}{|>{\RaggedRight\arraybackslash}p{\dimexpr\textwidth-2\tabcolsep-2\arrayrulewidth\relax}|}{Saldo insuficiente, tasa caducada, servicio no habilitado, rechazo o error externo: no acreditar fondos ficticios.}\\\hline
\end{longtable}

% CU11
\begin{longtable}{|>{\RaggedRight\arraybackslash}p{.15\anchoCUAcciones}|>{\RaggedRight\arraybackslash}p{.43\anchoCUAcciones}|>{\RaggedRight\arraybackslash}p{.26\anchoCUAcciones}|>{\RaggedRight\arraybackslash}p{.16\anchoCUAcciones}|}
\caption{Caso de uso CU11: Gestionar suscripciones premium (Stripe).}\label{tab:cu-cu11}\\
\hline
\cellcolor{negroTabla}\color{white}\bfseries Titulo del caso de uso & \multicolumn{1}{>{\RaggedRight\arraybackslash}p{.43\anchoCUAcciones}|}{Gestionar suscripciones premium (Stripe)} & \cellcolor{negroTabla}\color{white}\bfseries ID & CU11\\\hline
\multicolumn{4}{|>{\columncolor{negroTabla}\color{white}\bfseries\RaggedRight\arraybackslash}p{\dimexpr\textwidth-2\tabcolsep-2\arrayrulewidth\relax}|}{Descripcion general del caso de uso}\\\hline
\multicolumn{4}{|>{\RaggedRight\arraybackslash}p{\dimexpr\textwidth-2\tabcolsep-2\arrayrulewidth\relax}|}{Contratar, renovar o cancelar suscripción premium y actualizar derechos tras confirmación del pago. RF-34 a RF-36; tarifas preliminares sujetas a aprobación.}\\\hline
\multicolumn{4}{|>{\columncolor{negroTabla}\color{white}\bfseries\RaggedRight\arraybackslash}p{\dimexpr\textwidth-2\tabcolsep-2\arrayrulewidth\relax}|}{Entidades involucradas}\\\hline
\multicolumn{4}{|>{\RaggedRight\arraybackslash}p{\dimexpr\textwidth-2\tabcolsep-2\arrayrulewidth\relax}|}{Inversionista; Stripe}\\\hline
\multicolumn{4}{|>{\columncolor{negroTabla}\color{white}\bfseries\RaggedRight\arraybackslash}p{\dimexpr\textwidth-2\tabcolsep-2\arrayrulewidth\relax}|}{Precondiciones}\\\hline
\multicolumn{4}{|>{\RaggedRight\arraybackslash}p{\dimexpr\textwidth-2\tabcolsep-2\arrayrulewidth\relax}|}{Usuario autenticado; planes comerciales configurados; integración Stripe habilitada.}\\\hline
\multicolumn{4}{|>{\columncolor{negroTabla}\color{white}\bfseries\RaggedRight\arraybackslash}p{\dimexpr\textwidth-2\tabcolsep-2\arrayrulewidth\relax}|}{Flujo normal de eventos}\\\hline
\cellcolor{gray!18}\bfseries Paso & \cellcolor{gray!18}\bfseries Usuario / actor & \multicolumn{2}{>{\columncolor{gray!18}\RaggedRight\arraybackslash}p{\dimexpr.42\anchoCUAcciones+2\tabcolsep+\arrayrulewidth\relax}|}{\bfseries Sistema}\\\hline
1 & Selecciona plan y solicita pago o cambio de suscripción. & \multicolumn{2}{>{\RaggedRight\arraybackslash}p{\dimexpr.42\anchoCUAcciones+2\tabcolsep+\arrayrulewidth\relax}|}{Crea sesión de pago y presenta condiciones y precio vigente.}\\\hline
2 & Completa operación en interfaz segura del proveedor. & \multicolumn{2}{>{\RaggedRight\arraybackslash}p{\dimexpr.42\anchoCUAcciones+2\tabcolsep+\arrayrulewidth\relax}|}{Verifica evento firmado/confirmación de Stripe y actualiza vigencia y beneficios.}\\\hline
\multicolumn{4}{|>{\columncolor{negroTabla}\color{white}\bfseries\RaggedRight\arraybackslash}p{\dimexpr\textwidth-2\tabcolsep-2\arrayrulewidth\relax}|}{Postcondiciones principales del caso de uso}\\\hline
\multicolumn{4}{|>{\RaggedRight\arraybackslash}p{\dimexpr\textwidth-2\tabcolsep-2\arrayrulewidth\relax}|}{Estado de suscripción, comprobante o referencia y fecha de renovación.}\\\hline
\multicolumn{4}{|>{\columncolor{negroTabla}\color{white}\bfseries\RaggedRight\arraybackslash}p{\dimexpr\textwidth-2\tabcolsep-2\arrayrulewidth\relax}|}{Caminos de excepcion}\\\hline
\multicolumn{4}{|>{\RaggedRight\arraybackslash}p{\dimexpr\textwidth-2\tabcolsep-2\arrayrulewidth\relax}|}{Pago rechazado, webhook duplicado, confirmación tardía o cancelación: no activar beneficios sin validación.}\\\hline
\end{longtable}

% CU12
\begin{longtable}{|>{\RaggedRight\arraybackslash}p{.15\anchoCUAcciones}|>{\RaggedRight\arraybackslash}p{.43\anchoCUAcciones}|>{\RaggedRight\arraybackslash}p{.26\anchoCUAcciones}|>{\RaggedRight\arraybackslash}p{.16\anchoCUAcciones}|}
\caption{Caso de uso CU12: Consultar portafolio, rendimientos y comisiones.}\label{tab:cu-cu12}\\
\hline
\cellcolor{negroTabla}\color{white}\bfseries Titulo del caso de uso & \multicolumn{1}{>{\RaggedRight\arraybackslash}p{.43\anchoCUAcciones}|}{Consultar portafolio, rendimientos y comisiones} & \cellcolor{negroTabla}\color{white}\bfseries ID & CU12\\\hline
\multicolumn{4}{|>{\columncolor{negroTabla}\color{white}\bfseries\RaggedRight\arraybackslash}p{\dimexpr\textwidth-2\tabcolsep-2\arrayrulewidth\relax}|}{Descripcion general del caso de uso}\\\hline
\multicolumn{4}{|>{\RaggedRight\arraybackslash}p{\dimexpr\textwidth-2\tabcolsep-2\arrayrulewidth\relax}|}{Visualizar posiciones, valorizaciones, pérdidas/ganancias y comisiones según permisos. RF-37 a RF-40; rendimientos estimados no equivalen a liquidación.}\\\hline
\multicolumn{4}{|>{\columncolor{negroTabla}\color{white}\bfseries\RaggedRight\arraybackslash}p{\dimexpr\textwidth-2\tabcolsep-2\arrayrulewidth\relax}|}{Entidades involucradas}\\\hline
\multicolumn{4}{|>{\RaggedRight\arraybackslash}p{\dimexpr\textwidth-2\tabcolsep-2\arrayrulewidth\relax}|}{Inversionista, Comisionista (solo datos autorizados)}\\\hline
\multicolumn{4}{|>{\columncolor{negroTabla}\color{white}\bfseries\RaggedRight\arraybackslash}p{\dimexpr\textwidth-2\tabcolsep-2\arrayrulewidth\relax}|}{Precondiciones}\\\hline
\multicolumn{4}{|>{\RaggedRight\arraybackslash}p{\dimexpr\textwidth-2\tabcolsep-2\arrayrulewidth\relax}|}{Sesión activa; información financiera disponible y conciliada.}\\\hline
\multicolumn{4}{|>{\columncolor{negroTabla}\color{white}\bfseries\RaggedRight\arraybackslash}p{\dimexpr\textwidth-2\tabcolsep-2\arrayrulewidth\relax}|}{Flujo normal de eventos}\\\hline
\cellcolor{gray!18}\bfseries Paso & \cellcolor{gray!18}\bfseries Usuario / actor & \multicolumn{2}{>{\columncolor{gray!18}\RaggedRight\arraybackslash}p{\dimexpr.42\anchoCUAcciones+2\tabcolsep+\arrayrulewidth\relax}|}{\bfseries Sistema}\\\hline
1 & Abre portafolio o reporte de comisiones. & \multicolumn{2}{>{\RaggedRight\arraybackslash}p{\dimexpr.42\anchoCUAcciones+2\tabcolsep+\arrayrulewidth\relax}|}{Consulta posiciones y precios, calcula métricas e indica corte temporal.}\\\hline
2 & Filtra período o instrumento. & \multicolumn{2}{>{\RaggedRight\arraybackslash}p{\dimexpr.42\anchoCUAcciones+2\tabcolsep+\arrayrulewidth\relax}|}{Muestra detalle y restringe información según rol y vínculo.}\\\hline
\multicolumn{4}{|>{\columncolor{negroTabla}\color{white}\bfseries\RaggedRight\arraybackslash}p{\dimexpr\textwidth-2\tabcolsep-2\arrayrulewidth\relax}|}{Postcondiciones principales del caso de uso}\\\hline
\multicolumn{4}{|>{\RaggedRight\arraybackslash}p{\dimexpr\textwidth-2\tabcolsep-2\arrayrulewidth\relax}|}{Vista consolidada y trazable de saldos, rendimientos y comisiones.}\\\hline
\multicolumn{4}{|>{\columncolor{negroTabla}\color{white}\bfseries\RaggedRight\arraybackslash}p{\dimexpr\textwidth-2\tabcolsep-2\arrayrulewidth\relax}|}{Caminos de excepcion}\\\hline
\multicolumn{4}{|>{\RaggedRight\arraybackslash}p{\dimexpr\textwidth-2\tabcolsep-2\arrayrulewidth\relax}|}{Precio no disponible, posición sin conciliar o acceso indebido: mostrar advertencia o denegar datos.}\\\hline
\end{longtable}

% CU13
\begin{longtable}{|>{\RaggedRight\arraybackslash}p{.15\anchoCUAcciones}|>{\RaggedRight\arraybackslash}p{.43\anchoCUAcciones}|>{\RaggedRight\arraybackslash}p{.26\anchoCUAcciones}|>{\RaggedRight\arraybackslash}p{.16\anchoCUAcciones}|}
\caption{Caso de uso CU13: Parametrizar mercados, horarios y comisiones.}\label{tab:cu-cu13}\\
\hline
\cellcolor{negroTabla}\color{white}\bfseries Titulo del caso de uso & \multicolumn{1}{>{\RaggedRight\arraybackslash}p{.43\anchoCUAcciones}|}{Parametrizar mercados, horarios y comisiones} & \cellcolor{negroTabla}\color{white}\bfseries ID & CU13\\\hline
\multicolumn{4}{|>{\columncolor{negroTabla}\color{white}\bfseries\RaggedRight\arraybackslash}p{\dimexpr\textwidth-2\tabcolsep-2\arrayrulewidth\relax}|}{Descripcion general del caso de uso}\\\hline
\multicolumn{4}{|>{\RaggedRight\arraybackslash}p{\dimexpr\textwidth-2\tabcolsep-2\arrayrulewidth\relax}|}{Administrar parámetros operativos, horarios, feriados y porcentajes de comisión. RF-41 a RF-43; no puede ampliar capacidades reales de Alpaca.}\\\hline
\multicolumn{4}{|>{\columncolor{negroTabla}\color{white}\bfseries\RaggedRight\arraybackslash}p{\dimexpr\textwidth-2\tabcolsep-2\arrayrulewidth\relax}|}{Entidades involucradas}\\\hline
\multicolumn{4}{|>{\RaggedRight\arraybackslash}p{\dimexpr\textwidth-2\tabcolsep-2\arrayrulewidth\relax}|}{Administrador de la Plataforma}\\\hline
\multicolumn{4}{|>{\columncolor{negroTabla}\color{white}\bfseries\RaggedRight\arraybackslash}p{\dimexpr\textwidth-2\tabcolsep-2\arrayrulewidth\relax}|}{Precondiciones}\\\hline
\multicolumn{4}{|>{\RaggedRight\arraybackslash}p{\dimexpr\textwidth-2\tabcolsep-2\arrayrulewidth\relax}|}{Administrador autenticado con privilegios y MFA para cambios críticos.}\\\hline
\multicolumn{4}{|>{\columncolor{negroTabla}\color{white}\bfseries\RaggedRight\arraybackslash}p{\dimexpr\textwidth-2\tabcolsep-2\arrayrulewidth\relax}|}{Flujo normal de eventos}\\\hline
\cellcolor{gray!18}\bfseries Paso & \cellcolor{gray!18}\bfseries Usuario / actor & \multicolumn{2}{>{\columncolor{gray!18}\RaggedRight\arraybackslash}p{\dimexpr.42\anchoCUAcciones+2\tabcolsep+\arrayrulewidth\relax}|}{\bfseries Sistema}\\\hline
1 & Edita horario, día inhábil o esquema de comisión. & \multicolumn{2}{>{\RaggedRight\arraybackslash}p{\dimexpr.42\anchoCUAcciones+2\tabcolsep+\arrayrulewidth\relax}|}{Valida coherencia, restricciones del proveedor y efecto del cambio.}\\\hline
2 & Confirma modificación. & \multicolumn{2}{>{\RaggedRight\arraybackslash}p{\dimexpr.42\anchoCUAcciones+2\tabcolsep+\arrayrulewidth\relax}|}{Versiona configuración, registra auditoría y aplica vigencia definida.}\\\hline
\multicolumn{4}{|>{\columncolor{negroTabla}\color{white}\bfseries\RaggedRight\arraybackslash}p{\dimexpr\textwidth-2\tabcolsep-2\arrayrulewidth\relax}|}{Postcondiciones principales del caso de uso}\\\hline
\multicolumn{4}{|>{\RaggedRight\arraybackslash}p{\dimexpr\textwidth-2\tabcolsep-2\arrayrulewidth\relax}|}{Parámetros actualizados y trazabilidad de autor, fecha y versión.}\\\hline
\multicolumn{4}{|>{\columncolor{negroTabla}\color{white}\bfseries\RaggedRight\arraybackslash}p{\dimexpr\textwidth-2\tabcolsep-2\arrayrulewidth\relax}|}{Caminos de excepcion}\\\hline
\multicolumn{4}{|>{\RaggedRight\arraybackslash}p{\dimexpr\textwidth-2\tabcolsep-2\arrayrulewidth\relax}|}{Conflicto con reglas de Alpaca, valor inválido o falta de permiso: rechazar y conservar configuración anterior.}\\\hline
\end{longtable}

% CU14
\begin{longtable}{|>{\RaggedRight\arraybackslash}p{.15\anchoCUAcciones}|>{\RaggedRight\arraybackslash}p{.43\anchoCUAcciones}|>{\RaggedRight\arraybackslash}p{.26\anchoCUAcciones}|>{\RaggedRight\arraybackslash}p{.16\anchoCUAcciones}|}
\caption{Caso de uso CU14: Auditar operaciones y gestionar cuentas restringidas.}\label{tab:cu-cu14}\\
\hline
\cellcolor{negroTabla}\color{white}\bfseries Titulo del caso de uso & \multicolumn{1}{>{\RaggedRight\arraybackslash}p{.43\anchoCUAcciones}|}{Auditar operaciones y gestionar cuentas restringidas} & \cellcolor{negroTabla}\color{white}\bfseries ID & CU14\\\hline
\multicolumn{4}{|>{\columncolor{negroTabla}\color{white}\bfseries\RaggedRight\arraybackslash}p{\dimexpr\textwidth-2\tabcolsep-2\arrayrulewidth\relax}|}{Descripcion general del caso de uso}\\\hline
\multicolumn{4}{|>{\RaggedRight\arraybackslash}p{\dimexpr\textwidth-2\tabcolsep-2\arrayrulewidth\relax}|}{Inspeccionar eventos, identificar anomalías y restringir cuentas con justificación. RF-44 y RF-45; protección de datos y retención según política.}\\\hline
\multicolumn{4}{|>{\columncolor{negroTabla}\color{white}\bfseries\RaggedRight\arraybackslash}p{\dimexpr\textwidth-2\tabcolsep-2\arrayrulewidth\relax}|}{Entidades involucradas}\\\hline
\multicolumn{4}{|>{\RaggedRight\arraybackslash}p{\dimexpr\textwidth-2\tabcolsep-2\arrayrulewidth\relax}|}{Analista de Seguridad y Auditoría}\\\hline
\multicolumn{4}{|>{\columncolor{negroTabla}\color{white}\bfseries\RaggedRight\arraybackslash}p{\dimexpr\textwidth-2\tabcolsep-2\arrayrulewidth\relax}|}{Precondiciones}\\\hline
\multicolumn{4}{|>{\RaggedRight\arraybackslash}p{\dimexpr\textwidth-2\tabcolsep-2\arrayrulewidth\relax}|}{Analista autenticado y autorizado; registros de auditoría disponibles.}\\\hline
\multicolumn{4}{|>{\columncolor{negroTabla}\color{white}\bfseries\RaggedRight\arraybackslash}p{\dimexpr\textwidth-2\tabcolsep-2\arrayrulewidth\relax}|}{Flujo normal de eventos}\\\hline
\cellcolor{gray!18}\bfseries Paso & \cellcolor{gray!18}\bfseries Usuario / actor & \multicolumn{2}{>{\columncolor{gray!18}\RaggedRight\arraybackslash}p{\dimexpr.42\anchoCUAcciones+2\tabcolsep+\arrayrulewidth\relax}|}{\bfseries Sistema}\\\hline
1 & Consulta historial o investiga alerta. & \multicolumn{2}{>{\RaggedRight\arraybackslash}p{\dimexpr.42\anchoCUAcciones+2\tabcolsep+\arrayrulewidth\relax}|}{Presenta eventos íntegros y filtros con trazabilidad.}\\\hline
2 & Documenta hallazgo y solicita restricción de cuenta. & \multicolumn{2}{>{\RaggedRight\arraybackslash}p{\dimexpr.42\anchoCUAcciones+2\tabcolsep+\arrayrulewidth\relax}|}{Valida facultades, aplica bloqueo preventivo y registra motivo y responsable.}\\\hline
\multicolumn{4}{|>{\columncolor{negroTabla}\color{white}\bfseries\RaggedRight\arraybackslash}p{\dimexpr\textwidth-2\tabcolsep-2\arrayrulewidth\relax}|}{Postcondiciones principales del caso de uso}\\\hline
\multicolumn{4}{|>{\RaggedRight\arraybackslash}p{\dimexpr\textwidth-2\tabcolsep-2\arrayrulewidth\relax}|}{Evidencia de auditoría y estado de cuenta actualizado.}\\\hline
\multicolumn{4}{|>{\columncolor{negroTabla}\color{white}\bfseries\RaggedRight\arraybackslash}p{\dimexpr\textwidth-2\tabcolsep-2\arrayrulewidth\relax}|}{Caminos de excepcion}\\\hline
\multicolumn{4}{|>{\RaggedRight\arraybackslash}p{\dimexpr\textwidth-2\tabcolsep-2\arrayrulewidth\relax}|}{Registro incompleto, falta de privilegios o error de bloqueo: escalar incidencia y evitar cambios no auditados.}\\\hline
\end{longtable}

% CU15
\begin{longtable}{|>{\RaggedRight\arraybackslash}p{.15\anchoCUAcciones}|>{\RaggedRight\arraybackslash}p{.43\anchoCUAcciones}|>{\RaggedRight\arraybackslash}p{.26\anchoCUAcciones}|>{\RaggedRight\arraybackslash}p{.16\anchoCUAcciones}|}
\caption{Caso de uso CU15: Consultar tablero de indicadores y métricas del negocio.}\label{tab:cu-cu15}\\
\hline
\cellcolor{negroTabla}\color{white}\bfseries Titulo del caso de uso & \multicolumn{1}{>{\RaggedRight\arraybackslash}p{.43\anchoCUAcciones}|}{Consultar tablero de indicadores y métricas del negocio} & \cellcolor{negroTabla}\color{white}\bfseries ID & CU15\\\hline
\multicolumn{4}{|>{\columncolor{negroTabla}\color{white}\bfseries\RaggedRight\arraybackslash}p{\dimexpr\textwidth-2\tabcolsep-2\arrayrulewidth\relax}|}{Descripcion general del caso de uso}\\\hline
\multicolumn{4}{|>{\RaggedRight\arraybackslash}p{\dimexpr\textwidth-2\tabcolsep-2\arrayrulewidth\relax}|}{Consultar BI con volumen transaccional, ingresos por comisiones y métricas de retención. RF-46 y RF-47; no presentar estimaciones como cifras contables cerradas.}\\\hline
\multicolumn{4}{|>{\columncolor{negroTabla}\color{white}\bfseries\RaggedRight\arraybackslash}p{\dimexpr\textwidth-2\tabcolsep-2\arrayrulewidth\relax}|}{Entidades involucradas}\\\hline
\multicolumn{4}{|>{\RaggedRight\arraybackslash}p{\dimexpr\textwidth-2\tabcolsep-2\arrayrulewidth\relax}|}{Directivo del Negocio}\\\hline
\multicolumn{4}{|>{\columncolor{negroTabla}\color{white}\bfseries\RaggedRight\arraybackslash}p{\dimexpr\textwidth-2\tabcolsep-2\arrayrulewidth\relax}|}{Precondiciones}\\\hline
\multicolumn{4}{|>{\RaggedRight\arraybackslash}p{\dimexpr\textwidth-2\tabcolsep-2\arrayrulewidth\relax}|}{Directivo autorizado; datos agregados disponibles.}\\\hline
\multicolumn{4}{|>{\columncolor{negroTabla}\color{white}\bfseries\RaggedRight\arraybackslash}p{\dimexpr\textwidth-2\tabcolsep-2\arrayrulewidth\relax}|}{Flujo normal de eventos}\\\hline
\cellcolor{gray!18}\bfseries Paso & \cellcolor{gray!18}\bfseries Usuario / actor & \multicolumn{2}{>{\columncolor{gray!18}\RaggedRight\arraybackslash}p{\dimexpr.42\anchoCUAcciones+2\tabcolsep+\arrayrulewidth\relax}|}{\bfseries Sistema}\\\hline
1 & Selecciona indicadores y período. & \multicolumn{2}{>{\RaggedRight\arraybackslash}p{\dimexpr.42\anchoCUAcciones+2\tabcolsep+\arrayrulewidth\relax}|}{Agrega datos permitidos y muestra fuentes y fecha de corte.}\\\hline
2 & Aplica filtros y consulta desglose autorizado. & \multicolumn{2}{>{\RaggedRight\arraybackslash}p{\dimexpr.42\anchoCUAcciones+2\tabcolsep+\arrayrulewidth\relax}|}{Actualiza gráficos y exporta reporte si está habilitado.}\\\hline
\multicolumn{4}{|>{\columncolor{negroTabla}\color{white}\bfseries\RaggedRight\arraybackslash}p{\dimexpr\textwidth-2\tabcolsep-2\arrayrulewidth\relax}|}{Postcondiciones principales del caso de uso}\\\hline
\multicolumn{4}{|>{\RaggedRight\arraybackslash}p{\dimexpr\textwidth-2\tabcolsep-2\arrayrulewidth\relax}|}{Tablero de indicadores y reporte con fecha de actualización.}\\\hline
\multicolumn{4}{|>{\columncolor{negroTabla}\color{white}\bfseries\RaggedRight\arraybackslash}p{\dimexpr\textwidth-2\tabcolsep-2\arrayrulewidth\relax}|}{Caminos de excepcion}\\\hline
\multicolumn{4}{|>{\RaggedRight\arraybackslash}p{\dimexpr\textwidth-2\tabcolsep-2\arrayrulewidth\relax}|}{Datos incompletos, retraso de consolidación o acceso indebido: advertir y limitar información.}\\\hline
\end{longtable}

\endgroup

% ============================================================================


\section{Descripción y Justificación de la Arquitectura de Software Definida}

\subsection{Estilos Arquitecturales}

Para el desarrollo de AccionesUD se ha seleccionado una arquitectura de
\textbf{Monolito Modular con enfoque de Arquitectura Hexagonal
(Puertos y Adaptadores)}, como resultado del análisis de las
alternativas arquitectónicas realizado por el equipo de desarrollo.

Esta decisión, registrada en el acta de reunión del 31 de octubre
de 2026, responde a la necesidad de construir una plataforma web
que integre operaciones bursátiles, gestión de inversionistas,
asesoría mediante comisionistas, consulta de información financiera,
administración de fondos, suscripciones y auditoría, manteniendo
una organización interna clara y controlando la complejidad
técnica y operativa del proyecto.

\subsubsection{Estilo Arquitectural Principal: Monolito Modular}

El sistema se estructurará como un monolito modular, en el cual
las funcionalidades del backend se organizarán en módulos
delimitados por responsabilidades de negocio, integrados en una
misma aplicación y desplegados como una unidad lógica.

A diferencia de un monolito tradicional sin separación explícita,
esta organización permitirá establecer límites entre los módulos,
definir contratos de comunicación internos y reducir el
acoplamiento entre funcionalidades.

La división modular propuesta, alineada con el catálogo de
casos de uso del sistema, comprende los siguientes módulos:

\begin{itemize}

    \item \textbf{Módulo 01. Seguridad, Usuarios y Perfil
    de Inversión:}
    gestiona el registro de usuarios, autenticación,
    segundo factor (MFA), autorización basada en roles
    (RBAC), administración de sesiones y configuración
    de perfiles de inversión.

    \item \textbf{Módulo 02. Gestión de Comisionistas
    y Asesoría de Inversión:}
    administra la vinculación entre inversionistas
    y comisionistas, las calificaciones, la elaboración
    de propuestas de inversión y los procesos
    de autorización correspondientes.

    \item \textbf{Módulo 03. Información Bursátil,
    Búsqueda y Alertas:}
    proporciona acceso a instrumentos financieros,
    cotizaciones, información histórica, listas de
    observación y alertas configurables mediante
    la integración con el proveedor bursátil.

    \item \textbf{Módulo 04. Operaciones Bursátiles
    y Motor de Órdenes:}
    gestiona la creación, validación, autorización,
    tramitación, seguimiento y cancelación de
    órdenes bursátiles, respetando las reglas
    de negocio y las capacidades del proveedor.

    \item \textbf{Módulo 05. Gestión de Fondos,
    Portafolio y Suscripciones:}
    administra saldos, movimientos financieros,
    posiciones, rendimientos y suscripciones
    premium, mediante los servicios internos
    y las integraciones externas correspondientes.

    \item \textbf{Módulo 06. Gestión Operativa
    y Seguimiento del Negocio:}
    agrupa las funcionalidades operativas de
    seguimiento y consulta del estado de las
    actividades de la plataforma, sin asumir
    las responsabilidades de parametrización
    central, gobierno y cumplimiento asignadas
    al Módulo 07.

    \item \textbf{Módulo 07. Administración Central,
    Regulación y Parametrización
    (Core Admin \& Compliance):}
    proporciona al Administrador Central herramientas
    para configurar comisiones, calendarios bursátiles,
    horarios de operación y mecanismos de suspensión
    preventiva. Asimismo, permite consultar registros
    de auditoría transaccional y métricas directivas,
    bajo controles estrictos de autorización y
    segregación de responsabilidades.

    \item \textbf{Módulo 08. Experiencia Móvil Adaptativa
    y Rendimiento Cross-Platform (Mobile Web \& PWA):}
    implementa capacidades transversales de
    presentación para dispositivos móviles,
    incluyendo navegación adaptativa, componentes
    táctiles, formularios de órdenes optimizados,
    carga diferida y mecanismos de caché seguros.
    Su implementación se concentrará principalmente
    en el frontend y reutilizará los servicios
    de aplicación de los módulos de negocio.

\end{itemize}

Cada módulo deberá exponer interfaces internas definidas
para sus operaciones y evitar dependencias directas
innecesarias sobre la implementación de otros módulos.

La separación modular permitirá desarrollar y probar
funcionalidades de manera relativamente independiente,
aunque la aplicación conservará un ciclo de despliegue
conjunto. Por tanto, la modularidad no implica
independencia de escalamiento o despliegue entre módulos.

\subsubsection{Estilo Arquitectural Complementario:
Arquitectura Hexagonal}

Como complemento del monolito modular, se adoptará
la Arquitectura Hexagonal, también denominada
arquitectura de Puertos y Adaptadores.

Este enfoque busca mantener la lógica de negocio
independiente de los mecanismos de presentación,
persistencia y comunicación con servicios externos.

En los módulos donde resulte pertinente,
la organización interna contemplará:

\begin{itemize}

    \item \textbf{Dominio:}
    contiene entidades, objetos de valor, reglas
    e invariantes del negocio, sin dependencias
    directas de frameworks o proveedores externos.

    \item \textbf{Aplicación:}
    coordina los casos de uso, las operaciones
    transaccionales y las interacciones entre
    componentes mediante interfaces definidas.

    \item \textbf{Puertos:}
    establecen contratos para acceder a capacidades
    externas o exponer operaciones de aplicación,
    sin depender de implementaciones concretas.

    \item \textbf{Adaptadores:}
    implementan los puertos y permiten conectar
    el sistema con interfaces HTTP, bases de datos,
    proveedores bursátiles, pasarelas de pago
    y servicios de notificación.

\end{itemize}

Las dependencias del código deberán orientarse
hacia las abstracciones internas, evitando que
el dominio dependa directamente de FastAPI,
PostgreSQL, Alpaca, Stripe u otros mecanismos
de infraestructura.

La integración bursátil se realizará mediante
adaptadores específicos para Alpaca, proveedor
seleccionado por el equipo para los servicios
de información y operaciones bursátiles.

La disponibilidad de instrumentos, mercados,
tipos de órdenes, movimientos de fondos y
ejecución efectiva dependerá de las capacidades
y condiciones de acceso verificadas para dicho
proveedor.

De forma independiente, la gestión de suscripciones
podrá integrarse con Stripe mediante un adaptador
de pagos, sujeto a la validación de sus condiciones
técnicas y comerciales.

\subsubsection{Organización Tecnológica}

La solución contempla una separación entre
la aplicación frontend y el backend modular.

El frontend se desarrollará con \textbf{Next.js},
mientras que el backend utilizará \textbf{FastAPI
con Python}. La persistencia de datos se
implementará mediante \textbf{PostgreSQL},
con un esquema relacional normalizado y
controles de integridad.

Esta separación permite que la interfaz web
consuma las funcionalidades del backend
mediante contratos API definidos, sin acceder
directamente a las estructuras internas de
persistencia ni a las credenciales de
los proveedores externos.

Aunque frontend y backend podrán desplegarse
como componentes diferenciados, los módulos
internos del backend permanecerán integrados
en una única aplicación monolítica modular.

\subsubsection{Comunicación e Integración}

La comunicación entre el frontend y el backend
se realizará principalmente mediante interfaces
HTTP y contratos API documentados.

Dentro del backend, los módulos colaborarán
a través de interfaces y servicios de aplicación,
evitando llamadas HTTP entre módulos que
se ejecutan dentro del mismo proceso.

Para las integraciones externas se utilizarán
adaptadores responsables de gestionar
autenticación, validación de respuestas,
errores, límites de consumo y trazabilidad.

En operaciones que requieran procesamiento
diferido, como el seguimiento de órdenes,
las notificaciones o la conciliación de estados,
podrán emplearse mecanismos asíncronos,
siempre que su incorporación esté justificada
por los requisitos y recursos disponibles.

El sistema deberá diferenciar la recepción,
validación y aceptación de una solicitud
de su ejecución bursátil definitiva.

Asimismo, las órdenes fuera del horario
de negociación se tratarán conforme
a las reglas del mercado y a las
capacidades efectivamente soportadas
por Alpaca.

\subsubsection{Justificación de la Decisión Arquitectónica}

La elección del Monolito Modular con
Arquitectura Hexagonal se fundamenta
en los siguientes criterios:

\begin{itemize}

    \item \textbf{Complejidad controlada:}
    reduce la necesidad inicial de administrar
    múltiples servicios distribuidos y sus
    mecanismos de comunicación.

    \item \textbf{Costos de infraestructura:}
    permite iniciar con una infraestructura
    de despliegue relativamente sencilla,
    adecuada a las restricciones presupuestales
    y académicas del proyecto.

    \item \textbf{Mantenibilidad:}
    favorece la separación de responsabilidades,
    la localización de cambios y la evolución
    de reglas de negocio.

    \item \textbf{Testabilidad:}
    facilita las pruebas unitarias y de aplicación
    mediante la sustitución de adaptadores externos
    por implementaciones simuladas.

    \item \textbf{Integración controlada:}
    disminuye el acoplamiento directo
    con proveedores externos, especialmente
    Alpaca y los servicios de pagos.

    \item \textbf{Evolución arquitectónica:}
    permite evaluar una futura extracción
    de módulos hacia servicios independientes,
    si existen necesidades técnicas
    justificadas y recursos suficientes.

\end{itemize}

Durante el análisis arquitectónico se consideraron
alternativas basadas en microservicios y serverless.

Sin embargo, estas opciones implicaban mayores
exigencias de infraestructura, coordinación,
observabilidad y administración operativa
para el alcance inicial del proyecto.

Por ello, el monolito modular se considera
la alternativa más adecuada para equilibrar
mantenibilidad, complejidad de implementación,
costos y posibilidades de evolución.



\subsection{Diagrama de Arquitectura}
\figuraPendiente{Vista de alto nivel de la arquitectura de Acciones UD}{3.40}{fig:arquitectura}

% ============================================================================
\section{Descripción y Justificación de los Patrones de Diseño Adoptados}
Todavía no se han aprobado patrones concretos. Durante el diseño detallado se recomienda evaluar las siguientes opciones, sin presentarlas como decisiones definitivas:
\begin{itemize}
    \item \textbf{Adapter:} para desacoplar la lógica del negocio de los contratos de Alpaca, de los proveedores de cotizaciones y de los servicios de notificación.
    \item \textbf{Strategy:} para modelar validaciones y comportamientos variables según el tipo de orden y las condiciones acordadas de operación.
    \item \textbf{Observer o publicación/suscripción:} para propagar eventos de cambio de estado, alertas y notificaciones a componentes interesados.
    \item \textbf{Repository:} para separar el acceso a datos persistentes de los servicios que implementan las reglas del negocio.
\end{itemize}

Para cada patrón adoptado posteriormente se registrarán el problema que resuelve, los elementos que lo conforman, sus relaciones, las alternativas consideradas y sus consecuencias sobre los atributos de calidad.

% ============================================================================
\section{Puntos de Vista y Modelos Arquitecturales}
\subsection{Punto de Vista: Descripción del Problema}
Este punto de vista delimitará los actores, sistemas externos, conceptos de negocio y relaciones esenciales del dominio financiero que cubre Acciones UD.

\subsubsection{Punto de Vista: Contexto de la Solución (Diagrama de Contexto)}
Representará a Acciones UD como sistema de interés y sus interacciones con inversionistas, comisionistas, administradores, cumplimiento y servicios externos de datos, órdenes, fondos, pagos y notificaciones.
\figuraPendiente{Diagrama de contexto de la solución}{3.20}{fig:contexto}

\subsubsection{Punto de Vista: Objetos (Diagrama de Objetos)}
Ilustrará ejemplos consistentes de instancias y asociaciones del dominio, tales como usuario, comisionista, orden, activo, portafolio y posición; sus valores se escogerán después de definir el modelo de clases.
\figuraPendiente{Diagrama de objetos}{3.20}{fig:objetos}

\subsection{Punto de Vista: Interacción}
Describirá quién inicia cada operación y cómo colaboran los componentes para alcanzar un resultado observable.

\subsubsection{Diagrama de Casos de Uso}
Debe incorporar actores, límites del sistema y casos de uso validados; estos deben coincidir con las fichas del Anexo~\ref{an:casosuso}.
\figuraPendiente{Modelo de casos de uso}{3.20}{fig:casos}

\subsubsection{Diagrama de Secuencia}
Debe detallar, como mínimo, un flujo crítico de orden de compra o venta y sus variantes de error, incluyendo la solicitud al tercero y su respuesta.
\figuraPendiente{Diagrama de secuencia de una operación crítica}{3.20}{fig:secuencia}

\subsection{Punto de Vista: Funcional (Diagrama de Componentes)}
Presentará responsabilidades e interfaces de los componentes; permitirá rastrear qué módulos atienden cada requerimiento funcional.
\figuraPendiente{Diagrama de componentes}{3.20}{fig:componentes}

\subsection{Punto de Vista: Despliegue (Diagrama de Despliegue)}
Mostrará nodos, redes, servicios desplegados, mecanismos de conexión y dependencias con proveedores externos; la plataforma deberá ser accesible públicamente y no solo en un ambiente local.
\figuraPendiente{Diagrama de despliegue}{3.20}{fig:despliegue}

\subsection{Punto de Vista: Información (Diagrama(s) de Actividades)}
Modelará procesos de negocio y sus decisiones, por ejemplo la creación de una orden, validación de condiciones, aprobación opcional, envío y actualización de estados.
\figuraPendiente{Diagrama de actividades del proceso de órdenes}{3.20}{fig:actividades}

\subsection{Punto de Vista: Desarrollo (Diagrama de Clases)}
Formalizará las clases, atributos, responsabilidades, asociaciones y restricciones relevantes, manteniendo correspondencia con los casos de uso y la persistencia.
\figuraPendiente{Diagrama de clases}{3.20}{fig:clases}

\subsection{Punto de Vista: Modelo Relacional de la BBDD}
El diseño relacional deberá representar datos de usuarios, roles, instrumentos, mercados, órdenes, eventos y demás conceptos aprobados, con claves, integridad referencial y una normalización al menos hasta \textbf{tercera forma normal (3FN)}.
\figuraPendiente{Modelo relacional de la base de datos}{3.20}{fig:relacional}

% ============================================================================
\section{Justificación y Relaciones entre los Puntos de Vista}
Los puntos de vista deberán tratar la misma solución desde distintas preocupaciones, sin introducir contradicciones entre modelos. El diagrama de contexto fijará el límite del sistema y las dependencias externas; los casos de uso y diagramas de secuencia mostrarán las capacidades y la colaboración requerida para realizarlas; la vista funcional definirá los componentes responsables; y la vista de despliegue establecerá dónde se ejecutan y cómo se comunican.

Por su parte, los diagramas de objetos y clases documentarán estructuras del dominio, mientras que los diagramas de actividades representarán los flujos de negocio. El modelo relacional traducirá a estructuras persistentes aquellos datos que deben conservarse. La consistencia se revisará mediante trazabilidad: \textbf{requerimiento $\rightarrow$ caso de uso $\rightarrow$ componentes $\rightarrow$ interacción $\rightarrow$ datos/pruebas}. Toda decisión relevante deberá justificar cómo contribuye a los escenarios de calidad y a las restricciones registradas.

% ============================================================================
\clearpage
\section{Firmas de Aceptación}
La aceptación del documento y de sus decisiones pendientes se formalizará con los responsables designados por el proyecto. Esta sección no implica aprobación mientras no existan las firmas correspondientes.

\vspace{0.55cm}
\begin{tabularx}{\textwidth}{@{}X X@{}}
\rule{.90\linewidth}{.45pt} & \rule{.90\linewidth}{.45pt}\\
\textbf{Representante del cliente} & \textbf{Representante del equipo de desarrollo}\\
Nombre: \rule{.42\linewidth}{.3pt} & Nombre: \rule{.42\linewidth}{.3pt}\\[7pt]
Cargo: \rule{.45\linewidth}{.3pt} & Cargo: \rule{.45\linewidth}{.3pt}\\[7pt]
Fecha (DD/MM/AAAA): \rule{.20\linewidth}{.3pt} & Fecha (DD/MM/AAAA): \rule{.20\linewidth}{.3pt}\\
\end{tabularx}

\vspace{0.5cm}
\textbf{Observaciones de aceptación:}\\[0.3cm]
\rule{\textwidth}{.35pt}


\end{document}
